\section{Exact operator bridge to the released Navier--Stokes profiles}
\label{nsbridge:nsb:operator}

This bounded comparison uses the locally saved released Navier--Stokes source
(the source pages are 7--15 and 24--27 for the profile, stress and pulse
construction). Its stated physical viscosity is denoted by $\nu_{\rm NS}>0$.
The Boussinesq lane variables and physical viscosity remain $d$, $\mu$ and
$\nu=\sqrt{A_0}$; none is identified with $\nu_{\rm NS}$.

Use right-handed cylindrical coordinates $(r,\vartheta,z)$ and define
\[
 y_1=z,\qquad y_2=\frac{r^2}{2}>0,\qquad
 v=(u^z,ru^r)=J\nabla_y\Psi,
 \quad J=\begin{pmatrix}0&-1\\1&0\end{pmatrix}.
\]
The meridional divergence identity is
$\partial_{y_1}v_1+\partial_{y_2}v_2=0$. Define the circulation and
meridional vorticity variables
\[
 \Gamma=ru^\vartheta,
 \qquad \xi=\frac{\omega^\vartheta}{r},
 \qquad h=\frac{\Gamma}{2y_2}=\frac{u^\vartheta}{r},
 \qquad Q=\partial_{y_1}^2+2y_2\partial_{y_2}^2,
 \qquad M=Q+4\partial_{y_2}.
\]
The azimuthal momentum and meridional vorticity rows of the cylindrical
equations (with arbitrary smooth force components) become, on $y_2>0$,
\begin{align}
 D_v\Gamma&=\nu_{\rm NS}Q\Gamma+r f^\vartheta,
 \label{nsbridge:nsb:Gamma}\\
 D_v\xi&=\nu_{\rm NS}M\xi+\partial_{y_1}(h^2)
       +\operatorname{curl}_y(f^z,f^r/r),
 \label{nsbridge:nsb:xi}
\end{align}
where $D_v=\partial_t+v\cdot\nabla_y$ and
$\operatorname{curl}_y(a,b)=\partial_{y_1}b-\partial_{y_2}a$.

Indeed, $\partial_r=r\partial_{y_2}$ and
$\partial_r^2+r^{-1}\partial_r=2y_2\partial_{y_2}^2+2\partial_{y_2}$.
For $\Gamma$, the azimuthal viscous operator is
$\partial_r^2-r^{-1}\partial_r+\partial_z^2=Q$.
For $r\xi=\omega^\vartheta$, the cylindrical operator
$(\partial_r^2+r^{-1}\partial_r-r^{-2}+\partial_z^2)(r\xi)$,
divided by $r$, is $M\xi$; explicitly the radial terms are
$2y_2\xi_{y_2y_2}+4\xi_{y_2}$. The centrifugal term is
$r^{-2}\partial_z((u^\vartheta)^2)=\partial_{y_1}(\Gamma^2/(4y_2^2))
=\partial_{y_1}(h^2)$. These calculations prove every coefficient in
(\ref{nsbridge:nsb:Gamma})--(\ref{nsbridge:nsb:xi}), including the $r$ force conversion.

The velocity--vorticity relation is also retained. Since
$u^z=-\Psi_{y_2}$ and $u^r=\Psi_{y_1}/r$,
\[
 \xi=\frac{\partial_z u^r-\partial_r u^z}{r}
     =\frac{\Psi_{y_1y_1}}{2y_2}+\Psi_{y_2y_2}
     =\frac{Q\Psi}{2y_2}.
\]
Thus the two transport rows do not by themselves replace the elliptic
velocity reconstruction. Conversely, on a simply connected meridional
domain, a smooth $\Psi,\Gamma$ satisfying this relation and both transport
rows recovers $u$. Define its momentum residual without pressure as
$a=\partial_tu+(u\cdot\nabla)u-\nu_{\rm NS}\Delta u-f$.
The circulation row is $a^\vartheta=0$, while the vorticity row is
$\partial_z a^r-\partial_r a^z=0$. The one-form
$a^r\,\mathrm dr+a^z\,\mathrm dz$ is therefore closed, and
$p=-\int(a^r\,\mathrm dr+a^z\,\mathrm dz)$ along a path from a fixed
base point is path independent. Its derivatives are $p_r=-a^r$,
$p_z=-a^z$. This proves the local momentum reconstruction, with pressure
unique up to a function of time; boundary data are not removed.

The change from circulation to swirl rate is invertible on the exact retained
domain: $\Gamma=2y_2h$. Product differentiation gives
\[
 Q(2y_2h)=2y_2Mh,\qquad
 D_v(2y_2h)=2y_2D_vh+2v_2h.
\]
Therefore (\ref{nsbridge:nsb:Gamma}) is equivalent to
\begin{equation}
 D_vh=\nu_{\rm NS}Mh-\frac{v_2}{y_2}h+\frac{f^\vartheta}{r}.
 \label{nsbridge:nsb:h}
\end{equation}
The source profile uses precisely this relation: in its similarity variables,
$E$ is the azimuthal velocity profile and
$H_{\rm src}=\sqrt{2X}E$ is the scaled circulation profile (source equations
(4.3), (4.8)). Since $r=\sqrt{2qX}$ and
$u^\vartheta=q^{-A}E$, its physical circulation is
$\Gamma=q^{1/2-A}H_{\rm src}=q^{-h_{\rm src}}H_{\rm src}$,
where the source exponents satisfy $A=1/2+h_{\rm src}$. Here $q>0$ is the
geometric similarity coordinate and $h_{\rm src}>0$ is an exponent; neither
is a signed amplitude scale. This is the exact scale dictionary; it does not
identify $H_{\rm src}$ with the Boussinesq scalar amplitude.

The sign of the blowout amplitude must be kept separate from that positive
geometric coordinate. The axisymmetric equations have the exact involution
\[
 (u^r,u^\vartheta,u^z,p;f^r,f^\vartheta,f^z)
 \longmapsto
 (u^r,-u^\vartheta,u^z,p;f^r,-f^\vartheta,f^z).
 \label{nsbridge:nsb:sign}
\]
It sends $(\Gamma,h)$ to $(-\Gamma,-h)$, leaves $\xi$ and the meridional
flow unchanged, and preserves every displayed equation: the only quadratic
swirl term is $\partial_{y_1}(h^2)$, while the azimuthal equation is linear in
$(\Gamma,f^\vartheta)$. Consequently a source-selected positive profile
$E>0$ (source pp. 7--9) does not justify assuming that every blowout branch is
positive; its reflected branch has $E<0$, $\Gamma<0$ and the same blowup
magnitude. The exponent $q^{-h_{\rm src}}$ controls the magnitude in either
branch. The source's positive covariance coefficients and viscosity are
separate construction choices, not a sign claim about the velocity amplitude.

For comparison with the lane's controlled scalar equation, set
$B=h^2$. The globally smooth square equation, including its geometric and
viscous terms, is
\begin{equation}
 D_vB=\nu_{\rm NS}MB-2\nu_{\rm NS}
 |\nabla h|_Q^2-\frac{2v_2}{y_2}B+\frac{2h}{r}f^\vartheta,
 \qquad |\nabla h|_Q^2=h_{y_1}^2+2y_2h_{y_2}^2.
 \label{nsbridge:nsb:B}
\end{equation}
This follows by multiplying (\ref{nsbridge:nsb:h}) by $2h$ and using
$M(h^2)=2hMh+2|\nabla h|_Q^2$. At zeros of $h$ equation
(\ref{nsbridge:nsb:B}) remains smooth; dividing by $B$ is unnecessary and would lose
that domain. On a component where $B>0$, the equivalent expression is
$D_vB=\nu_{\rm NS}MB-\nu_{\rm NS}|\nabla B|_Q^2/(2B)
-2v_2B/y_2+2hf^\vartheta/r$.

The exact relation to the Boussinesq lane is consequently an operator
correspondence, not a scalar conjugacy. The lane's retained-viscosity scalar
row has $D_v\Theta=d\Delta\Theta+f_\Theta$ in Cartesian material
coordinates. A direct identification $\Theta=B$ would require, on the same
physical domain, both $M=\Delta$ after the nonlinear coordinate pullback and
$-2v_2B/y_2-2\nu_{\rm NS}|\nabla h|_Q^2+2hf^\vartheta/r$ to be absorbed into
one prescribed force. The first condition fails as an operator identity:
$M=\partial_{y_1}^2+2y_2\partial_{y_2}^2+4\partial_{y_2}$ has variable metric
and drift, whereas the lane's $\Delta$ is the Euclidean Cartesian Laplacian.
The geometric, gradient and force terms must therefore be retained or
calculated in the particular source state. Nonconstant profiles alone
would not prove that such terms never cancel.
The exact bridge is (\ref{nsbridge:nsb:Gamma})--(\ref{nsbridge:nsb:B}) with the
full residual terms retained. It neither claims that the lane's controls are
the source pulses nor changes the source's pressure reconstruction.

For the meridional velocity, the source's pulse construction supplies the
following further exact local dictionary (source pp. 11--15 and 74--77). If
$w$ is a divergence-free pulse and $\pi$ its pressure increment, then
\[
 R(u_B+w,p_B+\pi)=R(u_B,p_B)+L_{u_B}(w,\pi)
 +\nabla\!\cdot(w\otimes w),
\]
\[
 L_{u_B}(w,\pi)=\partial_tw+(u_B\cdot\nabla)w
 +(w\cdot\nabla)u_B-\nu_{\rm NS}\Delta w+\nabla\pi,
 \qquad \nabla\cdot w=0.
\]
For the moving local phase with wavevector $n$ and transverse amplitude
$t_m$, the source equation is
\[
 t_m'+\mathcal Kt_m+m^2d_N t_m+i k m n\,\pi_m=-f_m,
 \qquad n\cdot t_m=0,
 \qquad d_N=\varepsilon k^2|n|^2,
\]
with the source's moving shear matrix
$\mathcal K=\begin{pmatrix}0&-2F&0\\2F+RF_R&0&0\\G_R&0&0\end{pmatrix}$.
The transverse projection uses
\[
 A_\Phi=-\mathcal K+
 \frac{n(n^T\mathcal K-n'^T)}{|n|^2}.
\]
These are exact source pulse equations, including viscosity and pressure
projection. Their quadratic covariance is chosen to cancel the annular
stress; the source explicitly retains curl, cutoff, and correction interactions.
The Boussinesq lane's two-dimensional pair and its full profile residuals have
an exact algebraic form, but no identification with this three-dimensional
pulse equation follows without an additional embedding that supplies the
source's cylindrical geometry, moving plane, covariance cone and correction
cycle. That embedding is not asserted here.

The current bridge therefore proves a precise morphism of variables, operators,
forces and pulse residuals and records the exact obstruction to a direct scalar
identification. It establishes the next programme datum while preserving the
unproved status of any singular limit or third return.


\sloppy

\section{Scope and source equations}
This note records one local, exact comparison.  It does not identify the
coupled Boussinesq stages with the released Navier--Stokes construction and
makes no claim about a global Navier--Stokes solution.  The source is the
released PDF \texttt{openai\_navier\_stokes.pdf}, SHA-256
\texttt{8c8a94ad9ac824c8b605b9827cadf7beaca48bd10b380de3cfc872a2c37afa81}.
The source statements used here are Section 7, equations (7.2)--(7.8),
(7.12)--(7.22), and Section 9, equations (9.1)--(9.2), on PDF pages 74--75,
77--81, and 101--103. In the source notation, on one lifted pulse rectangle,
\begin{equation}
 \Phi=p\theta+p_zZ/\varepsilon+x_0R-v(pF+p_zG),\qquad
 n=\nabla_*\Phi,
\end{equation}
with
\begin{equation}
 n=(x_0-v(pF_R+p_zG_R),\ p/R,\ p_z-\varepsilon v(pF_Z+p_zG_Z)).
 \label{nsbridge:eq:nsource}
\end{equation}
For a nonzero integer harmonic $m$, the exact principal pulse equations are
\begin{equation}
 n\cdot t_m=0,\qquad
 t_m'+\mathsf Kt_m+m^2d_Nt_m+ikm n\pi_m=-f_m,
 \qquad d_N=\varepsilon k^2|n|^2,
 \label{nsbridge:eq:sourcepulse}
\end{equation}
where a prime is the pulse-coordinate derivative and
\begin{equation}
 \mathsf K=\begin{pmatrix}0&-2F&0\\2F+RF_R&0&0\\G_R&0&0\end{pmatrix},\qquad
 \mathsf A_\Phi=-\mathsf K+\frac{n(n^T\mathsf K-(n')^T)}{|n|^2}.
 \label{nsbridge:eq:sourceoperator}
\end{equation}
The source pressure coefficient is exactly
\begin{equation}
 \pi_m=\frac{i}{km|n|^2}\bigl(n\cdot\mathsf Kt_m-(n')\cdot t_m+n\cdot f_m\bigr).
 \label{nsbridge:eq:sourcepressure}
\end{equation}
For $m=1$, a real homogeneous pulse $t_1$ has the source energy identity
(7.22):
\begin{equation}
 \frac12(|t_1|^2)'=-g\cdot\bigl(t_{1,r}(t_{1,\theta},t_{1,z})\bigr)
 -\varepsilon k^2|n|^2|t_1|^2,
 \label{nsbridge:eq:sourceenergy}
\end{equation}
where $g=(RF_R,G_R)$ in the tangential coordinates.  Thus the source pulse
retains both shear transfer and viscous dissipation.

\section{The coupled physical pair}
For one activated layer of the coupled calculation, put
\begin{equation}
 a_c=-\frac{J\zeta\cdot G_{<q}}{Kr^2},\qquad b_c=K\zeta_1,\qquad
 \lambda_c^\Theta=\kappa K^2r^2,\qquad \lambda_c^\Omega=\varepsilon_c K^2r^2.
\end{equation}
Its exact retained-viscosity pair, including optional scalar and vorticity
controls, is
\begin{equation}
 \dot y=C_cy+c_c,\qquad y=\binom{\Theta}{\Omega},\qquad
 C_c=\begin{pmatrix}-\lambda_c^\Theta&a_c\\b_c&-\lambda_c^\Omega\end{pmatrix},\quad
 c_c=\binom{c_\theta}{c_\omega}.
 \label{nsbridge:eq:coupledpair}
\end{equation}
Here $G_{<q}$, $\zeta$, and $r$ evolve with the inherited background and phase
ODEs; in particular no fixed-background substitution is made.  Let $v$ be a
source pulse coordinate on an interval and let $q_v=\dd t/\dd v$ describe
its physical-time parametrization. The forward source parametrization has
$q_v>0$; the algebra below also retains a negative orientation without
changing the amplitude sign. Define $y'(v)=q_v(C_cy+c_c)$.
Here $K,r$ are the coupled frequency and phase length, not the source
carrier $k$ and cylindrical radius $R$. The original $A_0,\lambda_0,\mu$
and time scale $\nu=\sqrt{A_0}$ are unchanged. The source chart parameter
$\varepsilon$ is not identified with either physical diffusivity.

\section{Exact local residual map}
Let $B(v)\in\R^{3\times2}$ be a $C^1$ tangent frame with full column rank and
$n(v)^TB(v)=0$, so its range is exactly $n(v)^\perp$.  Let $B_\ell(v)$ be any
left inverse, $B_\ell B=I_2$; then $B B_\ell$ acts as the identity on
$n^\perp$.  Let $L(v)\in\R^{2\times2}$ be any $C^1$ coefficient map and embed
the coupled pair into a transverse source amplitude by
\begin{equation}
 t(v)=B(v)L(v)y(v).
 \label{nsbridge:eq:embedding}
\end{equation}
For an embedding with a non-tangent frame the transversality condition would be
\begin{equation}
 n(v)^TB(v)L(v)y(v)=0.
 \label{nsbridge:eq:transversecondition}
\end{equation}
For the tangent frame used here, the condition holds for every $L,y$.  Set
\begin{equation}
 H(v)=B_\ell(v)\bigl(\mathsf A_\Phi(v)B(v)-B'(v)\bigr).
 \label{nsbridge:eq:Hdef}
\end{equation}
Then the exact projected source residual of the embedded coupled pair is
\begin{align}
 \mathcal E_{L,c}:=B_\ell\proj_{n^\perp}
 \left[t'-\mathsf A_\Phi t+m^2d_Nt+f_m\right]
 &=\left[L'+q_vLC_c-HL+m^2d_NL\right]y+q_vLc_c \\
 &\quad+B_\ell\proj_{n^\perp}f_m.
 \label{nsbridge:eq:residualmap}
\end{align}
In particular, with $L=I$ and no added source, the defect is the explicit
matrix expression
\begin{equation}
 \mathcal E_{I,0}=\bigl(q_vC_c-H+m^2d_NI_2\bigr)y.
 \label{nsbridge:eq:identitydefect}
\end{equation}
This gives a local residual relation with every actual coefficient retained.

\begin{proof}
Differentiate \eqref{nsbridge:eq:embedding}:
$t'=B'L y+B L'y+B L y'$.  Since $y'=q_v(C_cy+c_c)$,
differentiate $n^TB=0$ to obtain $n^TB'=-(n')^TB$. Also
$n^T\mathsf A_\Phi B=-(n')^TB$, directly from
\eqref{nsbridge:eq:sourceoperator}. Thus $\mathsf A_\Phi B-B'$ has transverse
range and equals $BH$. Applying $B_\ell$ to the transverse terms and
$B_\ell\operatorname{proj}_{n^\perp}$ to the force proves
\eqref{nsbridge:eq:residualmap}. More explicitly,
\begin{equation}
 n^T(t'-\mathsf A_\Phi t+m^2d_Nt)=0,\qquad
 n^T(t'-\mathsf A_\Phi t+m^2d_Nt+f_m)=n^Tf_m.
\end{equation}
The unprojected force need not be transverse. The pressure below removes
precisely this remaining normal component.
\end{proof}

\section{Pressure reconstruction and full three-dimensional residual}
For any embedded transverse $t$ and arbitrary $f_m$, define
\begin{equation}
 \pi_{L,c}=\frac{i}{km|n|^2}
 \left(n\cdot\mathsf Kt-(n')\cdot t+n\cdot f_m\right).
 \label{nsbridge:eq:pressuremap}
\end{equation}
Then the full source residual satisfies the exact decomposition
\begin{align}
 \mathcal R_{L,c}
 &: =t'+\mathsf Kt+m^2d_Nt+ikm n\pi_{L,c}+f_m \\
 &=\proj_{n^\perp}\left(t'+\mathsf Kt+m^2d_Nt+f_m\right)
 =B\,\mathcal E_{L,c},
 \label{nsbridge:eq:fullresidual}
\end{align}
where the last equality uses the definition of $H$ and the transverse condition.
Thus the pressure removes precisely the normal component; it does not remove
$\mathcal E_{L,c}$, the moving-frame, shear, diffusion, or coupled-background
mismatch.  If one elects to force the source pulse to obey the released equation,
then the unique additional transverse forcing is
\begin{equation}
 f_{m,\mathrm{corr}}=-B\,\mathcal E_{L,c},
\end{equation}
added to the old $f_m$. The old $\mathcal E_{L,c}$ is evaluated before
this addition, so the new projected residual is zero.
Since $n^Tf_{m,\mathrm{corr}}=0$, pressure is unchanged.
This is an exact local
forced relation, not an identification of the physical constructions.

\section{Diffusion and energy comparison}
The source diffusion multiplier is $m^2d_N=m^2\varepsilon k^2|n|^2$ in
\eqref{nsbridge:eq:sourcepulse}.  The coupled pair has the two retained physical losses
$\lambda_c^\Theta$ and $\lambda_c^\Omega$, which are generally unequal to
$m^2d_N$.  Their signed discrepancy in \eqref{nsbridge:eq:identitydefect} is explicitly
\begin{equation}
 \bigl(m^2d_N-q_v\lambda_c^\Theta\bigr)\Theta,
 \qquad
 \bigl(m^2d_N-q_v\lambda_c^\Omega\bigr)\Omega
\end{equation}
when $H=0$ and $C_c$ is expanded entrywise; with general $H$ the exact
diagonal entries are those of $q_vC_c-H+m^2d_NI_2$.  No damping is dropped.
If one compares nonnegative loss magnitudes instead, the corresponding cost
gaps are $|m^2d_N-q_v\lambda_c^\Theta|$ and
$|m^2d_N-q_v\lambda_c^\Omega|$.
The required finite-interval coordinate map can be constructed explicitly.
Let $I=[v_0,v_1]$, and let $X_c,X_s$ be the identity-initial-value
fundamental matrices for $D_c=q_vC_c$ and $D_s=H-m^2d_NI_2$,
respectively. For either continuous matrix $D$, define
\[
 X_D(v)=I_2+\sum_{j\ge1}
 \int_{v_0\le w_j\le\cdots\le w_1\le v}
 D(w_1)\cdots D(w_j)\,\dd w_j\cdots\dd w_1 .
\]
With $M_D=\sup_I\|D\|$ and $T_I=v_1-v_0$, the $j$th integral has norm
at most $(M_DT_I)^j/j!$. Uniform convergence proves the integral equation
$X_D=I_2+\int_{v_0}^vDX_D$ and hence $X_D'=DX_D$.
The same iteration for $Z'=-ZD$, $Z(v_0)=I_2$ gives $ZX_D=I_2$
by differentiation, proving invertibility on $I$. For a specified
invertible $L_0$, set
\[
 L=X_sL_0X_c^{-1},\qquad L^{-1}=X_cL_0^{-1}X_s^{-1}.
\]
This has norm at most $\|L_0\|e^{(M_c+M_s)T_I}$; its inverse has the
same bound with $L_0^{-1}$. Differentiation proves
\begin{equation}
 L'=HL-q_vLC_c-m^2d_NL,
 \label{nsbridge:eq:intertwiner}
\end{equation}
so the homogeneous embedded pair has zero projected residual. This proves
an actual finite-interval conjugacy along the retained coefficient path.
It does not make the source phase,
pressure, cutoff, or nonlinear covariance equal to the coupled Boussinesq fields.
For general controls $c_c$, the forced term is $q_vLc_c$ in
\eqref{nsbridge:eq:residualmap}.

For an arbitrary embedded field (possibly complex), the exact energy balance
retains its full defect:
\begin{equation}
 \begin{split}
 \frac12(|t|^2)'={}&-\operatorname{Re}(\overline t^{\,T}\mathsf Kt)
 -m^2d_N|t|^2-\operatorname{Re}(\overline t^{\,T}f_m)\\
 &+\operatorname{Re}(\overline t^{\,T}B\mathcal E_{L,c}).
 \end{split}
\end{equation}
This follows by taking the real Hermitian product of
\eqref{nsbridge:eq:fullresidual} with $t$. Its pressure term vanishes because $n$ is
real and $n^Tt=0$. The $2F$ part of $\mathsf K$ is skew-symmetric, hence
its real work is zero; the shear work equals
$\operatorname{Re}(\overline{t_r}(g_\theta t_\theta+g_z t_z))$.
Only for a zero residual, zero force, $m=1$, and real $t$ does this reduce
to \eqref{nsbridge:eq:sourceenergy}. The real coupled pair has
\begin{equation}
 \frac12\frac{\dd}{\dd t}|y|^2
 =-\lambda_c^\Theta\Theta^2-\lambda_c^\Omega\Omega^2
 +(a_c+b_c)\Theta\Omega+c_\theta\Theta+c_\omega\Omega,
 \label{nsbridge:eq:coupledenergy}
\end{equation}
which is a two-component physical energy identity.  Equations
\eqref{nsbridge:eq:residualmap} and \eqref{nsbridge:eq:coupledenergy} exhibit exactly where the
source's three-dimensional shear flux and the coupled pair's two-dimensional
exchange differ.

\section{Proven scope}
The source uses a cylindrical, localized, divergence-free curl construction
(Section 7.4 and Section 9, PDF pages 85--87 and 102--103); the displayed
principal pulse equation is only one component of that construction.  The
comparison above proves an exact local map and its full pressure-projected
residual for every finite interval on which the denominators and frame left
inverse exist.  It preserves source viscosity, moving phase normal, pressure,
rotation/shear matrix, and all coupled physical diffusion coefficients.  It does
not prove that the coupled third-growth or steering stages produce the source's
pulse phase $\Phi$, covariance, cutoff tails, or smooth terminal force.  Those
remain separate bridge problems.

The sign conventions in this comparison are also explicit.  The quantities
$q_v>0$, $\varepsilon>0$, and $d_N=\varepsilon k^2|n|^2>0$ record the chosen
orientation of physical time and the positive viscous coefficient; they are
not assumptions on the sign of a pulse or blowout amplitude.  The equations
\eqref{nsbridge:eq:sourcepulse}--\eqref{nsbridge:eq:sourcepressure} are linear in $(t_m,f_m)$,
so the simultaneous replacement $(t_m,f_m,\pi_m)\mapsto
(-t_m,-f_m,-\pi_m)$ is an exact reflected branch.  Likewise, the local map
\eqref{nsbridge:eq:embedding} is sign-equivariant under $y\mapsto-y$ and
$c_c\mapsto-c_c$, with the coefficients fixed at the actual path.
This is a statement about the retained linear equations, not a sign symmetry
of the entire coupled nonlinear system.
The source covariance is quadratic, so reversing a wave leaves that
covariance unchanged. Reversing the full Navier--Stokes velocity need not
reverse its force, since $(u\cdot\nabla)u$ is even in $u$.
The full-flow map that reverses axisymmetric swirl is a spatial reflection;
its complete formula is proved in the accompanying covariance bridge.
No sign of a blowout or endpoint is inferred from positive geometric scales.


\section{Source objects and the finite calculation}
This calculation uses the preserved 165-page supplied manuscript
\emph{Finite Time Blowup for Navier--Stokes}, source equations
(7.23)--(7.42), (9.12)--(9.14). Its PDF SHA-256 is
\begin{center}\small
\texttt{8c8a94ad9ac824c8b605b9827cadf7beaca48bd10b380de3cfc872a2c37afa81}.
\end{center}
The source pulse fields, auxiliary rectangles, phases, and leading
covariance columns are the given objects; their global construction is
source input. Every finite map asserted below is proved here. The existing
coupled-stage proofs remain in the complete 102-page reader. The companion
pulse note proves the exact frame residual and a constructed finite-interval
intertwiner along its evolving coefficients. The source uses positive
$q(z,t)$, band scales $Q_\beta=2^{-\ell_\beta}$,
$\varepsilon_\beta=Q_\beta^h$, $A=1/2+h$; none is set to one here or
identified with the original coupled $A_0,\lambda_0,\mu,\nu=\sqrt{A_0}$
or its physical diffusivities.

All operations in this note are on a compact set inside $q>0$ and the
source active annulus. There are finitely many relevant bands and
locally finite source labels. Supports of newly prescribed inputs are
strictly inside the annulus, so no lower bound on a leading amplitude
at a vanishing shell edge is assumed. Endpoint and shell-edge estimates
are not inferred from compact-interior estimates.

\section{Physical curl of the actual retained amplitude}
Let $y=(\Theta,\Omega)^T$ be the actual smooth retained coupled solution,
$y'=q_v(C_cy+c_c)$, with the coefficients of the companion pulse note.
Use the smooth actual coefficients on the chosen slow/auxiliary patch;
all harmonic coefficients and cutoffs are independent of $\theta$ in
cylindrical components, and each carrier has nonzero integer angular
frequency. For its source tangent frame $B$, smooth coefficient map $L$,
and a fixed smooth real compact cutoff $\chi$ of this type, put
\[
 t_m=\chi BLy,\qquad
 \mathcal E_\chi=
 \chi\{[L'+q_vLC_c-HL+m^2d_NL]y+q_vLc_c\}+\chi'Ly,
 \quad H=B_\ell(\mathsf A_\Phi B-B').
\]
Here prime differentiates along the pulse coordinate; the other slow
and auxiliary derivatives remain in the spatial formulas below.
Smooth parameter dependence of the companion's constructed $L$ follows
directly from its integral equations on a fixed interval: for a nonzero
parameter multi-index $I$,
\[
 \partial^I X(v)=
 \int_{v_0}^v\left[D\,\partial^I X+
 \sum_{0<J\le I}\binom{I}{J}(\partial^JD)
                                 \partial^{I-J}X\right]\,\dd w .
\]
The initial parameter derivatives vanish for identity initial data.
Inductively the known forcing is smooth and bounded on the compact patch;
the same convergent ordered-integral iteration solves this equation.
This proves every finite mixed derivative of $X,X^{-1},L$.
Since $n^TB=0$, $n^Tt_m=0$. Differentiation proves
$t_m'-\mathsf A_\Phi t_m+m^2d_Nt_m=B\mathcal E_\chi$.
Thus even the pulse-cutoff term is explicit.

A physical version of the exact curl construction makes the intermediate
map to momentum residual unambiguous. For each chosen harmonic let
$\phi(x,t)=k_\beta m\Phi_\beta(x,t,Y(x,t))$ be its real evaluated phase,
let $\xi=\nabla_x\phi\ne0$, and let $a(x,t)=Q_\beta^{-A}t_m$ be the
physical transverse amplitude. The source coordinate change is
$r=\sqrt{Q_\beta}R$, $z=Q_\beta^D Z$, $1-t=Q_\beta T$,
with $D=1/2-h$ and normalized axial derivative
$D_z=\varepsilon_\beta\partial_Z$. It gives
$\xi=(k_\beta m/\sqrt{Q_\beta})n_\Phi$; all auxiliary chain derivatives
are included in $n_\Phi$. Hence $\xi\cdot a=0$. Define in Cartesian
components
\begin{equation}
 C=\frac{i\,\xi\times a}{|\xi|^2},\quad
 \mathcal A=C e^{i\phi},\quad
 \nabla\times\mathcal A=(a+r_a)e^{i\phi},\qquad
 r_a=\nabla\times C. \label{nsbridge:cb:curl}
\end{equation}
The vector triple-product identity gives
$i\xi\times C=a$, proving the last equality. It also shows that the
physical potential has exactly the source factor $Q_\beta^{1/2-A}$
relative to its chart coefficient
$i\,n_\Phi\times t_m/(k_\beta m|n_\Phi|^2)$.
Cartesian differentiation of $C$ includes derivatives of the cylindrical
basis; the same terms are the $R^{-1}$ terms of source (7.37)--(7.38).
Take real parts and sum the finite retained harmonics. This constructs
a smooth compact velocity $w_y$ with $\nabla\cdot w_y=0$.

Retain a finite smooth source state $(u_*,p_*)$, including its background,
actual leading waves and their curl corrections, and the real
physical pressure increments obtained from the companion pulse pressure
formula with $f_m=0$:
\[
 \pi_{\chi,m}=\frac{i}{k_\beta m|n|^2}
                  \bigl(n^T\mathsf Kt_m-(n')^Tt_m\bigr),\qquad
 \pi_y=\sum_{\beta,m}\operatorname{Re}
             (Q_\beta^{-2A}\pi_{\chi,m}e^{i\phi_{\beta,m}}).
\]
These are smooth and compact wherever $t_m$ is. This definition fixes the
pressure map without assuming the pulse residual vanishes.
Define the actual candidate residual
\begin{align}
 R_y={}&R_{\nu_{\rm NS}}(u_*+w_y,p_*+\pi_y), \nonumber\\
 R_{\nu_{\rm NS}}(u,p)
   :={}&\partial_tu+(u\cdot\nabla)u-\nu_{\rm NS}\Delta u+\nabla p .
 \label{nsbridge:cb:Ry}
\end{align}
This is a definition by explicit differentiated fields, not an assumed
source term. Its exact expansion is
\[
 R_y=R_{\nu_{\rm NS}}(u_*,p_*)
 +\partial_tw_y+(u_*\cdot\nabla)w_y+(w_y\cdot\nabla)u_*
 -\nu_{\rm NS}\Delta w_y+\nabla\pi_y
 +(w_y\cdot\nabla)w_y .
\]
Every cutoff, parent, phase, and curl derivative occurs in this formula.
Physical $\nu_{\rm NS}$ is retained and is distinct from every source
time-scaling convention. Define
$\Delta R_y=R_y-R_{\nu_{\rm NS}}(u_*,p_*)$.
We correct the new compact increment; this difference has compact support
since every added term does. The baseline residual remains explicitly.

The averaging needed next is the source extended-field operation, not a
spatial average of the evaluated field. Evaluate derivatives on the
extended coefficients first by the complete chain rule. In a fixed
chart set $\mathcal D_i=\partial_{x_i}
+\sum_\alpha(\partial_{x_i}Y_\alpha)\partial_{Y_\alpha}$ and similarly
$\mathcal D_t$. These are the pullbacks of commuting physical
derivatives, so $[\mathcal D_i,\mathcal D_j]=0$ on their lifted domain.
For smooth periodic coefficients their auxiliary derivative terms have
zero Haar average. In cylindrical components angular differentiation
also differentiates the basis. This defines an extended residual
$\widetilde R_y$ whose evaluation is exactly \eqref{nsbridge:cb:Ry}.
Average the tangential components of the increment at fixed slow variables:
\begin{equation}
 F_2=\av{\Delta\widetilde R_{y,\theta}},\qquad
 F_1=\av{\Delta\widetilde R_{y,z}},\qquad
 \av{g}=\int_{\mathbb T^2}\frac1{2\pi}
                    \int_0^{2\pi}g\,\dd\theta\,\dd Y . \label{nsbridge:cb:F}
\end{equation}
These $F_e$ are smooth and compact
in the chosen physical annulus. Equations \eqref{nsbridge:cb:curl}--\eqref{nsbridge:cb:F}
are the explicit intermediate map from the retained coupled pair to a
signed momentum residual. They do not identify it with the original
two-dimensional Boussinesq force.

\section{Radial moments, support, and the original chart scales}
For each fixed $(z,t)$, choose $0<r_-<r_+$ containing the compact
supports of both $F_e$ and the chosen bumps in the active annulus.
The source bump is
\[
 b_e(r,z,t)=q^{-(e+1)/2}\widehat b_e(r/\sqrt q),\qquad
 \int_0^\infty x^e\widehat b_e(x)\,\dd x=1,\qquad e=1,2,
\]
with interior support. The change $r=\sqrt q\,x$ proves
$\int r^eb_e\,\dd r=1$ exactly. Define
\begin{equation}
 M_e=\int_0^\infty r^eF_e(r)\,\dd r,\qquad
 \sigma_e(r)=-r^{-e}\int_0^r s^e(F_e(s)-b_e(s)M_e)\,\dd s .
 \label{nsbridge:cb:primitive}
\end{equation}
Then
\begin{equation}
 (\partial_r+e/r)\sigma_e=-F_e+b_eM_e,\qquad
 \int_0^\infty r^e(F_e-b_eM_e)\,\dd r=0. \label{nsbridge:cb:radial}
\end{equation}
Differentiate $r^e\sigma_e$ to prove the first identity; integrate and
use the bump normalization for the second. The integral is zero below
$r_-$ and above $r_+$, so $\sigma_e$ has compact smooth support in the
same enclosing interval. A homogeneous solution is $cr^{-e}$, and a
compact homogeneous solution has $c=0$; thus the primitive is unique
among compact solutions of its displayed equation. Moreover a compact
solution of $(\partial_r+e/r)\sigma=-F_e$ exists if and only if $M_e=0$:
necessity follows by integrating the weighted derivative, and
sufficiency follows from \eqref{nsbridge:cb:primitive}. The two moments are
therefore the exact surviving finite-dimensional terms.

For $C_e=\int_{r_-}^{r_+}s^e\,\dd s$,
$|M_e|\le C_e\|F_e\|_\infty$ and
\[
 \|\sigma_e\|_\infty\le r_-^{-e}C_e
        (\|F_e\|_\infty+\|b_e\|_\infty|M_e|).
\]
If $G_e=F_e-b_eM_e$, successive radial derivatives obey the exact recurrence
\[
 \partial_r^{j+1}\sigma_e=-\partial_r^jG_e
 -e\sum_{i=0}^j \binom{j}{i}(-1)^ii!r^{-i-1}
                              \partial_r^{j-i}\sigma_e .
\]
This proves every finite radial derivative bound inductively. Parameter
derivatives follow by differentiating the compact integral; all derivatives
of $q,b_e$ and the moving physical data are retained. On the present compact
set $q$ is bounded away from zero. No terminal uniform estimate is implied.

The source (9.12)--(9.14) uses physical
$F_e=Q^{-2A-1/2}E_e^*$, where
$E_2^*=\av{E_\theta}_Y$ and $E_1^*=\av{E_z}_Y$.
For our fields define $E_e^*=Q^{2A+1/2}F_e$ in that same convention.
With $r=\sqrt Q R$ and $Q$ fixed in a chart, the exact conversions are
\begin{gather}
 M_e=Q^{e/2-2A}M_e^*,\quad M_e^*=\int R^eE_e^*\,\dd R,\quad
 b_e^*(R)=Q^{(e+1)/2}b_e(\sqrt Q R),\nonumber\\
 \Sigma_e(R)=Q^{2A}\sigma_e(\sqrt Q R),\qquad
 (\partial_R+e/R)\Sigma_e=-E_e^*+b_e^*M_e^*. \label{nsbridge:cb:chart}
\end{gather}
Substitute $r=\sqrt Q R$ including its measure in the moment, and
differentiate $\Sigma_e$ to prove each factor. The input $F_e$ is already
auxiliary independent. Taking only an averaged moment of an unaveraged
$F(R,Y)$ would not suffice: $F=b_e(R)\cos(2\pi Y_1)$ has zero averaged
moment but a nonzero exterior primitive at each generic $Y$.

\section{The source covariance derivative and its signed right inverse}
At each slow point let $w_{\rm tan}=(w_\theta,w_z)$ and define
\[
 \Cov(w)=\av{w_rw_{\rm tan}},\qquad
 \Cross(w,v)=\av{w_rv_{\rm tan}+v_rw_{\rm tan}}.
\]
Direct expansion proves
$\Cov(w+v)=\Cov(w)+\Cross(w,v)+\Cov(v)$ and
$D\Cov(w)[v]=\Cross(w,v)$.

Use the exact source real pulses
$b_\sigma=\chi_g(\xi_g)\psi(v)t_\sigma^h\cos(k\Phi_\sigma)$,
$\sigma\in\{+,-\}$. They exclude the slow cutoff $\eta_\beta$ and
scalar amplitudes; $\chi_g$ is an auxiliary transverse cutoff.
The two signs have disjoint auxiliary supports. Put
$H_{\rm cov}=[\Cov(b_+)\mid\Cov(b_-)]$.
For slow scalar coefficients independent of both averages,
\[
 \Cov(a_+b_++a_-b_-)=H_{\rm cov}(a_+^2,a_-^2)^T.
\]
All cross products vanish by disjointness. The factor $1/2$ from
$\av{\cos^2(k\Phi_\sigma)}_\theta$ is already inside the columns.

On the source interior cone the fixed primary amplitudes satisfy
$a_\sigma=\sqrt{(H_{\rm cov}^{-1}T_{0,*})_\sigma}>0$,
$T_{0,*}=(Q/q)^{A+1/2}T_0$, and
$W_0=\sqrt\varepsilon\sum_\sigma a_\sigma b_\sigma$.
For any signed slow real vector $\Sigma$, define
\begin{equation}
 d_\Sigma=H_{\rm cov}^{-1}\Sigma/\varepsilon,\qquad
 \delta a_\sigma=(d_\Sigma)_\sigma/(2a_\sigma),\qquad
 L_\Sigma=\sqrt\varepsilon\sum_\sigma\delta a_\sigma b_\sigma .
 \label{nsbridge:cb:L}
\end{equation}
It follows exactly that
\begin{gather}
 \Cross(W_0,L_\Sigma)=\varepsilon H_{\rm cov}(2a\odot\delta a)=\Sigma,
 \label{nsbridge:cb:inverse}\\
 \Cov(W_0+L_\Sigma)
   =\varepsilon T_{0,*}+\Sigma+
       \varepsilon H_{\rm cov}(\delta a_+^2,\delta a_-^2)^T.
 \label{nsbridge:cb:quad}
\end{gather}
This proves a linear right inverse of the derivative. It is not an
exact inverse of the quadratic covariance map. There is no sign or
size restriction on $\Sigma$ in these identities; no square root of
$H_{\rm cov}^{-1}(T_{0,*}+\Sigma/\varepsilon)$ is taken.

With Euclidean and induced operator norms and
$a_{\min}=\min_\sigma|a_\sigma|>0$, they give the explicit bound
\begin{equation}
 \|\Cov(L_\Sigma)\|\le
 \frac{\|H_{\rm cov}\|\|H_{\rm cov}^{-1}\|^2}
              {4\varepsilon a_{\min}^2}\,\|\Sigma\|^2. \label{nsbridge:cb:bound}
\end{equation}
Indeed $\|\delta a\|\le\|H_{\rm cov}^{-1}\|\|\Sigma\|/
(2\varepsilon a_{\min})$ and
$\|(\delta a_\sigma^2)_\sigma\|_2\le\|\delta a\|_2^2$.
All higher derivative costs are explicit from Leibniz and
\[
 \partial^I H_{\rm cov}^{-1}
 =-H_{\rm cov}^{-1}
 \sum_{0<J\le I}\binom{I}{J}(\partial^JH_{\rm cov})
                           \partial^{I-J}H_{\rm cov}^{-1},
\]
and, for any nonzero scalar $a$,
$\partial^I(a^{-1})=-a^{-1}\sum_{0<J\le I}\binom{I}{J}
(\partial^Ja)\partial^{I-J}(a^{-1})$.
Together with \eqref{nsbridge:cb:L} these recurrences give finite bounds for
each order without assuming constant amplitudes or constant columns.

\section{Assembly, the actual covariance change, and residual cancellation}
For the physical target $\sigma=(\sigma_2,\sigma_1)$ from
\eqref{nsbridge:cb:primitive}, use the source squared slow partition
$\sum_\beta\eta_\beta^2=1$ on its support. Set
\begin{align}
 W_{0,\rm phys}&=\sum_\beta Q_\beta^{-A}\eta_\beta W_{0,\beta},\nonumber\\
 v_{\rm pr}&=\sum_\beta Q_\beta^{-A}\eta_\beta
                        L_\beta(Q_\beta^{2A}\sigma).
 \label{nsbridge:cb:assembly}
\end{align}
Overlapping slow labels have disjoint auxiliary supports. Taking cross
covariance and using \eqref{nsbridge:cb:inverse} gives
\[
 \Cross(W_{0,\rm phys},v_{\rm pr})
 =\sum_\beta Q_\beta^{-2A}\eta_\beta^2Q_\beta^{2A}\sigma=\sigma.
\]
Every band scale remains distinct. Likewise the leading physical factor is
$Q^{-2A}\varepsilon(Q/q)^{A+1/2}
=Q^{-A+h+1/2}q^{-A-1/2}=q^{-A-1/2}$ because $A=1/2+h$.

Apply the exact curl \eqref{nsbridge:cb:curl} to each coefficient of
$v_{\rm pr}$, including the slow cutoffs before curling. Write the
resulting divergence-free increment as $v_s=v_{\rm pr}+r_s$.
Let $t_{s,m}$ denote the chart transverse harmonic coefficients of
$\eta_\beta L_\beta(Q_\beta^{2A}\sigma)$, before multiplication by
$Q_\beta^{-A}$. For each such coefficient choose
$\pi_{s,m}=i(n^T\mathsf Kt_{s,m}-(n')^Tt_{s,m})/
(k_\beta m|n|^2)$ and sum
$\operatorname{Re}(Q_\beta^{-2A}\pi_{s,m}e^{i\phi_{\beta,m}})$
to define the smooth compact $\pi_s$. Its full gradients remain below.
Let the actual old wave $w$ consist of the wave part of $u_*$
plus the candidate $w_y$, including all existing curl corrections, and write
$w=W_{0,\rm phys}+e_w$. Define tensor covariance
$\mathcal T(w)=\av{w\otimes w}$ and
$\mathcal B_T(u,v)=\av{u\otimes v+v\otimes u}$.
Expansion, with $e_w=w-W_{0,\rm phys}$, proves
\begin{equation}
 \Delta\mathcal T=
 \mathcal B_T(W_{0,\rm phys},v_{\rm pr})
 +\mathcal B_T(e_w,v_{\rm pr})
 +\mathcal B_T(w,r_s)+\mathcal T(v_s).
 \label{nsbridge:cb:fullcov}
\end{equation}
Its first radial--tangential pair is $\sigma$.
All other entries and the remaining three tensors are retained.
For example their radial--tangential pair has bound
\[
 2\|e_w\|_{L^2_{\theta,Y}}\|v_{\rm pr}\|_{L^2_{\theta,Y}}
 +2\|w\|_{L^2_{\theta,Y}}\|r_s\|_{L^2_{\theta,Y}}
 +\|v_s\|_{L^2_{\theta,Y}}^2 ,
\]
by Cauchy--Schwarz. The compact-set estimates above and the explicit
curl differentiate to bound all these terms at each finite order.

For clarity the complete residual after adding $v_s$ and its selected
pressure $\pi_s$ is, without any averaging,
\begin{align}
 R_{\nu_{\rm NS}}(u+v_s,p+\pi_s)
 =R_{\nu_{\rm NS}}(u,p)
 &+\partial_tv_s+(u\cdot\nabla)v_s+(v_s\cdot\nabla)u \nonumber\\
 &-\nu_{\rm NS}\Delta v_s+\nabla\pi_s
 +(v_s\cdot\nabla)v_s . \label{nsbridge:cb:residual}
\end{align}
This follows by expanding the product and retains viscosity and every
pressure term. For a symmetric tensor $S$ the averaged tangential
divergence in physical cylindrical coordinates is
\[
 (\av{\nabla\cdot S})_\theta
   =(\partial_r+2/r)\av{S_{r\theta}}+\partial_z\av{S_{z\theta}},
 \qquad
 (\av{\nabla\cdot S})_z
   =(\partial_r+1/r)\av{S_{rz}}+\partial_z\av{S_{zz}} .
\]
To verify, write the Cartesian divergence in the orthonormal cylindrical
basis, using $\partial_\theta e_r=e_\theta$,
$\partial_\theta e_\theta=-e_r$; periodic angular derivatives average
to zero and the two basis terms give $2/r$ and $1/r$ respectively.
The auxiliary chain terms average to zero by periodicity, but the
physical slow derivatives remain.

Let $U=u_*+w_y$, $P_y=p_*+\pi_y$, and set
\[
 \mathcal L=
 \partial_tv_s+(U\cdot\nabla)v_s+(v_s\cdot\nabla)U
 -\nu_{\rm NS}\Delta v_s+\nabla\pi_s .
\]
For a fully explicit averaged remainder, write $U=u_{\rm mean}+w$,
where $u_{\rm mean}$ is its angular mean, and put
$\mathcal V=\mathcal L-(w\cdot\nabla)v_s-(v_s\cdot\nabla)w$.
Then the average of \eqref{nsbridge:cb:residual} is
$\av{R_y}+\av{\mathcal V}+\nabla\cdot\Delta\mathcal T$.
Let $S_0=\mathcal B_T(W_{0,\rm phys},v_{\rm pr})$ and
$S_{\rm rem}=\mathcal B_T(e_w,v_{\rm pr})
 +\mathcal B_T(w,r_s)+\mathcal T(v_s)$.
Consequently the radial divergence of the first pair in
\eqref{nsbridge:cb:fullcov} changes the prescribed averaged residual $F_e$ into
$b_eM_e$ by \eqref{nsbridge:cb:radial}. The axial tensor divergences, the other
three tensors, the exact linear terms and the pressure terms of
\eqref{nsbridge:cb:residual} remain in the new residual. Equivalently the new
averaged tangential residual is exactly
\begin{align*}
 \av{R_{\rm new}}_{\rm tan}
 ={}&\av{R_{\nu_{\rm NS}}(u_*,p_*)}_{\rm tan}
 +(b_2M_2,b_1M_1)+\av{\mathcal V}_{\rm tan}\\
 &+\partial_z((S_0)_{z\theta},(S_0)_{zz})
   +(\nabla\cdot S_{\rm rem})_{\rm tan}.
\end{align*}
Every term here is defined by the displayed differentiated fields and
tensors. None is assumed to vanish.
The map from $y$ through \eqref{nsbridge:cb:F}, \eqref{nsbridge:cb:primitive},
\eqref{nsbridge:cb:L}, and \eqref{nsbridge:cb:assembly} is now an actual composite
construction. It retains the signed residual produced by the coupled
candidate and the source's original physical and chart units.

\section{Signs and the exact full-flow reflection}
The maps $F_e\mapsto M_e\mapsto\sigma_e\mapsto L_\Sigma$ are linear
and reverse sign with their inputs; $\Cov(L_\Sigma)$ is even.
The two primary amplitudes may also be replaced independently by
$s_\sigma a_\sigma$, $s_\sigma\in\{-1,1\}$, provided the same fixed
choice is used in each denominator: the cross identity is unchanged.
This does not negate the nonlinear Navier--Stokes equation.

Here is the exact full-flow map. In Cartesian coordinates let
$O=\operatorname{diag}(1,-1,1)$ and define
\[
 u^\sharp(x,t)=O u(Ox,t),\quad
 p^\sharp(x,t)=p(Ox,t),\quad f^\sharp(x,t)=O f(Ox,t).
\]
The chain rule gives
$\nabla u^\sharp=O(\nabla u)(Ox,t)O$,
$\Delta u^\sharp=O(\Delta u)(Ox,t)$, and
$(u^\sharp\cdot\nabla)u^\sharp
 =O((u\cdot\nabla)u)(Ox,t)$.
Thus $\nabla\cdot u^\sharp=0$ and
$R_{\nu_{\rm NS}}(u^\sharp,p^\sharp)=O R_{\nu_{\rm NS}}(u,p)(Ox,t)$
with the same positive viscosity.
In cylindrical components this is
\[
 (u_r^\sharp,u_\theta^\sharp,u_z^\sharp)(r,\theta,z,t)
 =(u_r,-u_\theta,u_z)(r,-\theta,z,t),
\]
and the force transforms identically. Pressure reflects in its argument.
For an axisymmetric field this reverses exactly its swirl; meridional
components remain fixed. The entire momentum tensor is mapped to
$O\mathcal T O^T$, so its radial--tangential pair transforms by
$(T_{r\theta},T_{rz})\mapsto(-T_{r\theta},T_{rz})$ after angular
averaging. Vorticity transforms as
$\omega^\sharp(x,t)=\det(O)\,O\omega(Ox,t)$, obtained by applying
the same derivative identities to the Levi-Civita formula for curl.
The map is an involution, preserves the kinetic-energy integral
(by $|\det O|=1$), compact support, smoothness, and blowup magnitude
along the reflected path. It preserves zero initial data when present.
It proves transport of any established axisymmetric signed swirl
divergence to the opposite sign; it does not itself prove an endpoint.

The source five-moment correction, auxiliary mean inverse, and infinite
summation are subsequent source operations. Their full estimates are
not re-proved here. The earlier draft formulas for them are withdrawn:
in particular the fifth moment row is
$\int(2RG\gamma_d-RV\Delta v)\,\dd R=-J_z$, without an extra outer
$R$, and its reserved profile in a $Q$ chart includes $(Q/q)^A$.
The finite map proved above retains the two moments and all other
residual terms, so it establishes neither a third shooting return nor
an infinite modified viscous sequence.


\section{The released result and its actual signed path}
\label{nsprop:source}
The source in this section is the supplied 166-page edition of
\emph{Finite Time Blowup for Navier--Stokes}, whose SHA-256 is
\begin{center}\footnotesize
\texttt{0e779481c4da40bd28d1e642e1d8ca57447d129610df28dfa5a11e9af8ae228f}.
\end{center}
Its Theorem 1.1, on page 1, states the following result for the standard
three-dimensional incompressible Navier--Stokes equation. For every
fixed physical viscosity $\nu_{\rm NS}>0$, the source constructs
\[
 f\in C^\infty_c(\mathbb R^3\times(0,\infty);\mathbb R^3),
 \qquad u,p\in C^\infty(\mathbb R^3\times[0,1)),
\]
with a fixed compact spatial support $K$ for $u$ and $p$, such that
\begin{equation}
 \begin{gathered}
 \partial_tu+(u\cdot\nabla)u-\nu_{\rm NS}\Delta u+\nabla p=f,
 \quad \nabla\cdot u=0,\quad u(\cdot,0)=0,\\
 \quad \sup_{t<1}\|u(t)\|_2<\infty,\quad
 \limsup_{t\uparrow1}\|u(t)\|_\infty=\infty .
 \end{gathered}
 \label{nsprop:source_theorem}
\end{equation}
It also excludes a global smooth solution with the same force and datum
and uniformly bounded kinetic energy. Its Corollary 10.6, pages 125--126,
gives a periodic construction on $\mathbb R^3/\mathbb Z^3$ with smooth
time-compact force, zero initial datum and no global smooth continuation.
The source identifies these with alternatives (C) and (D) in the
referenced Millennium problem statement. The force is nonzero, as the
source explicitly notes on page 124. These are statements of the cited
source theorem; the finite algebra checks in this lane do not constitute
an independent verification of the complete 166-page proof.

The source proves more than an unsigned norm divergence. In its
viscosity-one construction, equations (10.20)--(10.21), page 124, give
\begin{equation}
 \begin{split}
 A&=\tfrac12+h,\qquad 0<h<1/100,\qquad X_{\rm in}>0,\quad e_0>0,\\
 x_\tau&=(\sqrt{2X_{\rm in}\tau},0,0),\qquad t_\tau=1-\tau,\\
 u_\theta(x_\tau,t_\tau)
   &=\tau^{-A}(e_0+\epsilon(\tau)),\qquad
      |\epsilon(\tau)|\le C\tau^{2h}\quad(0<\tau<\tau_0).
 \end{split}\label{nsprop:positive_path}
\end{equation}
Here $C,\tau_0$ are finite constants supplied by the source asymptotic;
$e_0$ and $X_{\rm in}$ retain their source values. If $C>0$, choose
$0<\tau_1\le\tau_0$ so that $C\tau_1^{2h}\le e_0/2$; if $C=0$,
take $\tau_1=\tau_0$. It follows directly that
\[
 u_\theta(x_\tau,t_\tau)\ge \frac{e_0}{2}\tau^{-A}
       \longrightarrow+\infty\qquad(0<\tau<\tau_1).
\]
The observation point moves to the origin. This assertion does not
replace that path by the fixed point $x=0$. On the path, $q=\tau>0$,
and at its cylindrical angle $\theta=0$ one has $e_\theta=e_2$.
The geometric domain $q>0$, the positive exponent $h$, and the signed
velocity value are three retained parts of the same explicit formula.

\section{Exact propagation at retained viscosities and orientations}
\label{nsprop:map}
For any smooth fields $u,p$ define the physical momentum residual
\[
 R_{\nu_b}(u,p)=\partial_tu+(u\cdot\nabla)u-\nu_b\Delta u+\nabla p.
\]
Let $\nu_b>0$ and $\nu_t>0$ be the base and target physical viscosities,
and retain
\[
 \rho=\sqrt{\nu_t/\nu_b},\qquad O^TO=OO^T=I,\qquad a\in\mathbb R^3.
\]
Define the affine spatial coordinate $y=O^T(x-a)/\rho$ and the fields
\begin{equation}
 u^\sharp(x,t)=\rho O u(y,t),\quad
 p^\sharp(x,t)=\rho^2p(y,t),\quad
 f^\sharp(x,t)=\rho O f(y,t).
 \label{nsprop:affine_map}
\end{equation}
No time coordinate is changed. These formulas give a bijection on smooth
field triples, with inverse
\[
 u(y,t)=\rho^{-1}O^Tu^\sharp(a+\rho Oy,t),\quad
 p(y,t)=\rho^{-2}p^\sharp(a+\rho Oy,t),\quad
 f(y,t)=\rho^{-1}O^Tf^\sharp(a+\rho Oy,t).
\]
For the Jacobian convention $(Du)_{ij}=\partial_j u_i$, differentiation
gives, without dropping an orientation or viscosity,
\begin{align*}
 Du^\sharp&=O(Du)(y,t)O^T,&
 \nabla\cdot u^\sharp&=(\nabla\cdot u)(y,t),\\
 \partial_tu^\sharp&=\rho O\partial_tu(y,t),&
 (u^\sharp\cdot\nabla_x)u^\sharp&=\rho O((u\cdot\nabla_y)u)(y,t),\\
 \Delta_xu^\sharp&=\rho^{-1}O\Delta_yu(y,t),&
 \nabla_xp^\sharp&=\rho O\nabla_yp(y,t).
\end{align*}
The nonlinear identity follows by multiplying $Du^\sharp$ by
$u^\sharp$ and using $O^TO=I$. The Laplacian identity follows by
contracting the two spatial derivative indices and using orthogonality.
Since $\nu_t=\rho^2\nu_b$, these identities prove
\begin{equation}
 R_{\nu_t}(u^\sharp,p^\sharp)-f^\sharp
       =\rho O\bigl(R_{\nu_b}(u,p)-f\bigr)(y,t).
 \label{nsprop:residual_map}
\end{equation}
In particular the map preserves the exact equation as well as each
nonzero residual with its exact factor.

For vorticity, the cross-product identity
$(Ov)\times(Ow)=\det(O)O(v\times w)$ gives
\begin{equation}
 \omega^\sharp(x,t)=\det(O)O\omega(y,t).
 \label{nsprop:vorticity}
\end{equation}
One proof of this identity takes its scalar product with $Oz$:
both sides give $\det(O)\det[v,w,z]$ for every $z$. Applying it to
the contracted first derivatives proves the vorticity formula; the
factor $\rho$ in velocity cancels the inverse factor in a derivative.

Spatial support becomes $a+\rho OK$. The Jacobian has absolute
determinant $\rho^3$, for both signs of $\det O$. Hence, for
$1\le p<\infty$,
\begin{equation}
 \|u^\sharp(t)\|_p=\rho^{1+3/p}\|u(t)\|_p,\qquad
 \|u^\sharp(t)\|_\infty=\rho\|u(t)\|_\infty,\qquad
 \|u^\sharp(t)\|_2^2=\rho^5\|u(t)\|_2^2 .
 \label{nsprop:norms}
\end{equation}
The Frobenius norm of $O(Du)O^T$ equals that of $Du$. Consequently,
with integrals over the same time interval,
\[
 \nu_t\int\|\nabla u^\sharp\|_2^2\,dt
     =\rho^5\nu_b\int\|\nabla u\|_2^2\,dt .
\]
For arbitrary directions $v_1,\ldots,v_j$, the exact force derivative is
\[
 D_x^j\partial_t^m f^\sharp[v_1,\ldots,v_j](x,t)
 =\rho^{1-j}O
   (D_y^j\partial_t^m f)[O^Tv_1,\ldots,O^Tv_j](y,t).
\]
Thus smoothness, compact space-time support and vanishing of every
derivative at a terminal point are preserved. The initial datum remains
zero. A global smooth bounded-energy competitor for the transformed
datum and force would, under the displayed inverse, be a competitor
for the base datum and force, because its energy is multiplied by the
fixed factor $\rho^{-5}$. This proves propagation of the source
noncontinuation conclusion through the map.

\section{The source positive branch and the exact negative branch}
\label{nsprop:branches}
Apply the map to the source field in \eqref{nsprop:positive_path}.
Here $\nu_b=1$ specifies the source object being used; physical target
viscosity $\nu_t$ remains arbitrary and $\rho=\sqrt{\nu_t}$ is retained.
Take $a=0$ and the two particular orthogonal matrices
\[
 O_+=I,\qquad O_-=\operatorname{diag}(1,-1,1).
\]
Both fix $e_1$, so they give the same spatial observation path
$x^\sharp_\tau=\rho x_\tau$ and the same time $1-\tau$.
Since $O_+e_2=e_2$ and $O_-e_2=-e_2$, their exact signed values are
\begin{equation}
 u^\sharp_{\pm,\theta}(\rho x_\tau,1-\tau)
     =\pm\rho\,\tau^{-A}(e_0+\epsilon(\tau))
        \longrightarrow\begin{cases}+\infty,&+,\\-\infty,&-.\end{cases}
 \label{nsprop:two_paths}
\end{equation}
All support, energy, force and viscosity statements in the preceding
section apply to both triples. For the negative branch the force is
$\rho O_-f(O_-x/\rho,t)$. Equality of this force with the positive
branch force is not needed and is not asserted.

For a general, possibly nonaxisymmetric field the reflection alone has
the exact cylindrical component formula
\[
 (u^r,u^\theta,u^z)^\sharp(r,\theta,z,t)
    =(u^r,-u^\theta,u^z)(r,-\theta,z,t).
\]
Indeed $O_-e_r(-\theta)=e_r(\theta)$,
$O_-e_\theta(-\theta)=-e_\theta(\theta)$ and $O_-e_z=e_z$.
This proves the formula even in the presence of all oscillatory
corrections. For the axisymmetric background the argument $-\theta$
can be dropped because its components are independent of $\theta$.
For a velocity covariance tensor the exact Cartesian transformation is
$T^\sharp(x,t)=\rho^2 O T(y,t)O^T$: expand the outer product of
$\rho Ow$ with itself and use linearity of the retained average.
For its angularly averaged symmetric tensor, the reflection basis
calculation at $\rho=1$ gives
\[
 (T_{r\theta},T_{rz})^\sharp=(-T_{r\theta},T_{rz}).
\]
Under the additional viscosity dilation its components have factor
$\rho^2$ at the corresponding points. Thus the two covariance columns
and their target transform together. Positive coefficients in a source
representation $T=c_1v_1+c_2v_2$, $c_i>0$, remain the same coefficients
of the transformed columns; no change of coefficient sign is required
to realize the negative-swirl branch.

The source's periodic construction also has this reflection: $O_-$
preserves $\mathbb Z^3$, so $[x]\mapsto[O_-x]$ is well defined on
$\mathbb R^3/\mathbb Z^3$. The derivative identities apply to periodic
lifts. This gives both signed orientations of the source Corollary 10.6.
Arbitrary $O\in O(3)$ need not preserve this particular lattice; only
the lattice-preserving reflection just specified is used for this
periodic assertion.

\section{Propagation into the coupled calculation}
\label{nsprop:programme}
The preceding sections use the released whole-space and periodic
constructions at the precise places where those source results apply.
The original coupled Boussinesq work is due to Levent Alp\"oge and
Tristan Buckmaster. This lane's retained coupled system, its finite
controlled stages and its residual maps remain the separate objects
defined in the preceding parts of this reader.

For every compact candidate $(u_*+w_y,p_*+\pi_y)$ in the covariance
part, \eqref{nsprop:residual_map} propagates its full computed residual.
Linearity of the transformation and the exact nonlinear identity give
in particular
\[
 \Delta R_y^\sharp(x,t)=\rho O\Delta R_y(y,t).
\]
Thus the finite candidate's cutoff forces, curl errors and covariance
remainders all have an exact transformed counterpart. A nonzero
remainder is carried to a nonzero remainder, because the map is
invertible. The source's residual flatness at the singular point belongs
to its completed local source construction; its extended smooth force
need not vanish on the entire time-one slice. The next section calculates the
surviving moments of the actual finite candidate, so that propagation
continues through its computed terms rather than assuming that its
different coupled coefficients inherit the source's infinite estimates.

The earlier suggestion of a negative blowout imposes no condition on
the programme. Formula \eqref{nsprop:positive_path} is the source-selected
sign; formula \eqref{nsprop:two_paths} proves the exact relation of its
two orientations. Both formulas retain a positive physical viscosity
and the source coordinate domain $q>0$.


\section{The two moments of the actual compact curl increment}
\label{nsmom:setup}
The objects in this calculation are the finite smooth source state
$(u_*,p_*)$ and the actual added curl field $v=w_y$ with its selected
pressure $\pi=\pi_y$ from the signed covariance bridge. Their physical
residual increment is
\begin{equation}
 \Delta R=\partial_t v+(u_*\cdot\nabla)v+(v\cdot\nabla)u_*
 +(v\cdot\nabla)v-\nu_{\rm NS}\Delta v+\nabla\pi .
 \label{nsmom:increment}
\end{equation}
The physical viscosity $\nu_{\rm NS}$ is unchanged. Both velocities are
divergence free: the source state uses its actual divergence-free
background and curl increments, and $v$ is the curl of its specified
compact potential. No residual equation for $u_*$ is required in
\eqref{nsmom:increment}; its residual has been subtracted exactly.
All radial supports involving $v$ or $\pi$ lie in a fixed compact
subinterval of $(0,\infty)$ on the compact space--time patch at issue.

Here the source extended field is indispensable. Its absolute auxiliary
coordinate is $Y\in\mathbb T^2=\mathbb R^2/\mathbb Z^2$, and physical
evaluation uses precisely source (6.2)--(6.4):
\begin{gather}
 J_g=\begin{pmatrix}3&1\\1&5\end{pmatrix},\quad
 \Lambda_g=4-\sqrt2,\quad T_g=4+\sqrt2,\quad b_g=\sqrt2-1,\nonumber\\
 \mathbf v_r=(1,-b_g),\quad\mathbf v_t=(b_g,1),\quad
 \rho_g=\frac{\log\Lambda_g}{\log T_g},\quad
 \kappa_s=10^{-5},\quad d_r=2((1+h)\rho_g-h\kappa_s),\nonumber\\
 Y(r,t)=\mathbf v_r r^{d_r}+\mathbf v_t t\pmod{\mathbb Z^2},\qquad
 \mathscr D_r=\partial_r+d_r r^{d_r-1}\mathbf v_r\cdot\partial_Y,
 \quad\mathscr D_t=\partial_t+\mathbf v_t\cdot\partial_Y.
 \label{nsmom:phase}
\end{gather}
These are the source's original quantities, not the coupled physical
frequency or diffusivity. Angular and axial physical derivatives are
$\partial_\theta$ and $\partial_z$, because $Y(r,t)$ depends on neither
variable. Every field below is first differentiated using these lifted
physical derivatives. A tilde is suppressed for this extended
representation. Evaluation at \eqref{nsmom:phase} recovers the physical
field and its physical derivatives by the chain rule.

The incompressibility needed below holds at every independent $Y$.
Indeed the actual finite source construction retains the form
$u_* =\operatorname{curl}_{\mathscr D}\mathcal A_*+
B_*e_\theta$ with $\partial_\theta B_*=0$, and the added field is
$v=\operatorname{curl}_{\mathscr D}\mathcal A_y$. The retained source
potentials here include its finite mean and wave curl increments.
The lifted derivatives in \eqref{nsmom:phase}, together with
$\partial_\theta$ and $\partial_z$, commute, and
$\mathscr D_r r=1$. Substitution in the cylindrical curl and
divergence formulas therefore gives
$\operatorname{div}_{\mathscr D}\operatorname{curl}_{\mathscr D}
\mathcal A=0$ identically in $(r,\theta,z,t,Y)$: each mixed derivative
cancels its opposite, including the derivative of $r^{-1}$.
Also $\operatorname{div}_{\mathscr D}(B_*e_\theta)
=r^{-1}\partial_\theta B_*=0$.
Thus extended incompressibility is proved from the actual retained
potential representation, rather than inferred from its restriction to
the physical phase map.

Use the normalized averages
\[
 \langle a\rangle_\theta=\frac1{2\pi}\int_0^{2\pi}a\,d\theta,
 \qquad \langle a\rangle_Y=\int_{\mathbb T^2}a\,dY,
 \qquad\langle a\rangle=\langle\langle a\rangle_\theta\rangle_Y,
 \qquad\int_{\mathbb T^2}1\,dY=1.
\]
They act on cylindrical components before physical evaluation, exactly
as source (3.7)--(3.9). All claimed integral identities below concern
these specified averages; the exact relation to evaluated moments is
also proved below.

\subsection{The actual harmonic curl has zero angular mean}
\label{nsmom:zero_mean}
For each retained source harmonic, its extended potential and pressure
have cylindrical components of the form
\[
 C_j(r,z,t,Y)e^{i(n_j\theta+\psi_j(r,z,t,Y))},\qquad
 P_j(r,z,t,Y)e^{i(n_j\theta+\psi_j(r,z,t,Y))},\qquad
 n_j=k_{\beta_j}m_jp_j\in\mathbb Z\setminus\{0\}.
\]
The coefficients include the actual cutoff $\chi$, frame $B$, map $L$,
coupled solution $y=(\Theta,\Omega)^T$, band powers and pressure formula.
Their cylindrical components are independent of $\theta$. In particular,
no sign of $\Theta$, $\Omega$ or the real part of a harmonic is prescribed.
For a vector potential $\mathcal A$ the exact cylindrical curl is
\[
 (\operatorname{curl}\mathcal A)_r=r^{-1}\partial_\theta\mathcal A_z
                         -\partial_z\mathcal A_\theta,
 \quad
 (\operatorname{curl}\mathcal A)_\theta=\partial_z\mathcal A_r
                         -\mathscr D_r\mathcal A_z,
 \quad
 (\operatorname{curl}\mathcal A)_z=(\mathscr D_r+r^{-1})\mathcal A_\theta
                         -r^{-1}\partial_\theta\mathcal A_r.
\]
These formulas include the derivatives of the cylindrical basis.
Each displayed operation retains the same factor $e^{in_j\theta}$:
the coefficients of $\mathscr D_r$ are independent of $\theta$, and
$\partial_\theta$ multiplies that factor by $in_j$. Integration of the
factor over the angular circle gives zero. Taking real parts and the
actual finite harmonic sum proves
\begin{equation}
 \langle v_r\rangle_\theta=\langle v_\theta\rangle_\theta
 =\langle v_z\rangle_\theta=\langle\pi\rangle_\theta=0
 \quad\hbox{at every }(r,z,t,Y).
 \label{nsmom:mean_zero}
\end{equation}
This includes every curl and cutoff correction, not merely the
principal amplitude. The same identities hold after physical
evaluation because $Y(r,t)$ is independent of $\theta$.

For every smooth periodic coefficient $a$, integration of its auxiliary
derivatives gives
\begin{equation}
 \langle\mathscr D_r a\rangle=\partial_r\langle a\rangle,
 \quad\langle\mathscr D_t a\rangle=\partial_t\langle a\rangle,
 \quad\langle\partial_z a\rangle=\partial_z\langle a\rangle,
 \quad\langle\partial_\theta a\rangle=0.
 \label{nsmom:commute}
\end{equation}
There is no omitted variable coefficient in this assertion:
$d_r r^{d_r-1}\mathbf v_r$ is independent of $Y$. Applying the first
identity again proves
$\langle\mathscr D_r^2a\rangle=\partial_r^2\langle a\rangle$.
Explicitly, if $c(r)=d_r r^{d_r-1}\mathbf v_r$, then
\[
 \mathscr D_r^2 a=\partial_r^2a+2c\cdot\partial_Y\partial_r a
 +c'\cdot\partial_Ya+(c\cdot\partial_Y)^2a;
\]
each of the three additional terms has zero auxiliary average.

\subsection{Exact removal of the linear and mean-flow terms}
\label{nsmom:linear}
Put
\[
 \bar u_*:=\langle u_*\rangle_\theta,
 \qquad w_*:=u_*-\bar u_*.
\]
The field $\bar u_*$ may retain auxiliary dependence; none is discarded.
Since it is independent of $\theta$, \eqref{nsmom:mean_zero} gives
\[
 \langle\bar u_*\otimes v+v\otimes\bar u_*\rangle=0.
\]
Define the actual symmetric cross-and-quadratic tensor
\begin{equation}
 S:=\langle w_*\otimes v+v\otimes w_*+v\otimes v\rangle.
 \label{nsmom:stress}
\end{equation}
Products of all original source waves with the added curl field are
retained in this formula, including old and new curl remainders.

For two divergence-free fields $a,b$, Cartesian component expansion
gives
\[
 \nabla\cdot(a\otimes b+b\otimes a)
 =(b\cdot\nabla)a+(a\cdot\nabla)b,
 \qquad\nabla\cdot(v\otimes v)=(v\cdot\nabla)v.
\]
For example
$\sum_j\partial_j(a_i b_j)=\sum_jb_j\partial_ja_i+
a_i\sum_j\partial_jb_j$; the last term is zero. The lifted derivatives
obey the same product rule, so the identities hold on the extended
domain before evaluation. The three advection terms in
\eqref{nsmom:increment} are therefore exactly the divergence of
$u_*\otimes v+v\otimes u_*+v\otimes v$.

For completeness, set
$\mathscr L_0=\mathscr D_r^2+r^{-1}\mathscr D_r+
r^{-2}\partial_\theta^2+\partial_z^2$.
The vector Laplacian and pressure gradient in the two components are
\[
 (\Delta v)_\theta=(\mathscr L_0-r^{-2})v_\theta
                       +2r^{-2}\partial_\theta v_r,
 \qquad(\Delta v)_z=\mathscr L_0v_z,
 \qquad(\nabla\pi)_\theta=r^{-1}\partial_\theta\pi,
 \qquad(\nabla\pi)_z=\partial_z\pi.
\]
Equations \eqref{nsmom:mean_zero}--\eqref{nsmom:commute} prove separately
\begin{equation}
 \langle\mathscr D_t v\rangle_{\rm tan}=0,
 \qquad -\nu_{\rm NS}\langle\Delta v\rangle_{\rm tan}=0,
 \qquad\langle\nabla\pi\rangle_{\rm tan}=0.
 \label{nsmom:linear_zero}
\end{equation}
Thus viscosity and pressure were included before averaging; their
vanishing here is a proved consequence of the actual carriers.

If $T$ is symmetric, differentiation of
$e_r=(\cos\theta,\sin\theta,0)$,
$e_\theta=(-\sin\theta,\cos\theta,0)$ gives
\begin{align*}
 (\nabla\cdot T)_\theta
   &=(\mathscr D_r+2/r)T_{r\theta}
       +r^{-1}\partial_\theta T_{\theta\theta}+\partial_zT_{z\theta},\\
 (\nabla\cdot T)_z
   &=(\mathscr D_r+1/r)T_{rz}
       +r^{-1}\partial_\theta T_{\theta z}+\partial_zT_{zz}.
\end{align*}
Indeed $\partial_\theta e_r=e_\theta$ and
$\partial_\theta e_\theta=-e_r$ supply the two radial connection
terms in the first line and the single radial connection term in the
second. Averaging these formulas and using
\eqref{nsmom:commute}, \eqref{nsmom:stress}, and \eqref{nsmom:linear_zero}
proves the exact residual identities
\begin{equation}
 F_2:=\langle\Delta R_\theta\rangle
    =(\partial_r+2/r)S_{r\theta}+\partial_zS_{z\theta},
 \qquad
 F_1:=\langle\Delta R_z\rangle
    =(\partial_r+1/r)S_{rz}+\partial_zS_{zz}.
 \label{nsmom:F}
\end{equation}
In particular every mean-flow cross term vanishes as the divergence
of its zero averaged tensor, even when the mean depends on $Y$.

\subsection{The surviving axial fluxes and their signs}
\label{nsmom:moments}
Define the two moments with the original physical radial measure:
\begin{equation}
 M_2(z,t)=\int_0^\infty r^2F_2(r,z,t)\,dr,
 \qquad M_1(z,t)=\int_0^\infty rF_1(r,z,t)\,dr.
 \label{nsmom:Mdef}
\end{equation}
Their exact values are
\begin{align}
 M_2&=\partial_z J_2, &
 J_2&=\int_0^\infty r^2
    \langle w_{*,z}v_\theta+v_z w_{*,\theta}+v_zv_\theta\rangle\,dr,
 \label{nsmom:M2}\\
 M_1&=\partial_z J_1, &
 J_1&=\int_0^\infty r
    \langle2w_{*,z}v_z+v_z^2\rangle\,dr.
 \label{nsmom:M1}
\end{align}
To prove them, multiply \eqref{nsmom:F} by the respective weight:
$r^2(\partial_r+2/r)S_{r\theta}=\partial_r(r^2S_{r\theta})$
and $r(\partial_r+1/r)S_{rz}=\partial_r(rS_{rz})$.
The weighted boundary values at both ends of $(0,\infty)$ are zero,
because each entry of $S$ contains at least one compactly supported
factor $v$. Differentiation under the radial integral is valid by the
same compact support and smoothness. This proves
\eqref{nsmom:M2}--\eqref{nsmom:M1} including their signs and factors.

There is no pointwise cancellation of the axial derivatives in these
formulas. Since the added field has compact axial support, $J_1$ and
$J_2$ have compact axial support and
\begin{equation}
 \int_{\mathbb R}M_e(z,t)\,dz=0\quad(e=1,2).
 \label{nsmom:global_zero}
\end{equation}
This follows by evaluating the boundary values of $J_e$ at the two
axial ends. It does not assert $M_e(z,t)=0$ at an individual axial
point. The compact radial primitive in the covariance bridge must
therefore retain its $b_eM_e$ terms with these actual $M_e$.

At fixed source state, replacing the added field and pressure by their
negatives sends the old-wave cross terms in \eqref{nsmom:M2}--\eqref{nsmom:M1}
to their negatives and leaves the quadratic terms unchanged. Thus
negating an amplitude is not a sign reversal of either full moment.
Under the full physical reflection
$O=\operatorname{diag}(1,-1,1)$, applied to the source state and added
field together as $a^\sharp(x,t)=Oa(Ox,t)$, the cylindrical radial and
axial components retain their signs and the angular components reverse
sign at $(r,-\theta,z,t)$. Changing the angular integration variable
therefore proves
\[
 J_2^\sharp=-J_2,\quad M_2^\sharp=-M_2,
 \qquad J_1^\sharp=J_1,\quad M_1^\sharp=M_1.
\]

An explicit transverse-curl example shows why compactness and nonzero
angular frequency alone do not force either moment to vanish. This
example tests precisely those structural properties, and is not a
replacement for the retained source state. Let $H$ be any nonzero real
$C_c^\infty((r_-,r_+))$ function with $0<r_-<r_+$, let
$g\in C_c^\infty(\mathbb R)$ be nonconstant, let $n$ be a nonzero
integer, and put $f=rH'$. Use zero old state and the real potential
\[
 \mathcal A_r=fg\sin(n\theta),\qquad
 \mathcal A_\theta=0,\qquad
 \mathcal A_z=Hg\cos(n\theta).
\]
It has the required transverse representation:
$\phi=n\theta$, $\xi=(n/r)e_\theta$,
$C=(-ifg,0,Hg)$ in cylindrical components,
$a=i\xi\times C$, and
$C=i\xi\times a/|\xi|^2$ because $\xi\cdot C=0$.
The curl is exactly
\[
 v_r=-\frac{nHg}{r}\sin(n\theta),\qquad
 v_\theta=fg'\sin(n\theta)-H'g\cos(n\theta),\qquad
 v_z=-\frac{nfg}{r}\cos(n\theta).
\]
Consequently
\[
 M_2=ngg'\int_0^\infty r^2(H')^2\,dr,
 \qquad
 M_1=n^2gg'\int_0^\infty r(H')^2\,dr.
\]
For example $g(z)=\exp(-1/(1-z^2))$ on $|z|<1$, extended by zero,
has $g(z)g'(z)\ne0$ at every $0<|z|<1$. Both integrals are strictly
positive: a compactly supported smooth function with $H'=0$ everywhere
would be constant and hence zero. Thus both moments can be nonzero,
with their displayed signed values, within the exact transverse-curl
architecture.

\subsection{Bounds in the actual fields and original chart units}
\label{nsmom:bounds}
At fixed $(z,t)$ set
\[
 \|a\|_e^2=\int_0^\infty r^e\langle|a|^2\rangle\,dr,
 \qquad e=1,2.
\]
When an old-field norm is used below it can be restricted to the
enclosing radial support of $v$ and $\partial_zv$. This avoids imposing
any decay condition on an unused exterior part of the old field.
The following component bounds contain only actual fields:
\begin{align}
 |J_2|\le{}&\|w_{*,z}\|_2\|v_\theta\|_2
 +\|w_{*,\theta}\|_2\|v_z\|_2+\|v_z\|_2\|v_\theta\|_2,
 \label{nsmom:J2bound}\\
 |J_1|\le{}&2\|w_{*,z}\|_1\|v_z\|_1+\|v_z\|_1^2,
 \label{nsmom:J1bound}\\
 |M_2|\le{}&\|\partial_z w_{*,z}\|_2\|v_\theta\|_2
 +\|w_{*,z}\|_2\|\partial_zv_\theta\|_2
 +\|\partial_zv_z\|_2\|w_{*,\theta}\|_2\nonumber\\
 &+\|v_z\|_2\|\partial_z w_{*,\theta}\|_2
 +\|\partial_zv_z\|_2\|v_\theta\|_2
 +\|v_z\|_2\|\partial_zv_\theta\|_2,
 \label{nsmom:M2bound}\\
 |M_1|\le{}&2\bigl(\|\partial_z w_{*,z}\|_1\|v_z\|_1
 +\|w_{*,z}\|_1\|\partial_zv_z\|_1
 +\|v_z\|_1\|\partial_zv_z\|_1\bigr).
 \label{nsmom:M1bound}
\end{align}
Each is the Cauchy--Schwarz inequality on the product measure
$r^e\,dr\,d\theta/(2\pi)\,dY$, after applying the product rule in
\eqref{nsmom:M2}--\eqref{nsmom:M1}. For example, the derivative of
$2w_{*,z}v_z+v_z^2$ is
$2(\partial_z w_{*,z})v_z+2w_{*,z}\partial_zv_z+
2v_z\partial_zv_z$, proving every coefficient in
\eqref{nsmom:M1bound}.
For a $(z,t)$ multi-index $I$, every further derivative is exactly
obtained by replacing each product $ab$ in $J_e$ by
\[
 \partial^{I+(1,0)}(ab)
 =\sum_{K\le I+(1,0)}\binom{I+(1,0)}{K}
       (\partial^Ka)(\partial^{I+(1,0)-K}b)
\]
inside its original radial integral. This identity follows inductively
from the one-derivative product rule, and gives the corresponding
finite-order moment bound by applying the same Cauchy--Schwarz
inequality separately to every displayed summand.

The field $v$ in these bounds is not an abstract replacement amplitude.
For each actual band and harmonic write
\begin{gather}
 t_j=\chi_jB_jL_jy_j,\qquad
 c_j=\frac{i\,n_{\Phi,j}\times t_j}
                  {k_{\beta_j}m_j|n_{\Phi,j}|^2},\qquad
 b_j=t_j+\operatorname{curl}_{*,\mathrm{coeff}}c_j,\nonumber\\
 v=\sum_j\operatorname{Re}
       \bigl(Q_{\beta_j}^{-A}b_j e^{ik_{\beta_j}m_j\Phi_j}\bigr),
 \qquad A=\tfrac12+h,\quad D=\tfrac12-h.
 \label{nsmom:actual_amplitude}
\end{gather}
Here $\operatorname{curl}_{*,\mathrm{coeff}}$ differentiates the chart
coefficient, including the cutoff, auxiliary chain rule and cylindrical
basis, but not the displayed exponential. This is exactly the
remainder of the physical potential
$Q_{\beta_j}^{1/2-A}c_j e^{ik_{\beta_j}m_j\Phi_j}$: physical curl
introduces the factor $Q_{\beta_j}^{-1/2}$, and differentiating the
exponential produces $t_j$ by the transverse triple-product identity.
Explicitly, on the absolute auxiliary torus in a fixed band,
\[
 \mathscr D_R=\partial_R+
 \sqrt Q\,d_r r^{d_r-1}\mathbf v_r\cdot\partial_Y,
 \qquad r=\sqrt Q R,\qquad \mathscr D_Z^*=\varepsilon\partial_Z.
\]
For the actual coefficient $c=(c_r,c_\theta,c_z)$, which has
$\partial_\theta c=0$, the operator is exactly
\[
 \operatorname{curl}_{*,\mathrm{coeff}}c
 =\bigl(-\varepsilon\partial_Zc_\theta,
       \varepsilon\partial_Zc_r-\mathscr D_Rc_z,
       (\mathscr D_R+R^{-1})c_\theta\bigr).
\]
Coefficients originally on a band auxiliary torus are pulled back to
the absolute torus before applying this formula; this retains the
integer-covering chain derivative in $\partial_Yc$.
In the same fixed chart the exact axial derivative is
\begin{align}
 \partial_zv
 &=\sum_j\operatorname{Re}\left[
 Q_{\beta_j}^{-A-D}
 \bigl(\partial_{Z_j}b_j+
   ik_{\beta_j}m_j(\partial_{Z_j}\Phi_j)b_j\bigr)
 e^{ik_{\beta_j}m_j\Phi_j}\right].
 \label{nsmom:actual_derivative}
\end{align}
The band power is fixed during chart differentiation and
$\partial_z=Q_{\beta_j}^{-D}\partial_{Z_j}$; auxiliary phase evaluation
adds no axial derivative because of \eqref{nsmom:phase}. Thus
\eqref{nsmom:actual_amplitude}--\eqref{nsmom:actual_derivative} substitute
the actual retained pair and all its curl terms into the bounds.
For example, writing
$\|b\|_{e,\beta}^2=\int R^e\langle|b|^2\rangle\,dR$ in that band,
the triangle inequality and $|\operatorname{Re}z|\le|z|$ give
\begin{align}
 \|v\|_e
 &\le\sum_j Q_{\beta_j}^{(e+1)/4-A}\|b_j\|_{e,\beta_j},
 \label{nsmom:scale_bound}\\
 \|\partial_zv\|_e
 &\le\sum_j Q_{\beta_j}^{(e+1)/4-A-D}
 \|\partial_{Z_j}b_j+
 ik_{\beta_j}m_j(\partial_{Z_j}\Phi_j)b_j\|_{e,\beta_j}.
 \label{nsmom:scale_derivative_bound}
\end{align}
The exponent $(e+1)/4$ is exactly the square root of the radial measure
factor $r^e\,dr=Q^{(e+1)/2}R^e\,dR$. Distinct bands and harmonics
remain distinct in these sums. No limiting estimate as $q\downarrow0$
has been inferred from the finite compact-patch norms.

For a single fixed chart, put
$F_e=Q^{-2A-1/2}E_e^*$ as in the covariance bridge and
$S^*=Q^{2A}S$. Direct change of radial variable gives
\[
 M_e=Q^{e/2-2A}M_e^*,\qquad
 M_e^*=\int R^eE_e^*\,dR,
 \qquad
 M_e^*=\varepsilon\partial_Z\int R^e S^*_{z,\tau_e}\,dR,
 \quad\varepsilon=Q^h,
\]
where $\tau_2=\theta$ and $\tau_1=z$. Indeed the axial flux has
factor $Q^{(e+1)/2-2A}$, its axial derivative contributes $Q^{-D}$,
and dividing by $Q^{e/2-2A}$ leaves $Q^{1/2-D}=Q^h$.
This proves the precise relation to the source chart moment without
altering its axial anisotropy.

\subsection{Exact relation to moments after phase evaluation}
\label{nsmom:evaluated}
Let
\[
 \mathcal S=\langle w_*\otimes v+v\otimes w_*+v\otimes v\rangle_\theta,
 \qquad \mathcal S^\circ=\mathcal S-\langle\mathcal S\rangle_Y.
\]
Thus $\langle\mathcal S\rangle_Y=S$ and
$\langle\mathcal S^\circ\rangle_Y=0$. For the evaluated physical
fields, denote the angular-only residual moments by $M_e^{\rm ev}$.
The proof of \eqref{nsmom:mean_zero} holds after evaluation. The
physical divergence, vector Laplacian and gradient identities therefore
prove, without any auxiliary averaging,
\[
 M_e^{\rm ev}
   =\partial_z\int_0^\infty r^e
        \mathcal S_{z,\tau_e}(r,z,t,Y(r,t))\,dr.
\]
The physical radial boundary terms still vanish by compact support;
the chain rule is included in the physical radial derivative before
that integration. Subtracting \eqref{nsmom:M2}--\eqref{nsmom:M1} proves
the exact comparison
\begin{equation}
 M_e^{\rm ev}-M_e
 =\partial_z\int_0^\infty r^e
     \mathcal S^\circ_{z,\tau_e}(r,z,t,Y(r,t))\,dr.
 \label{nsmom:evaluation_difference}
\end{equation}
For example its absolute value is bounded by
$\int r^e\sup_Y|\partial_z\mathcal S^\circ_{z,\tau_e}|\,dr$,
because $Y(r,t)$ is independent of $z$. Zero auxiliary Haar average
does not make the right side of \eqref{nsmom:evaluation_difference}
zero at the physical phase. The displayed map supplies precisely the
remaining auxiliary-oscillation flux; it is not silently identified
with the source fast mean.


\section{Time derivatives of the actual signed third growth}
\label{tgdc:section}

This calculation is on the actual interval already constructed in
\emph{The actual third growth target on the moving second hold} and
\emph{Exact third growth on the moving second holding state}. It gives
quantitative derivatives of that trajectory for the source-frame map.
All derivatives in this section are with respect to coupled physical
time $t_c$, and $\tau=t_c-t_b$ denotes elapsed coupled time. The notation
does not use the cylindrical spatial radius for elapsed time.

\subsection{Original data, evolving parent, and proved interval}

Retain the source parameters
\begin{equation}\label{tgdc:original}
\begin{gathered}
 \nu=\sqrt{A_0},\qquad
 \mu=\lambda_0=\lceil\sqrt{A_0}/2\rceil,\qquad
 K_j=\mu\widehat\lambda_j,\\
 S_3=\frac{A_0}{\mu}\widehat\lambda_3^{-k_3-6},\qquad
 P_3^*=\frac{A_0}{\mu}\widehat\lambda_3^{-7/8}=S_3e^{L_3}.
\end{gathered}
\end{equation}
The physical common diffusivity is $d>0$; it is distinct from $\nu$.
Use the actual second return $t_b$, its values $h_b,A_{1,b},P_{2,b}$,
and the original angles $s_*,s_2,s_3$. Put
\begin{equation}\label{tgdc:constants}
\begin{gathered}
 a=\sin s_2,\quad \chi_b=\sqrt{h_b^2+a^2},\quad
 \beta_b=\arctan(a/h_b),\quad
 b=\sin(s_3+\beta_b),\quad g_b=\cos(s_3+\beta_b),\\
 c=bh_b-ag_b=\chi_b\sin s_3,\qquad
 S_*=A_0\sin s_*,\quad C_*=A_0\cos s_*.
\end{gathered}
\end{equation}
Here $a,b,c,S_*,C_*>0$, and the letter $b$ is the original oriented
geometric constant, whereas $b_3$ below is an amplitude coefficient.
The moving parent is exactly
\begin{equation}\label{tgdc:parent}
\begin{aligned}
 \chi^2&=h^2+a^2,& \dot h&=a\Omega_1,&
 \dot\Omega_1&=aA_1/\chi-dK_1^2\Omega_1,\\
 \dot A_1&=-S_*\Omega_1-dK_1^2A_1,&
 \dot P_2&=-dK_2^2\chi^2P_2,&
 \dot g&=b\Omega_1,\\
 g&=g_b+(b/a)(h-h_b),&
 N_3&=b(A_1+C_*)+gS_*+cK_2P_2,&
 R_3&=g^2+b^2.
\end{aligned}
\end{equation}
The retained feedback is $\dot\alpha=-a^2\Omega_1/\chi^2$ and
$\Omega_2=0$. In particular the older temperature and first vorticity
are not frozen. The cancellation in the complete numerator is
\begin{equation}\label{tgdc:numerator}
 \dot N_3=-dE_3,\qquad
 E_3=bK_1^2A_1+cK_2^3\chi^2P_2.
\end{equation}

Write $N_i=N_3(t_b)$, $a_i=N_i/K_3$,
$b_i=K_3\sin s_3$, $\gamma_i=\sqrt{a_i b_i}$ and
$u_i=\sqrt{b_i/a_i}$. The exact new equations, including their
original signed seed, are
\begin{equation}\label{tgdc:actual_pair}
\begin{gathered}
 y=\binom{\Theta_3}{\Omega_3},\qquad
 \dot y=C_3y,\qquad
 C_3=\begin{pmatrix}-\lambda_3&a_3\\b_3&-\lambda_3\end{pmatrix},\\
 a_3=\frac{N_3}{K_3R_3},\qquad
 b_3=\frac{K_3c}{\chi},\qquad
 \lambda_3=dK_3^2R_3,\qquad
 y(t_b)=-S_3\binom1{u_i}.
\end{gathered}
\end{equation}
There are no new amplitude controls in this growth pair:
$c_\vartheta=c_\omega=0$. The activation of the compact wave is a
separate original factor $\eta(\Gamma_3\tau)$; it does not replace
\eqref{tgdc:actual_pair} by a forced coefficient trajectory.

For clarity, the following facts about this actual interval were proved
in the third-threshold construction. They identify the trajectory to
which the estimates below apply:
\begin{equation}\label{tgdc:proved_boxes}
\begin{gathered}
 I_3=[t_b,t_3^a],\quad t_3^a-t_b=\tau_a<14/\sigma_1,\quad
 P_3=-\Theta_3>0,\quad P_3(t_b)=S_3,\quad P_3(t_3^a)=P_3^*,\\
 u=\Omega_3/\Theta_3,\qquad
 \frac{u_i}{\sqrt2}\leq u\leq2u_i,\qquad
 \frac{a_i}4\leq a_3\leq a_i,\qquad
 \frac{b_i}2\leq b_3\leq b_i,\\
 1\leq R_3\leq\overline R:=\left(\frac{265}{256}\right)^2<2,
 \quad 1\leq\frac\chi{\chi_b}\leq\frac{259}{256},\quad
 0<dK_3^2\leq\frac{\gamma_i}{16\sqrt2},\\
 \frac{\gamma_i}{8\sqrt2}\leq (\log P_3)'\leq2\gamma_i,
 \quad 0<A_1\leq A_{1,b},\quad 0<P_2\leq P_{2,b},\quad
 0<\Omega_1<9\sigma_1,\\
 h_b\leq h\leq h_b+144a,\qquad
 g_b\leq g\leq g_b+144b,\qquad
 \tfrac12N_i\leq N_3\leq N_i,\qquad 1\leq\chi^2<10.
\end{gathered}
\end{equation}
The value of $d$ is precisely the previously selected range
$0<d\leq d_{3,\max}$, retaining its $K_3^{-2}$ dependence.
Thus \eqref{tgdc:proved_boxes} does not assert a new unrestricted
viscosity range or assume a further return.

\subsection{Signed amplitudes and first derivatives with the diffusion retained}

\begin{theorem}\label{tgdc:first_derivatives}
For every $t_c\in I_3$, both $\Theta_3(t_c)$ and $\Omega_3(t_c)$
are negative and strictly decreasing. Their exact component rates are
\begin{equation}\label{tgdc:exact_first}
 \dot\Theta_3=-P_3(a_3u-\lambda_3),\qquad
 \dot\Omega_3=-P_3(b_3-\lambda_3u).
\end{equation}
They obey the bounds
\begin{equation}\label{tgdc:component_bounds}
\begin{gathered}
 S_3e^{\gamma_i\tau/(8\sqrt2)}
 \leq P_3(t_c)\leq\min\{P_3^*,S_3e^{2\gamma_i\tau}\},\\
 \frac{u_i}{\sqrt2}P_3\leq |\Omega_3|\leq2u_iP_3,\\
 \left(\frac{\gamma_i}{4\sqrt2}-dK_3^2\overline R\right)P_3
 \leq |\dot\Theta_3|\leq2\gamma_iP_3,\\
 u_i\left(\frac{\gamma_i}{2}-2dK_3^2\overline R\right)P_3
 \leq |\dot\Omega_3|\leq\gamma_i u_iP_3.
\end{gathered}
\end{equation}
In particular both lower derivative bounds are positive, and
\begin{equation}\label{tgdc:norm_bounds}
\begin{gathered}
 P_3\sqrt{1+u_i^2/2}\leq |y|_2
       \leq P_3\sqrt{1+4u_i^2}\leq P_3^*\sqrt{1+4u_i^2},\\
 |\dot y|_2\leq\gamma_iP_3\sqrt{4+u_i^2}
       \leq\gamma_iP_3^*\sqrt{4+u_i^2}.
\end{gathered}
\end{equation}
The original beginning and attained endpoint have the exact derivatives
\begin{equation}\label{tgdc:endpoint_derivatives}
\begin{aligned}
 \dot y(t_b)&=-S_3(\gamma_i-dK_3^2)\binom1{u_i},\\
 y(t_3^a)&=-P_3^*\binom1{u(t_3^a)},\\
 \dot y(t_3^a)&=-P_3^*
 \binom{a_3(t_3^a)u(t_3^a)-dK_3^2R_3(t_3^a)}
       {b_3(t_3^a)-dK_3^2R_3(t_3^a)u(t_3^a)}.
\end{aligned}
\end{equation}
\end{theorem}
\begin{proof}
The signs and quotient in \eqref{tgdc:proved_boxes} give
$y=-P_3(1,u)^T$. Substituting this expression in the exact pair proves
\eqref{tgdc:exact_first}. Integrating the already proved logarithmic
rate from $t_b$ and using the attained endpoint and strict increase
of $P_3$ gives the first line of \eqref{tgdc:component_bounds}.
The quotient box gives its second line. Since $a_i u_i=\gamma_i$
and $b_i=\gamma_i u_i$, direct multiplication of the positive boxes
gives
\[
 \frac{\gamma_i}{4\sqrt2}\leq a_3u\leq2\gamma_i,
 \qquad
 \frac{\gamma_i u_i}{2}\leq b_3\leq\gamma_i u_i,
 \qquad 0<\lambda_3\leq dK_3^2\overline R.
\]
Consequently the displayed lower bound for $a_3u-\lambda_3$ follows.
It is at least $\gamma_i/(8\sqrt2)>0$, because
$\overline R<2$ and $dK_3^2\leq\gamma_i/(16\sqrt2)$.
Similarly,
\[
 b_3-\lambda_3u
 \geq u_i(\gamma_i/2-2dK_3^2\overline R)
 >\gamma_i u_i(1/2-1/(4\sqrt2))
 >\gamma_i u_i/4>0.
\]
Thus both derivatives in \eqref{tgdc:exact_first} are strictly negative.
Dropping their nonnegative diffusion losses only for the purpose of
an upper bound gives the two upper derivative estimates; the exact
rates continue to contain the losses. Squaring the two amplitude and
two derivative bounds, adding them, and taking the positive square
root proves \eqref{tgdc:norm_bounds}. Finally $R_3(t_b)=1$,
$u(t_b)=u_i$, $a_3(t_b)=a_i$ and $b_3(t_b)=b_i$ give the first
endpoint identity. Substitution of the proved target value gives
the other two. In particular no later endpoint was substituted for
the actual $u(t_3^a)$.
\end{proof}

\subsection{Explicit coefficient rates}

Set $\overline H=h_b+144a$, $\overline G=g_b+144b$ and
$\overline W=9\sigma_1$. Differentiating the complete coefficients
gives
\begin{equation}\label{tgdc:coefficient_rates}
\begin{aligned}
 \dot a_3&=-a_3\left(\frac{dE_3}{N_3}
                            +\frac{2bg\Omega_1}{R_3}\right),\\
 \dot b_3&=-b_3\frac{ah\Omega_1}{\chi^2},\qquad
 \dot\lambda_3=2dK_3^2bg\Omega_1.
\end{aligned}
\end{equation}
Every quantity on the right comes from the moving parent. In particular
the exact numerator derivative \eqref{tgdc:numerator} is used before
any upper estimate. Because $K_1<K_2$, $\chi^2<10$ and all summands
of $N_3$ are positive, $E_3\leq10K_2^2N_3$. Thus
\begin{equation}\label{tgdc:coefficient_rate_bounds}
\begin{gathered}
 |\dot a_3|\leq
 a_i(10dK_2^2+2b\overline G\overline W),\qquad
 |\dot b_3|\leq
 b_i a\overline H\overline W\chi_b^{-2},\\
 |\dot\lambda_3|\leq
 2dK_3^2 b\overline G\overline W.
\end{gathered}
\end{equation}
These formulas also display the distinct inherited $K_2^2$ loss and
new $K_3^2$ loss. They do not replace either by the source scaling
parameter $\nu$.

\subsection{A closed all-orders derivative calculation}

Here are finite recurrences whose values are explicit upper bounds
from the actual birth data and the already proved interval. They use
ordinary derivatives, with no change to physical time or amplitudes.
For integers $n\geq0$ and $m\geq2$ define the product expression
\[
 \mathcal P_n(X^{(1)},\ldots,X^{(m)})
 =\sum_{j_1+\cdots+j_m=n}
       \frac{n!}{j_1!\cdots j_m!}
       X^{(1)}_{j_1}\cdots X^{(m)}_{j_m}.
\]
This notation is a finite sum, and repeated arguments mean repeated
factors. Define the nonnegative arrays $H,W,A,P,Q,G,N,J$ initially by
\begin{equation}\label{tgdc:majorant_initial}
\begin{gathered}
 H_0=\overline H,\quad W_0=\overline W,\quad
 A_0^{\rm bd}=A_{1,b},\quad P_0=P_{2,b},\quad
 Q_0=\chi_b^{-1},\quad G_0=\overline G,\quad
 N_0=N_i,\quad J_0=1.
\end{gathered}
\end{equation}
In the following recurrences the array $A$ has entries $A_n^{\rm bd}$;
the superscript distinguishes its zeroth entry from the unchanged
physical constant $A_0$ in \eqref{tgdc:original}. Write
\[
 X_n=\mathcal P_n(H,H)+a^2\mathbf1_{n=0}.
\]
Having computed all eight arrays up to index $n$, compute index $n+1$
by the explicit formulas
\begin{equation}\label{tgdc:parent_majorants}
\begin{aligned}
 H_{n+1}&=aW_n,\\
 W_{n+1}&=a\mathcal P_n(A,Q)+dK_1^2W_n,\\
 A_{n+1}^{\rm bd}&=S_*W_n+dK_1^2A_n^{\rm bd},\\
 P_{n+1}&=dK_2^2\mathcal P_n(X,P),\\
 Q_{n+1}&=a\mathcal P_n(H,W,Q,Q,Q),\\
 G_{n+1}&=bW_n,\\
 N_{n+1}&=dbK_1^2A_n^{\rm bd}
                 +dcK_2^3\mathcal P_n(X,P),\\
 J_{n+1}&=2b\mathcal P_n(G,W,J,J).
\end{aligned}
\end{equation}
For the phase-length array use
\[
 R_0^{\rm bd}=\overline R,\qquad
 R_n^{\rm bd}=\mathcal P_n(G,G)\quad(n\geq1).
\]
The coefficient arrays are
\begin{equation}\label{tgdc:coefficient_majorants}
\begin{gathered}
 \mathfrak A_0=a_i,\quad \mathfrak B_0=b_i,\quad
 \mathfrak D_0=dK_3^2\overline R,\\
 \mathfrak A_n=K_3^{-1}\mathcal P_n(N,J),\qquad
 \mathfrak B_n=K_3cQ_n,\qquad
 \mathfrak D_n=dK_3^2R_n^{\rm bd}\quad(n\geq1).
\end{gathered}
\end{equation}
No optimization of these finite sums is needed for their validity.
In particular $\mathfrak A_1$ may be replaced by the smaller of its
value above and the first bound in \eqref{tgdc:coefficient_rate_bounds},
and similarly for the other two first coefficient derivatives.

The amplitude arrays start with the actual target and the sharper
first-derivative bounds already proved:
\begin{equation}\label{tgdc:amplitude_majorants}
\begin{gathered}
 T_0=P_3^*,\qquad O_0=2u_iP_3^*,\qquad
 T_1=2\gamma_iP_3^*,\qquad O_1=\gamma_i u_iP_3^*,\\
 T_{n+1}=\sum_{j=0}^n\binom nj
      (\mathfrak D_jT_{n-j}+\mathfrak A_jO_{n-j}),\qquad n\geq1,\\
 O_{n+1}=\sum_{j=0}^n\binom nj
      (\mathfrak B_jT_{n-j}+\mathfrak D_jO_{n-j}),\qquad n\geq1.
\end{gathered}
\end{equation}
The choice to begin these recurrences at $n=1$ preserves the stronger
first-derivative estimates instead of substituting a larger matrix-norm
bound.

\begin{theorem}\label{tgdc:all_orders}
For every integer $n\geq0$ the finite arrays above satisfy on $I_3$
\begin{equation}\label{tgdc:all_order_bounds}
\begin{gathered}
 |\partial_{t_c}^n\Theta_3|\leq T_n,\qquad
 |\partial_{t_c}^n\Omega_3|\leq O_n,\qquad
 |\partial_{t_c}^ny|_2\leq\sqrt{T_n^2+O_n^2},\\
 |\partial_{t_c}^na_3|\leq\mathfrak A_n,\qquad
 |\partial_{t_c}^nb_3|\leq\mathfrak B_n,\qquad
 |\partial_{t_c}^n\lambda_3|\leq\mathfrak D_n.
\end{gathered}
\end{equation}
The exact, signed amplitude derivative recurrence underlying these
bounds is
\begin{equation}\label{tgdc:exact_derivative_recurrence}
 \partial_{t_c}^{n+1}y
     =\sum_{j=0}^n\binom nj
            (\partial_{t_c}^jC_3)(\partial_{t_c}^{n-j}y),
 \qquad n\geq0.
\end{equation}
In particular the exact second derivative contains the changing
coefficient and the complete diffusion square:
\begin{equation}\label{tgdc:second_derivative}
 \ddot y=
 \begin{pmatrix}
 \lambda_3^2+a_3b_3-\dot\lambda_3&\dot a_3-2\lambda_3a_3\\
 \dot b_3-2\lambda_3b_3&\lambda_3^2+a_3b_3-\dot\lambda_3
 \end{pmatrix}y.
\end{equation}
\end{theorem}
\begin{proof}
Introduce the exact auxiliary functions $q=\chi^{-1}$ and
$j=R_3^{-1}$, both well-defined by \eqref{tgdc:proved_boxes}.
Direct differentiation of their specified inverses gives
\[
 \dot q=-ah\Omega_1q^3,\qquad
 \dot j=-2bg\Omega_1j^2.
\]
These identities retain the actual phase lengths; the symbols do not
replace them by constant values. The complete equations for
$(h,\Omega_1,A_1,P_2,q,g,N_3,j)$ are therefore the signed polynomial
system consisting of \eqref{tgdc:parent}, \eqref{tgdc:numerator} and
these two equations, with $\chi^2=h^2+a^2$ in the parent losses.
Their zero-order absolute values are bounded respectively by
$(H_0,W_0,A_0^{\rm bd},P_0,Q_0,G_0,N_0,J_0)$:
the claims for $q$ and $j$ use $\chi\geq\chi_b$ and $R_3\geq1$.

For any $m$ differentiable factors the ordinary $n$th derivative of
their product is $\mathcal P_n$ applied to their derivative arrays.
To prove this formula without assuming it, start with $n=0$.
Differentiating a term of order $n$ once chooses one factor whose
index increases. At indices summing to $n+1$, the resulting coefficient
is the sum of $n!/(j_1!\cdots(j_k-1)!\cdots j_m!)$ over indices
$k$ with $j_k>0$. It is
$n!(j_1+\cdots+j_m)/(j_1!\cdots j_m!)=(n+1)!/(j_1!\cdots j_m!)$.
This proves the formula by induction.

Apply this exact product derivative identity to the signed parent
equations and then the triangle inequality. For instance
\[
 \partial_{t_c}^{n+1}N_3
  =-dbK_1^2\partial_{t_c}^n A_1
   -dcK_2^3\partial_{t_c}^n\{(h^2+a^2)P_2\},
\]
and its absolute value is bounded by the displayed recurrence for
$N_{n+1}$. This example shows explicitly that neither inherited
diffusion term is omitted. The other seven equations give exactly
\eqref{tgdc:parent_majorants}. Induction in $n$ proves all eight parent
derivative bounds. The identities $R_3=g^2+b^2$,
$a_3=N_3j/K_3$, $b_3=K_3cq$ and $\lambda_3=dK_3^2R_3$ then prove
\eqref{tgdc:coefficient_majorants} for $n\geq1$ by the same product
formula. Their zero-order bounds were already proved in
\eqref{tgdc:proved_boxes}.

Repeatedly differentiate the exact pair to obtain
\eqref{tgdc:exact_derivative_recurrence}. Its componentwise absolute
value gives \eqref{tgdc:amplitude_majorants}, with the zero and first
orders supplied by Theorem~\ref{tgdc:first_derivatives}. Induction
now proves the two amplitude bounds at every order, and the Euclidean
norm bound follows by summing squares. Finally
$\ddot y=(\dot C_3+C_3^2)y$ and direct multiplication of the two by two
matrix give \eqref{tgdc:second_derivative}. Every recurrence uses
finitely many sums at each specified order. Thus each resulting bound
is a determinate expression in the original physical constants and
actual birth data, rather than an unevaluated supremum of a missing
coefficient derivative.
\end{proof}

\subsection{Derivative transport through the specified source clock}

Let $L_s>0$ be the source slot length used for this actual interval,
and define its specified clock by
\[
 \alpha_c=\tau_a/L_s,\qquad
 t_c=t_b+\alpha_c v,\qquad 0\leq v\leq L_s.
\]
The two actual endpoints are retained. Since $\alpha_c$ is constant
for the fixed trajectory and slot, repeated chain differentiation
gives, exactly,
\begin{equation}\label{tgdc:clock}
 \partial_v^ny(t_b+\alpha_cv)
      =\alpha_c^n\partial_{t_c}^ny(t_b+\alpha_cv),\qquad
 |\partial_v^ny|_2\leq\alpha_c^n\sqrt{T_n^2+O_n^2}.
\end{equation}
In a source chart with fixed label $Q$, exponent $h_{\rm src}$,
and $\partial_t v=Q^{-1-h_{\rm src}}$, the derivative contribution
acting on this amplitude alone is consequently
\[
 \partial_t^ny(t_b+\alpha_cv)
  =(\alpha_cQ^{-1-h_{\rm src}})^n
                      \partial_{t_c}^ny(t_b+\alpha_cv).
\]
This identity holds with the stated fixed chart parameters and affine
clock; derivatives of any additional frame, cutoff, or varying
parameter must be included by the ordinary product rule. In particular
neither the clock factor nor the physical frequency has been replaced
by one. These finite derivative estimates establish no uniform
terminal extension for an infinite sequence.


\section{The actual third-growth pair in the source clock and frame}
\label{tgframe:section}

This calculation uses equations (6.2)--(6.12), (7.2)--(7.12), and
Lemma 7.4 of the retained 166-page source. Its SHA-256 is displayed
in two consecutive segments:
\[
\begin{gathered}
\texttt{0e779481c4da40bd28d1e642e1d8ca57}\\
\texttt{447d129610df28dfa5a11e9af8ae228f}.
\end{gathered}
\]
The coupled input is the actual third-growth interval proved in
\path{finite_stage/third_growth_entry.tex}. The source's pulse
equation is written at physical viscosity one. We retain an arbitrary
positive physical viscosity $\nu_{\rm NS}$ when computing a candidate
residual; when identifying the source's particular homogeneous pulse,
we explicitly set $\nu_{\rm NS}=1$ and harmonic $m=1$.
The original coupled time scale $\nu=\sqrt{A_0}$, the physical coupled
diffusivity $d$, and the source chart parameter $\varepsilon=Q^h$
are distinct.
The source growth constant denoted $\lambda_0$ in $\lambda(v)$ below
is the source's own growth constant. It is not identified with the
coupled integer frequency $\lambda_0=\mu=\lceil\sqrt{A_0}/2\rceil$.

\subsection{Exact physical differentiation of the source pulse clock}

Fix a source label and one of its local lifts. Set
\[
 b_g=\sqrt2-1,\quad v_r=(1,-b_g),\quad v_t=(b_g,1),\quad
 J_g=\begin{pmatrix}3&1\\1&5\end{pmatrix},\quad
 T_g=4+\sqrt2,\quad\Lambda_g=4-\sqrt2.
\]
Direct multiplication gives $J_gv_t=T_gv_t$ and
$J_gv_r=\Lambda_gv_r$, while $v_t\cdot v_r=0$ and
$|v_t|^2=|v_r|^2=1+b_g^2$. The source retains
\[
\begin{gathered}
 d_r=2\left((1+h)\frac{\log\Lambda_g}{\log T_g}-h\kappa_s\right),\qquad
 \kappa_s=10^{-5},\\
 Q=2^{-\ell},\quad
 i=\left\lfloor\log_{T_g}\frac{Q^{-1-h}}{S_*}\right\rfloor,\quad
 c_i=T_g^iQ^{1+h}.
\end{gathered}
\]
On the local lift
\[
 Y_i-c_\gamma-k_{\rm copy}=\xi_gv_r+\eta_gv_t,\qquad
 v=\frac{\eta_g+r_0}{c_i},\qquad L_s=\frac{2r_0}{c_i},
 \quad
 \eta_g=\frac{v_t\cdot(Y_i-c_\gamma-k_{\rm copy})}{1+b_g^2}.
\]
The label and integer $k_{\rm copy}$ are held fixed on this lift.
Consequently $L_i\eta_g=0$, $N_i\eta_g=1$, where
$L_i=v_r\cdot\partial_{Y_i}$ and $N_i=v_t\cdot\partial_{Y_i}$.
The clock has no explicit dependence on the slow variables $R,Z,T$.
The evaluated chart derivatives are exactly
\[
 D_r=\partial_R+M_id_rR^{d_r-1}L_i,\qquad
 D_z=\varepsilon\partial_Z,\qquad
 t_*=-\varepsilon\partial_T+c_iN_i,
 \quad M_i=\Lambda_g^iQ^{d_r/2}.
\]
It follows by applying these displayed operators, without estimating
any coefficient, that
\begin{equation}\label{tgframe:clockchart}
 D_rv=D_zv=0,\qquad t_*v=1.
\end{equation}
There is also an explicit physical verification. Evaluate at the
source map $Y=v_rr^{d_r}+v_tt$ modulo the lattice. Locally the lift of
$Y_i=J_g^iY$ differs from
$\Lambda_g^iv_rr^{d_r}+T_g^iv_tt$ by a constant lattice vector.
The dual formula for $\eta_g$ therefore gives
$\eta_g=T_g^it+C_\gamma$ with a locally constant $C_\gamma$.
Thus, on the evaluated local rectangle,
\begin{equation}\label{tgframe:clockphysical}
 \partial_rv=\partial_zv=\partial_\theta v=0,\qquad
 \partial_tv=\frac{T_g^i}{c_i}=Q^{-1-h}.
\end{equation}
The $\partial_t$ in this equation is the derivative of the evaluated
physical function. On the extended domain the corresponding derivative
is $\partial_t+N_{\rm abs}$. These statements are identical after
evaluation; neither drops the auxiliary derivative.

Let $t_b$ and $r_a>0$ be the actual second return time and the actual
third target-crossing elapsed time. Define
\begin{equation}\label{tgframe:coupledclock}
 t_c(v)=t_b+\alpha v,\qquad \alpha=\frac{r_a}{L_s}>0.
\end{equation}
This parametrizes the entire proved coupled growth interval, with its
original endpoints. In physical source variables its derivative is
$\partial_tt_c=\alpha Q^{-1-h}$ and all three spatial derivatives vanish.
The positive orientation of this clock asserts no sign for a velocity.
The coupled amplitudes below are strictly negative.

\subsection{The actual frame, its inverse, and the direct candidate}

Write the original source phase normal as
\[
 n=\left(x_0-v(pF_R+p_zG_R),\ \frac pR,\
             p_z-\varepsilon v(pF_Z+p_zG_Z)\right).
\]
All $p,p_z,x_0,k,Q$ are the original fixed label constants, including
$kp\in\mathbb Z\setminus\{0\}$. Set
\[
\begin{gathered}
 K_a=\frac{n_{\tan}}{|n_{\tan}|},\quad
 N_a=((K_a)_z,-(K_a)_\theta),\quad
 s_a=\frac{n_r}{|n_{\tan}|},\\
 c(v)=c_0\sqrt{1+s(v)^2},
 \quad
 s(v)=\sigma\left(\frac{u_*}{2}+\frac{u_*v}{L_s}\right).
\end{gathered}
\]
Here $c_0<0$ is the unchanged source constant, and $K_a,N_a$
are identified with their tangential three-vectors. In particular
$c(v)\ne0$. Because $p/R\ne0$, every formula is defined throughout
the source rectangle. The actual source frame is
\[
 U=[e_r-s_aK_a,\ N_a],\qquad
 M=\begin{pmatrix}1&1\\c&-c\end{pmatrix},\qquad B=UM.
\]
The two columns of $U$ are orthogonal, with squared lengths
$g_a^2=1+s_a^2$ and $1$. Both are orthogonal to $n$.
Since $\det M=-2c\ne0$, $B$ is a bijection from $\mathbb R^2$
onto $n^\perp$. The source left inverse, defined on all of
$\mathbb R^3$, is
\begin{equation}\label{tgframe:leftinverse}
 B_\ell=M^{-1}\begin{pmatrix}e_r^T\\N_a^T\end{pmatrix},
 \qquad
 M^{-1}=\frac12\begin{pmatrix}1&c^{-1}\\1&-c^{-1}\end{pmatrix}.
\end{equation}
Indeed the last two rows applied to $U$ give $I_2$.

The eigenvalues of
\[
 B^TB=
 \begin{pmatrix}g_a^2+c^2&g_a^2-c^2\\
                 g_a^2-c^2&g_a^2+c^2\end{pmatrix}
\]
are $2g_a^2$ and $2c^2$, with orthogonal eigenvectors $(1,1)$
and $(1,-1)$. The Euclidean operator norm and the inverse norm on
the tangent plane are therefore exactly
\begin{equation}\label{tgframe:framenorms}
 \|B\|=\sqrt2\max(g_a,|c|),\qquad
 \|(B|_{\mathbb R^2})^{-1}\|_{n^\perp\to\mathbb R^2}
       =\frac1{\sqrt2\min(g_a,|c|)}.
\end{equation}
The ambient norm of the particular extension \eqref{tgframe:leftinverse}
is $(1/\sqrt2)\max(1,|c|^{-1})$; it need not equal the inverse norm
restricted to $n^\perp$.

Retain the actual coupled data
\[
 K_3=\mu\widehat\lambda_3,\quad
 S_3=\frac{A_0}{\mu}\widehat\lambda_3^{-k_3-6},\quad
 P_3^*=\frac{A_0}{\mu}\widehat\lambda_3^{-7/8}=S_3e^{L_3},
 \quad u_i=\sqrt{\frac{b_i}{a_i}}
           =\frac{K_3}{\sigma_2\sqrt{n_i}}.
\]
On the interval \eqref{tgframe:coupledclock}, the already constructed
solution is
\[
 y_3(t_c)=\binom{\Theta_3}{\Omega_3}
        =-P(v)\binom{1}{u(v)},\qquad
 S_3\le P(v)\le P_3^*,\qquad
 \frac{u_i}{\sqrt2}\le u(v)\le2u_i,
\]
with $P(0)=S_3$, $P(L_s)=P_3^*$ and $u(0)=u_i$.
These inequalities follow from the actual evolving-parent system,
rather than from a new frame approximation.
The direct candidate $t_{\rm dir}=By_3$ consequently has the exact
components
\begin{align}
 (t_{\rm dir})_r&=-P(1+u),\nonumber\\
 (t_{\rm dir})_{\tan}
       &=P s_a(1+u)K_a-Pc(1-u)N_a,\label{tgframe:directcomponents}\\
 |t_{\rm dir}|^2
       &=P^2\bigl[(1+s_a^2)(1+u)^2+c^2(1-u)^2\bigr],\nonumber\\
 (t_{\rm dir})_r(t_{\rm dir})_{\tan}
       &=-P^2s_a(1+u)^2K_a+P^2c(1-u^2)N_a.\nonumber
\end{align}
These follow by multiplying $My_3=-P(1+u,c(1-u))^T$ and using
the orthogonality of the two columns of $U$. In particular
\[
 |(t_{\rm dir})_r|\ge S_3(1+u_i/\sqrt2),\qquad
 |t_{\rm dir}|\le\sqrt2\max(g_a,|c|)P_3^*\sqrt{1+4u_i^2}.
\]
At the actual target,
\begin{equation}\label{tgframe:endpointomega}
 \frac{P_3^*u_i}{\sqrt2}\le|\Omega_3(t_b+r_a)|
                      \le2P_3^*u_i,\qquad
 P_3^*u_i=\frac{A_0}{\sigma_2\sqrt{n_i}}
                          \widehat\lambda_3^{1/8}.
\end{equation}
The last equality retains $\mu$ until it cancels against its identical
factor in $K_3$. Thus a decreasing temperature target does not give a
decreasing two-component target.

\subsection{The complete principal mismatch, including the cutoff}

With a prime denoting $v$ differentiation at fixed slow variables,
retain
\[
 \mathsf K=\begin{pmatrix}0&-2F&0\\2F+RF_R&0&0\\G_R&0&0\end{pmatrix},
 \quad
 A_\Phi=-\mathsf K+
        \frac{n(n^T\mathsf K-(n')^T)}{|n|^2},
 \quad H=B_\ell(A_\Phi B-B').
\]
Differentiating $n^TB=0$ shows
$n^TB'=-(n')^TB$. Direct multiplication gives
$n^TA_\Phi B=-(n')^TB$ as well. Therefore
$A_\Phi B-B'$ takes values in $n^\perp$, and
\begin{equation}\label{tgframe:frametransport}
 A_\Phi B-B'=BH.
\end{equation}
The source defines the exact error matrix
$E=H-\operatorname{diag}(\lambda,-\lambda)$ with
$\lambda=\lambda_0/\sqrt{1+s^2}$. Its source estimate
$\|E\|\le C/S_*$ is not an assertion that $E$ is zero.

The actual coupled matrix and arbitrary-viscosity source loss are
\[
\begin{gathered}
 C_c=\begin{pmatrix}-d_c&a_3\\b_3&-d_c\end{pmatrix},
 \quad d_c=dK_3^2R_3,\quad
 a_3=\frac{N_3}{K_3R_3},\quad b_3=\frac{K_3c_{\rm det}}{\chi_{\rm old}},\\
 \quad
 \delta_\nu=\nu_{\rm NS}m^2\varepsilon k^2|n|^2,\quad
 \Delta_\nu=\delta_\nu-\alpha d_c.
\end{gathered}
\]
Here $c_{\rm det}$ and $\chi_{\rm old}$ are the third-growth determinant
constant and older phase length; these names distinguish them from the
source $c(v)$ and the cutoff below. Every coupled quantity is evaluated
at $t_c(v)$. The actual third growth has zero added amplitude controls,
so $y_3'=\alpha C_cy_3$ exactly.

Let $\chi=\eta(R,Z,T)\chi_g(\xi_g)\psi(v)$ be a smooth scalar
cutoff with the source support, and set $t_\chi=\chi By_3$.
The homogeneous-pressure formula, with zero principal source, is
\[
 \pi_\chi=\frac{i}{km|n|^2}
        \bigl(n\cdot\mathsf Kt_\chi-n'\cdot t_\chi\bigr).
\]
The cutoff leaves $n^Tt_\chi=0$. Differentiating this identity
and substituting the pressure gives
\begin{equation}\label{tgframe:residual}
 t_\chi'+\mathsf Kt_\chi+\delta_\nu t_\chi+ikmn\pi_\chi
 =B\mathcal E_\chi,\qquad
 \mathcal E_\chi=\chi(\alpha C_c-H+\delta_\nu I)y_3+\chi'y_3.
\end{equation}
For completeness, the normal component of the sum before pressure is
$-n'\cdot t_\chi+n\cdot\mathsf Kt_\chi$; multiplication by $ikmn$
times the displayed pressure cancels it. The remaining vector is
$t_\chi'-A_\Phi t_\chi+\delta_\nu t_\chi$. Use
\eqref{tgframe:frametransport} and $y_3'=\alpha C_cy_3$ to obtain
\eqref{tgframe:residual}.

In original scalar entries this is
\begin{equation}\label{tgframe:residualentries}
 \mathcal E_\chi=-P
 \left\{\chi
 \binom{\Delta_\nu-\lambda-E_{11}+(\alpha a_3-E_{12})u}
       {\alpha b_3-E_{21}+(\Delta_\nu+\lambda-E_{22})u}
       +\chi'\binom1u\right\}.
\end{equation}
All entries are signed. In particular the damping mismatch is
$\delta_\nu-\alpha dK_3^2R_3$, not the sum of the losses.
At fixed slow variables and fixed $\xi_g$, $\chi'=\eta\chi_g\psi'$.
The physical derivative of $\chi$ also includes
$-\varepsilon(\partial_T\eta)\chi_g\psi$ in $t_*\chi$; that slow term
belongs to the complete physical residual and has not been included in,
or removed by, the principal equation \eqref{tgframe:residual}.
A direct finite bound is
\[
 |\mathcal E_\chi|
 \le P\sqrt{1+u^2}\left[
 |\chi|\bigl(|\Delta_\nu|+\lambda+\|E\|
                  +\alpha\max(a_3,b_3)\bigr)+|\chi'|\right].
\]
This is the triangle inequality and the exact norm
$\left\|\begin{smallmatrix}0&a_3\\b_3&0\end{smallmatrix}\right\|
=\max(a_3,b_3)$, with $a_3,b_3>0$.
It retains the carrier, frame, cutoff and all physical diffusion factors.

\subsection{An explicit map of the actual seed to the actual source pulse}

For this subsection only, take $m=1$, $\nu_{\rm NS}=1$, so
$\delta_1=\varepsilon k^2|n|^2$. This is precisely the source
homogeneous pulse equation with its given background. Put
\[
 d_{\rm ref}=\varepsilon k^2B_s^2(1+s^2),\qquad
 {\cal P}_\sigma(v)=
   \exp\left(\int_{L_s/2}^{v}(\lambda(w)-d_{\rm ref}(w))\,dw\right).
\]
Every coefficient is finite on the closed interval. Thus
$\mathcal P_\sigma(0)>0$ exactly, without a Gaussian approximation.
Source Lemma 7.4 specifies the initial coordinate vector
\[
 z_{\sigma,0}=\binom{{\cal P}_\sigma(0)}0,\qquad
 t^h_\sigma(0)=B(0)z_{\sigma,0}.
\]
The fixed source label may affect this strictly positive scalar.

For any nonzero real target vector $z_0=(z_1,z_2)^T$, define
\begin{equation}\label{tgframe:L0general}
 L_0=-\frac1{S_3(1+u_i^2)}
 \begin{pmatrix}z_1&-z_2\\z_2&z_1\end{pmatrix}
 \begin{pmatrix}1&u_i\\-u_i&1\end{pmatrix}.
\end{equation}
This is an explicit invertible linear map with
\[
 L_0[-S_3(1,u_i)^T]=z_0,\qquad
 \det L_0=\frac{|z_0|^2}{S_3^2(1+u_i^2)}>0,
\]
and inverse
\begin{equation}\label{tgframe:L0inverse}
 L_0^{-1}=-\frac{S_3}{|z_0|^2}
       \begin{pmatrix}1&-u_i\\u_i&1\end{pmatrix}
       \begin{pmatrix}z_1&z_2\\-z_2&z_1\end{pmatrix}.
\end{equation}
Indeed the second matrix of \eqref{tgframe:L0general} sends
$(1,u_i)^T$ to $(1+u_i^2,0)^T$, and each matrix times its
transpose is a positive scalar times $I_2$.
The same multiplication verifies the inverse and proves the exact norms
\begin{equation}\label{tgframe:L0norms}
 \|L_0\|=\frac{|z_0|}{S_3\sqrt{1+u_i^2}},\qquad
 \|L_0^{-1}\|=\frac{S_3\sqrt{1+u_i^2}}{|z_0|}.
\end{equation}
In particular the source target gives
\[
 L_0=-\frac{{\cal P}_\sigma(0)}{S_3(1+u_i^2)}
        \begin{pmatrix}1&u_i\\-u_i&1\end{pmatrix},
 \quad
 \|L_0\|=
 \frac{{\cal P}_\sigma(0)\,\mu\widehat\lambda_3^{k_3+6}}
 {A_0\sqrt{1+K_3^2/(\sigma_2^2n_i)}}.
\]
No seed or source initial amplitude has been set to one.

Define the two identity-initial-value fundamental matrices by
\[
 X_c'=\alpha C_cX_c,\quad X_c(0)=I_2,\qquad
 X_s'=(H-\delta_1I_2)X_s,\quad X_s(0)=I_2.
\]
Here is a direct existence, invertibility and uniqueness construction.
For a continuous matrix $D$ on $[0,L_s]$ let
\[
 X_D(v)=I_2+
 \sum_{j\ge1}\int_{0\le w_j\le\cdots\le w_1\le v}
       D(w_1)\cdots D(w_j)\,dw_j\cdots dw_1 .
\]
If $\sup\|D\|=M_D$, the $j$th term is bounded by
$(M_DL_s)^j/j!$. Uniform convergence and regrouping the integrals
give $X_D=I_2+\int_0^vD(w)X_D(w)\,dw$, hence the differential
equation. The corresponding right-ordered series for
$Z'=-ZD$, $Z(0)=I_2$, converges by the same bound.
Differentiation gives $(ZX_D)'=0$, hence $ZX_D=I_2$.
For square matrices this is invertibility. If two solutions agree
initially, the difference satisfies the zero-initial-value integral
equation; iteration bounds its norm by a fixed finite bound times
$(M_Dv)^j/j!$ for every $j$, which tends to zero.
This proves uniqueness. Smooth coefficient dependence follows from
the same series with the following explicit bound. On a compact
coefficient-parameter neighborhood let $M_N\ge1$ bound every coefficient
derivative through order $N$. Differentiating a product of $j$ factors
$N$ times allocates each derivative to one of $j$ factors, so there are
at most $j^N$ terms, each of norm at most $M_N^j$. The differentiated
$j$th integral is therefore bounded by
$j^N(M_NL_s)^j/j!$, a summable sequence. Uniform convergence of these
derivative series justifies termwise differentiation: integrate each
first-derivative series along a parameter segment and interchange the
uniform sum with that integral, then iterate through order $N$.
The right-ordered inverse series has the same bound. Pulse derivatives
then follow from the differential equations. This proves smoothness,
with the explicit derivative equations below providing further bounds.

Now set
\begin{equation}\label{tgframe:intertwiner}
 L(v)=X_s(v)L_0X_c(v)^{-1},\qquad
 L^{-1}=X_cL_0^{-1}X_s^{-1}.
\end{equation}
Differentiation gives
\[
 L'=(H-\delta_1I_2)L-\alpha LC_c.
\]
The original solution is $y_3(v)=X_c(v)y_3(0)$.
Consequently $L(v)y_3(v)=X_s(v)z_{\sigma,0}$, and
\begin{equation}\label{tgframe:exactpulsematch}
 B(v)L(v)y_3(v)=t^h_\sigma(v)
\end{equation}
by uniqueness of the source homogeneous equation and the exact initial
seed calculation. Thus this finite map reaches the actual source pulse.
It does not merely send the coupled pair into an unspecified transverse
solution. Multiplying both sides by the identical source cutoffs
preserves equality of the amplitudes, and applying identical linear
curl and pressure constructions preserves equality of the resulting
fields. Multiplication by the source prescribed scalar leading-wave
amplitude also preserves that equality. The cutoff principal residual
is then $\chi'BLy_3$; the remaining source construction uses its own
estimates for that exact source object.

\subsection{All scale costs of the finite intertwiner}

Here and below $\|\cdot\|$ is the Euclidean matrix norm.
Put $D_i=\operatorname{diag}(1,u_i)$,
$\kappa_i=\max(u_i,u_i^{-1})$, and
\[
 \beta_c(v)=\frac12\left(a_3(t_c(v))u_i+
                                  \frac{b_3(t_c(v))}{u_i}\right).
\]
The actual coefficient bounds $a_3\le a_i$, $b_3\le b_i$ and
$a_iu_i=b_i/u_i=\gamma_i$ give $0<\beta_c\le\gamma_i$.
Conjugating the coupled matrix by the constant $D_i$ gives
\[
 D_i^{-1}C_cD_i=
 \begin{pmatrix}-d_c&a_3u_i\\b_3/u_i&-d_c\end{pmatrix}.
\]
Its symmetric part has eigenvalues $-d_c\pm\beta_c$.
For any solution $w$ in these coordinates, differentiating $|w|^2$
therefore gives
\[
 2\alpha(-d_c-\beta_c)|w|^2
 \le (|w|^2)'
 \le2\alpha(-d_c+\beta_c)|w|^2.
\]
Multiplication by the corresponding scalar integrating factors proves
the following bounds, including the lower singular-value estimate
needed for the inverse:
\[
 \|X_c(v)\|\le\kappa_i
       e^{\int_0^v\alpha(\beta_c-d_c)\,dw},\qquad
 \|X_c(v)^{-1}\|\le\kappa_i
       e^{\int_0^v\alpha(\beta_c+d_c)\,dw}.
\]
For $D_s=H-\delta_1I_2=
\operatorname{diag}(\lambda,-\lambda)-\delta_1I_2+E$,
its symmetric part is bounded between
$(-\lambda-\delta_1-\|E\|)I_2$ and
$(\lambda-\delta_1+\|E\|)I_2$. The identical energy argument gives
\[
 \|X_s(v)\|\le
 e^{\int_0^v(\lambda-\delta_1+\|E\|)\,dw},\qquad
 \|X_s(v)^{-1}\|\le
 e^{\int_0^v(\lambda+\delta_1+\|E\|)\,dw}.
\]
Combining these four proved estimates with \eqref{tgframe:intertwiner}
yields
\begin{align}
 \|L(v)\|&\le\kappa_i\|L_0\|
 e^{\int_0^v[\lambda-\delta_1+\|E\|
                       +\alpha(\beta_c+d_c)]\,dw},\nonumber\\
 \|L(v)^{-1}\|&\le\kappa_i\|L_0^{-1}\|
 e^{\int_0^v[\lambda+\delta_1+\|E\|
                       +\alpha(\beta_c-d_c)]\,dw}.
 \label{tgframe:intertwinernorms}
\end{align}
All quantities on the right are the retained, finite coefficient paths.
In particular the exact coupled diffusion integral at the endpoint is
\[
 \int_0^{L_s}\alpha d_c\,dv
 =dK_3^2\int_0^{r_a}R_3(t_b+r)\,dr
 \in[dK_3^2r_a,\ 2dK_3^2r_a].
\]
The equality is the substitution $r=\alpha v$; the interval is the
proved actual phase-length bound. The determinant is also exact:
\begin{equation}\label{tgframe:detL}
 \det L(v)=\frac{|z_{\sigma,0}|^2}{S_3^2(1+u_i^2)}
 \exp\left(\int_0^v[
        \operatorname{tr}E-2\delta_1+2\alpha d_c]\,dw\right).
\end{equation}
For a differentiable $2$ by $2$ matrix, expanding its two-column
determinant and inserting $X'=DX$ gives
$(\det X)'=(\operatorname{tr}D)\det X$. Its initial determinant
is one. Apply this calculation to $X_s$ and $X_c$, and use
$\operatorname{tr}H=\operatorname{tr}E$ and
$\operatorname{tr}C_c=-2d_c$ to prove \eqref{tgframe:detL}.

Higher pulse derivatives have the explicit recurrence
\[
 L^{(n+1)}=
 \sum_{j=0}^n\binom nj\left[
 D_s^{(j)}L^{(n-j)}-L^{(n-j)}D_c^{(j)}\right],
 \qquad D_c=\alpha C_c .
\]
For a coefficient parameter derivative $\partial_\xi$ the exact
inhomogeneous equation is
\[
 (\partial_\xi L)'=D_s\partial_\xi L
       -(\partial_\xi L)D_c
       +(\partial_\xi D_s)L-L(\partial_\xi D_c).
\]
Its initial value is the derivative of the explicit
\eqref{tgframe:L0general}. To see directly that this produces a finite
bound, multiply on the left by $X_s^{-1}$ and on the right by $X_c$,
differentiate, and integrate the last two terms with these same
multipliers. The preceding norm estimates bound every resulting
integral on $[0,L_s]$. Iteration proves the analogous assertion for
every finite mixed derivative; all product-rule terms and initial
derivatives are retained. This also proves smoothness on each compact
source slow neighborhood where the frame is defined.

The factor $\mathcal P_\sigma(0)/S_3$ and the two explicit propagation
costs in \eqref{tgframe:intertwinernorms} can be very large or very small.
There is no supplied relation identifying $L_s$ with $L_3$.
Thus these equations prove a finite invertible map and exact matching
to the source object, while exposing the cost of that map. They do not
prove a uniform weighted-class estimate for $L$ or for the direct
candidate $L=I$. The source's particular pulse retains its own source
estimates because of the equality \eqref{tgframe:exactpulsematch}.
Transport of the full viscosity-one witness to another physical
viscosity uses the already proved affine orthogonal and viscosity map.
Alternatively one may replace $\delta_1$ by $\delta_\nu$ in this
finite ODE construction, but the resulting changed-background
homogeneous pulse is not asserted to be the source's published witness.


\section{Actual third-growth budgets in the physical curl and force}
\label{tgq:section}
We now use the actual third-growth pair and the original source frame
of the preceding calculations, with the direct coefficient choice
$L=I_2$. This fixes a particular candidate within the earlier general
bridge. All source band, box, sign and nonzero harmonic indices below
are retained, and only a finite collection is used on the present compact
interior patch. The cutoff of each term has compact support inside its
source rectangle and active annulus, with smooth extension by zero.
The original coupled datum, its selected return time $t_b$ and its
growth duration $r_a$ are fixed before assigning these source labels;
they have no dependence on the source spatial or auxiliary variables.
The source time and the coupled time have the exact dictionary proved
above:
\[
 t_{c,j}=t_b+\alpha_jv_j,\qquad
 \alpha_j=\frac{r_a}{L_{s,j}},\qquad
 \nabla_x t_{c,j}=0,\qquad
 \partial_t t_{c,j}=\alpha_jQ_j^{-1-h}.
\]
These derivative identities concern the complete physical phase
evaluation on each lifted rectangle. They remain valid when differentiating
a field supported strictly inside that rectangle: the field vanishes on
a collar where its extension would otherwise be needed.

Write $P(t_c)=-\Theta_3(t_c)$, $u(t_c)=\Omega_3(t_c)/\Theta_3(t_c)$.
The actual finite third growth gives
\[
 S_3\le P\le P_3^*,\qquad
 \frac{u_i}{\sqrt2}\le u\le2u_i,\qquad
 u_i=\sqrt{b_i/a_i}=\frac{K_3}{\sigma_2\sqrt{n_i}},
 \qquad \gamma_i=\sqrt{a_ib_i}.
\]
The original quantities remain
\[
 K_3=\mu\widehat\lambda_3,\qquad
 S_3=\frac{A_0}{\mu}\widehat\lambda_3^{-k_3-6},\qquad
 P_3^*=\frac{A_0}{\mu}\widehat\lambda_3^{-7/8},\qquad
 \nu=\sqrt{A_0}.
\]
In particular the following two component budgets have different
original frequency powers:
\begin{equation}
 \begin{split}
 Y_1&=P_3^*=\frac{A_0}{\mu}\widehat\lambda_3^{-7/8},\\
 Y_2&=2u_iP_3^*
      =\frac{2A_0}{\sigma_2\sqrt{n_i}}\widehat\lambda_3^{1/8},\\
 V_1&=2\gamma_iP_3^*,\qquad
 V_2=\gamma_i u_iP_3^*
      =\frac{\gamma_i A_0}{\sigma_2\sqrt{n_i}}
                       \widehat\lambda_3^{1/8}.
 \end{split}\label{tgq:budgets}
\end{equation}
Here the indices $1,2$ mean the two amplitude components, not physical
vector directions or radial moment weights. The derivative calculation
above proves
$|y_a(t_c)|\le Y_a$ and $|\dot y_a(t_c)|\le V_a$ on the entire
original interval, including its original negative seed.
For the second component at the actual threshold, the sharper two-sided
conclusion is
\[
 \frac{A_0}{\sigma_2\sqrt{2n_i}}\widehat\lambda_3^{1/8}
 \le |\Omega_3(t_3^a)|
 \le \frac{2A_0}{\sigma_2\sqrt{n_i}}\widehat\lambda_3^{1/8}.
\]
All these identities follow by multiplying the original quotient bounds
by $P$ and using $K_3=\mu\widehat\lambda_3$; the factor $\mu$ is
canceled only in the explicitly displayed product $u_iP_3^*$, not
removed from the underlying objects or their definitions.

\subsection{Two complete spatial shape fields for each retained harmonic}
\label{tgq:shapes}
For harmonic $j$, retain $k_jm_j\ne0$, where $k_j$ is the source
carrier and $m_j$ its integer harmonic. Neither equals the coupled
integer $k_3$ by definition. Let $B_j$ be the source frame, $n_j$ its
normal, and $\chi_j$ the full real slow, transverse and pulse cutoff.
All are the original specified functions on that label. Put, for
$a=1,2$ and the standard vectors $e_a\in\mathbb R^2$,
\begin{equation}
 \begin{split}
 T_{ja}&=\chi_jB_je_a,\qquad
 C_{ja}=\frac{i\,n_j\times T_{ja}}{k_jm_j|n_j|^2},\\
 W_{ja}&=\operatorname{Re}
       \bigl[(T_{ja}+\operatorname{curl}_{*,j,\mathrm{coeff}}C_{ja})
                              e^{ik_jm_j\Phi_j}\bigr],\\
 \varpi_{ja}
  &=\operatorname{Re}\left[
    \frac{i\{n_j^T\mathsf K_jT_{ja}-(n'_j)^TT_{ja}\}}
                {k_jm_j|n_j|^2}\,e^{ik_jm_j\Phi_j}\right],\\
 Z_{ja}&=Q_j^{-A}W_{ja},\qquad
 \Pi_{ja}=Q_j^{-2A}\varpi_{ja},\qquad A=\tfrac12+h .
 \end{split}\label{tgq:shape_def}
\end{equation}
Vectors in a cross product have their cylindrical components in the
oriented orthonormal physical basis. The coefficient curl is the full
operator from the preceding moment calculation. Equivalently,
\[
 Z_{ja}=\operatorname{Re}\operatorname{curl}_x
      \bigl[Q_j^{1/2-A}C_{ja}e^{ik_jm_j\Phi_j}\bigr].
\]
The transverse triple product proves this equivalence, including every
cutoff derivative. Thus each $Z_{ja}$ is a smooth compactly supported
divergence-free physical vector field; the same divergence identity
holds before auxiliary evaluation, with the commuting lifted
derivatives. The pressure coefficient is linear in $T_{ja}$ and includes
the full moving-normal term.

The actual added physical fields, with the two original amplitudes,
are exactly
\begin{equation}
 w_y=\sum_j\sum_{a=1}^2 y_a(t_{c,j}) Z_{ja},\qquad
 \pi_y=\sum_j\sum_{a=1}^2 y_a(t_{c,j})\Pi_{ja}.
 \label{tgq:factorization}
\end{equation}
Indeed $\nabla_xt_{c,j}=0$ gives
$\operatorname{curl}(y_a(t_{c,j})\mathcal A)=
y_a(t_{c,j})\operatorname{curl}\mathcal A$. It also gives
$\nabla(y_a(t_{c,j})\Pi)=y_a(t_{c,j})\nabla\Pi$.
The pressure formula in \eqref{tgq:shape_def} is linear, so it agrees
with the selected pressure for the amplitude $\chi_jB_jy(t_{c,j})$.
This proves \eqref{tgq:factorization}; in particular it does not discard
a spatial derivative of a rapidly growing amplitude. That derivative
has been calculated and is zero in these precise source coordinates.

The exact spatial derivative is
\[
 \partial_zw_y=\sum_{j,a}Q_j^{-A-D}y_a(t_{c,j})
          \partial_{Z_j}W_{ja},\qquad D=\tfrac12-h,
\]
where $\partial_{Z_j}W_{ja}$ differentiates its coefficient,
its exponential phase and all source frame and cutoff factors.
The absolute auxiliary evaluation map is independent of $z$.
There is no derivative of the actual coupled time in this formula.

\subsection{Quantitative radial moments with the original component powers}
\label{tgq:moment_budgets}
For a physical field $f$, use the precise norm
\[
 \|f\|_{[e]}^2=\int_0^\infty r^e\langle |f|^2\rangle\,dr,
 \qquad e=1,2,
\]
with the normalized angular and auxiliary averages of the moment
calculation. In a fixed band let
$\|F\|_{[e],j}^2=\int R_j^e\langle|F|^2\rangle\,dR_j$.
All norms are at the corresponding fixed slow axial coordinate and
physical time. Define the actual finite shape budgets
\begin{equation}
 \begin{split}
 \mathcal V_e&=
  \sum_{j,a}Y_a Q_j^{(e+1)/4-A}\|W_{ja}\|_{[e],j},\\
 \mathcal Z_e&=
  \sum_{j,a}Y_a Q_j^{(e+1)/4-A-D}
                              \|\partial_{Z_j}W_{ja}\|_{[e],j}.
 \end{split}\label{tgq:shape_budgets}
\end{equation}
The factors $Y_1,Y_2$ are the two explicit original-frequency expressions
in \eqref{tgq:budgets}, not a common unspecified constant. Triangle
inequality, the pointwise amplitude bounds, and
$r^e\,dr=Q_j^{(e+1)/2}R_j^e\,dR_j$ prove
\[
 \|w_y\|_{[e]}\le\mathcal V_e,\qquad
 \|\partial_zw_y\|_{[e]}\le\mathcal Z_e .
\]
No orthogonality of distinct harmonics has been assumed in these bounds.
One may evaluate every norm by the explicit source fields
\eqref{tgq:shape_def}; compactness makes each finite.

Let $w_*=u_*-\langle u_*\rangle_\theta$ be the actual old wave,
including its curl corrections, and set
\[
 \mathcal W_e=\|w_*\|_{[e]},\qquad
 \mathcal W_{e,z}=\|\partial_zw_*\|_{[e]},
\]
restricting the old-field radial integrals to a fixed enclosing support
of the new fields and their derivatives. The exact moments already
proved are
\begin{align*}
 M_2&=\partial_z\int r^2
   \langle w_{*,z}w_{y,\theta}+w_{y,z}w_{*,\theta}
                         +w_{y,z}w_{y,\theta}\rangle\,dr,\\
 M_1&=\partial_z\int r
   \langle2w_{*,z}w_{y,z}+w_{y,z}^2\rangle\,dr .
\end{align*}
Apply the product rule to every term and then Cauchy--Schwarz on its
original weighted product measure. This gives the concrete simultaneous
bounds
\begin{equation}
 |M_e|\le
 2\bigl(\mathcal W_{e,z}\mathcal V_e+
         \mathcal W_e\mathcal Z_e+
         \mathcal V_e\mathcal Z_e\bigr),\qquad e=1,2.
 \label{tgq:moment_bound}
\end{equation}
For $M_2$, the first two cross products contribute one term each to
each of the first two summands, and the self product contributes two
terms to the last. For $M_1$ the explicit coefficient $2$ and the
derivative of the square give exactly those same three coefficients.
Thus \eqref{tgq:moment_bound} retains the complete cross-and-quadratic
dependence on the original component budgets.

\subsection{The full physical residual with all clock and diffusion costs}
\label{tgq:full_residual}
Fix the actual finite source state $(u_*,p_*)$ used in the bridge and
an arbitrary physical viscosity $\nu_{\rm NS}>0$ in the residual.
It remains distinct from the coupled $d$ and $\nu=\sqrt{A_0}$.
Define the explicitly differentiated shape residual
\[
 \mathcal K_{ja}=
 \partial_tZ_{ja}+(u_*\cdot\nabla)Z_{ja}
 +(Z_{ja}\cdot\nabla)u_*-\nu_{\rm NS}\Delta Z_{ja}+\nabla\Pi_{ja}.
\]
With $C_3(t_c)=
\begin{pmatrix}-dK_3^2R_3&a_3\\ b_3&-dK_3^2R_3\end{pmatrix}$,
the full residual increment of \eqref{tgq:factorization} is
\begin{equation}
 \begin{split}
 \Delta R_y={}&
 \sum_{j,a} y_a(t_{c,j})\mathcal K_{ja}\\
 &+\sum_{j,a}\alpha_jQ_j^{-1-h}
                   (C_3y)_a(t_{c,j})Z_{ja}\\
 &+\sum_{j,a,k,b} y_a(t_{c,j})y_b(t_{c,k})
                                  (Z_{ja}\cdot\nabla)Z_{kb}.
 \end{split}\label{tgq:physical_residual}
\end{equation}
To prove it, substitute \eqref{tgq:factorization} into the complete
Navier--Stokes residual and subtract $R_{\nu_{\rm NS}}(u_*,p_*)$.
The time product rule gives
$\partial_ty_a=\alpha_jQ_j^{-1-h}(C_3y)_a$.
Every spatial derivative of $y_a(t_{c,j})$ vanishes, so neither
advection nor diffusion nor pressure adds an unrecorded derivative
of that factor. The three linear spatial terms, the time derivative
of $Z_{ja}$ and the pressure gradient are precisely $\mathcal K_{ja}$.
The remaining quadratic product is the last full sum in
\eqref{tgq:physical_residual}. This proves the identity with every
old-wave interaction, cutoff, curl term and pressure included.
Distinct labels with disjoint source supports have zero corresponding
products; no such cancellation is used for harmonics of the same label.
The baseline $R_{\nu_{\rm NS}}(u_*,p_*)$ must be added back to recover
the candidate's total residual.

The source units make every factor in this identity explicit.
Let $\nabla_{*,j}$ and $\Delta_{*,j}$ be the full cylindrical spatial
operators, including basis derivatives and the lifted auxiliary
derivatives, with
$\nabla_x=Q_j^{-1/2}\nabla_{*,j}$ and
$\Delta_x=Q_j^{-1}\Delta_{*,j}$ on that chart.
Put $U_{*,j}=Q_j^Au_*$ and
$\mathfrak t_{*,j}=Q_j^{1+h}\partial_t$ with the full physical
chain rule. Its extended expression is
$-\varepsilon_j\partial_{T_j}+c_{i,j}N_{i,j}$ as in source (6.6).
Define
\[
 \mathfrak K_{ja}=
 \mathfrak t_{*,j}W_{ja}
 +(U_{*,j}\cdot\nabla_{*,j})W_{ja}
 +(W_{ja}\cdot\nabla_{*,j})U_{*,j}
 -\nu_{\rm NS}\varepsilon_j\Delta_{*,j}W_{ja}
 +\nabla_{*,j}\varpi_{ja}.
\]
Each term is obtained by differentiating the explicit shapes in
\eqref{tgq:shape_def}, and
\[
 \mathcal K_{ja}=Q_j^{-2A-1/2}\mathfrak K_{ja}.
\]
For the time term the exponent identity is $A+1+h=2A+1/2$;
for viscosity it is
$Q_j^{-A-1}=Q_j^{-2A-1/2}\varepsilon_j$.
The nonlinear and pressure terms have the same stated residual unit.
This proves every factor rather than replacing the source or coupled
viscosity by a common normalized number.

Using Euclidean vector norms and the induced norm of the differentiated
vector map, \eqref{tgq:physical_residual} and
\eqref{tgq:budgets} give the following pointwise bound in actual
shape fields:
\begin{equation}
 \begin{split}
 |\Delta R_y|\le{}&
 \sum_{j,a}Q_j^{-2A-1/2}
       \bigl(Y_a|\mathfrak K_{ja}|+\alpha_jV_a|W_{ja}|\bigr)\\
 &+\sum_{j,a,k,b}Y_aY_bQ_j^{-A}Q_k^{-A-1/2}
                   |W_{ja}|\,\|\nabla_{*,k}W_{kb}\|.
 \end{split}\label{tgq:residual_bound}
\end{equation}
Both $Q_j$ and $Q_k$ remain in every mixed-band term.
The expression $\nabla_{*,k}W_{kb}$ includes the cylindrical frame
connection, so the bound controls the actual Cartesian directional
derivative. The clock satisfies the explicit original-stage bounds
\[
 \frac{L_3}{2\gamma_iL_{s,j}}\le\alpha_j
       \le\frac{8\sqrt2L_3}{\gamma_iL_{s,j}}.
\]
Consequently its derivative costs in the first line can, if desired,
be bounded using the original constants by
\[
 \alpha_jV_1\le
  \frac{16\sqrt2L_3}{L_{s,j}}
       \frac{A_0}{\mu}\widehat\lambda_3^{-7/8},\qquad
 \alpha_jV_2\le
  \frac{8\sqrt2L_3}{L_{s,j}}
       \frac{A_0}{\sigma_2\sqrt{n_i}}\widehat\lambda_3^{1/8}.
\]
These are finite quantitative costs for the actual inserted third stage.
They do not assert decay of the resulting residual at a singular limit.

The source pulse equation (7.5) itself belongs to its viscosity-one
source object. At a general $\nu_{\rm NS}$ in the displayed candidate
residual, the principal wave damping is
$\nu_{\rm NS}m_j^2\varepsilon_jk_j^2|n_j|^2$.
Alternatively, the exact viscosity transport already proved can be
applied to the complete viscosity-one source field and this complete
candidate. That transport rescales support, force, residual and
vorticity together. Merely changing this one damping coefficient
does not assert that the unchanged source background is a solution
with the same force at the changed viscosity.

Equations \eqref{tgq:shape_def}, \eqref{tgq:factorization},
\eqref{tgq:moment_bound} and \eqref{tgq:physical_residual} are the
explicit substitution of the actual third-growth amplitudes into
the source-coordinate construction. Every surviving moment and
residual has an actual differentiated field and a retained-scale
bound. The exact intertwiner above supplies the relation to the
source's homogeneous pulse; neither this direct candidate nor that
finite coordinate identity is silently promoted to a completed
infinite modified coupled construction.


\section{The actual wave increment and the source's three compatibility targets}
\label{momop:section}

The source is the retained 166-page manuscript. Its SHA-256 is given
in these two consecutive segments:
\[
\begin{gathered}
\texttt{0e779481c4da40bd28d1e642e1d8ca57}\\
\texttt{447d129610df28dfa5a11e9af8ae228f}.
\end{gathered}
\]
The source operations used here are (8.3)--(8.27) and the complete
four-step cycle of Proposition 9.6, on printed pages 89--111.
Its coefficient classes are Definitions 6.4--6.5 and Proposition 6.6,
printed pages 69--72. The source estimates quoted below are attributed
to those statements. The following calculation of the actual candidate
targets and all pressure-induced changes is a direct finite calculation.
It does not infer source iteration exponents from finite smoothness.
The source correction cycle is its viscosity-one cycle. The physical
increment identities below separately retain any fixed
$\nu_{\rm NS}>0$; using another physical viscosity in those identities
does not identify the resulting candidate with the source's witness.

\subsection{The actual covariance increment before pressure reconstruction}

Use right-handed cylindrical coordinates $(r,\theta,z)$. An overbar
means the normalized angular average followed by normalized auxiliary
Haar average, taken at fixed slow variables before phase evaluation.
The old source state has actual zero-angular-mean wave $w_*$, including
all its curl corrections. The new candidate is the proved finite field
\[
 w_y=\sum_{j,a}y_a(t_{c,j})Z_{ja},\qquad
 \pi_y=\sum_{j,a}y_a(t_{c,j})\Pi_{ja},\qquad
 \langle w_y\rangle_\theta=\langle\pi_y\rangle_\theta=0.
\]
The $Z_{ja},\Pi_{ja}$ are the complete physical shapes from the
third-growth quantitative calculation; $y_a$ is the actual coupled
pair, with $t_{c,j}=t_b+(r_a/L_{s,j})v_j$. No amplitude factor is
removed from an auxiliary average: it can depend on $v_j(Y)$.
Write the exact double-averaged covariance increment as
\begin{equation}\label{momop:C}
 C_{ab}=
 \overline{w_{*,a}w_{y,b}+w_{y,a}w_{*,b}+w_{y,a}w_{y,b}},
 \qquad a,b\in\{r,\theta,z\}.
\end{equation}
Equivalently, substitute the displayed finite sum for both occurrences
of $w_y$ inside this average. In particular $C$ retains every old-wave
cross term and every same-label quadratic term. It is symmetric.
All its radial integrals below are finite because the new fields are
smooth and supported in a compact interior annulus. The old field
need only be restricted to this support.

The old source finite correction state has the two exact mean-velocity
constraints
\[
 \int_0^\infty r^2\overline v\,dr=0,\qquad
 \int_0^\infty r\overline\gamma\,dr=0 .
\]
These are source (9.10) in physical variables. The present initial
addition $w_y$ leaves $(\beta,v,\gamma)$ unchanged because its angular
mean is zero. The temporal mean operation has zero auxiliary mean,
and the first two rows of the five-row operation below give zero
increments of these integrals. Thus all specified operations preserve
their initially zero values. Preservation, by itself, does not change
a nonzero entering value to zero.

The already calculated raw tangential residual moments are
\[
 M_2=\partial_zJ_2,\qquad M_1=\partial_zJ_1,\qquad
 J_2=\int_0^\infty r^2C_{z\theta}\,dr,\quad
 J_1=\int_0^\infty rC_{zz}\,dr.
\]
Here the candidate has been added with its zero-angular-mean pressure
$\pi_y$, so the mean pressure is still the old mean pressure.
The source algorithm next reconstructs the mean pressure from the
updated velocity. This operation changes the axial residual moment.

To compute it, first let $\mathcal C_{ab}$ denote the angular average
of the tensor inside \eqref{momop:C}, retaining its auxiliary dependence.
Let $\mathcal D_r$ be the physical radial derivative on the extended
domain, including the auxiliary phase-map derivative. Its auxiliary
average is ordinary $\partial_r$. Since the mean velocity is unchanged
by $w_y$, subtracting the radial conservative equations gives
\begin{equation}\label{momop:delta_g}
 \Delta g_r=-(\mathcal D_r+r^{-1})\mathcal C_{rr}
                   -\partial_z\mathcal C_{zr}
                   +r^{-1}\mathcal C_{\theta\theta}.
\end{equation}
Indeed the radial transport of a divergence-free velocity is
$(\mathcal D_r+r^{-1})u_r^2+
r^{-1}\partial_\theta(u_\theta u_r)+
\partial_z(u_zu_r)-u_\theta^2/r$.
Angular averaging removes the angular derivative and all products
containing exactly one old mean factor and one zero-angular-mean wave.
The source defines $g_r$ as the negative of the mean radial residual
apart from the pressure gradient, which gives exactly the signs in
\eqref{momop:delta_g}. The time and viscous averages of $w_y$
are zero, so no physical viscosity has been dropped in this subtraction.
The formula remains true for any fixed $\nu_{\rm NS}>0$ in the
physical residual.

Set $g=\overline{\Delta g_r}$. Auxiliary derivative integrals vanish,
and hence
\begin{equation}\label{momop:gbar}
 g=-(\partial_r+r^{-1})C_{rr}-\partial_zC_{zr}
                                    +r^{-1}C_{\theta\theta}.
\end{equation}
The increment of the source's pressure defect is therefore the
following explicit new scalar, rather than either raw tangential moment:
\begin{equation}\label{momop:deltaP}
 \Delta P=\int_0^\infty g\,dr
 =\int_0^\infty\frac{C_{\theta\theta}-C_{rr}}r\,dr
                           -\partial_z\int_0^\infty C_{zr}\,dr .
\end{equation}
The equality integrates the derivative $\partial_rC_{rr}$;
both endpoints vanish by compact support. The weighted second
moment of this same source is
\begin{equation}\label{momop:gsecond}
 \int_0^\infty r^2g\,dr
 =\int_0^\infty r(C_{rr}+C_{\theta\theta})\,dr
                   -\partial_z\int_0^\infty r^2C_{zr}\,dr .
\end{equation}
For its first term, integration by parts gives
$-\int r^2\partial_rC_{rr}=2\int rC_{rr}$, after which the
explicit $-C_{rr}/r$ contribution subtracts one copy.
These two identities retain the radial and axial coupling exactly.

\subsection{Compact mean pressure and the corrected axial flux moment}

Retain the source concentration scale and its full axial derivative:
\[
 z=q^D\eta,\qquad 1-t=q(1-\eta^2),\qquad
 A=\tfrac12+h,\quad D=\tfrac12-h,\quad
 q_z=\frac{2\eta q^{1-D}}{1-2h\eta^2}.
\]
Choose the source fixed bump $\widehat\rho$ in its reserved mean
interval $I_m$, with $\int\widehat\rho(x)\,dx=1$, and put
\[
 \rho_{\rm phys}=q^{-1/2}\widehat\rho(r/\sqrt q),\qquad
 c_{\rho,\rm phys}=\frac12\int r^2\rho_{\rm phys}\,dr
             =\frac q2\int x^2\widehat\rho(x)\,dx .
\]
Thus $(c_{\rho,\rm phys})_z=(q_z/q)c_{\rho,\rm phys}$;
it is generally not constant in $z$.

For precision the source's compact physical primitive is as follows.
The fixed slow cutoff $\chi_m$ is zero on the inner collar and one
on the outer collar of the active annulus. For an auxiliary-dependent
scalar $f$ define
\[
 \begin{split}
 I f(r,Y)&=\int_0^r
 f(s,Y+(s^{d_r}-r^{d_r})v_r)\,ds,\\
 J f(r,Y)&=\int_0^\infty
 f(s,Y+(s^{d_r}-r^{d_r})v_r)\,ds,\qquad
 I_cf=If-\chi_mJf.
 \end{split}
\]
Here $d_r=2((1+h)\rho_g-h\kappa_s)$ is the exact source exponent,
$\rho_g=\log(4-\sqrt2)/\log(4+\sqrt2)$,
$v_r=(1,-(\sqrt2-1))$, and $\kappa_s=10^{-5}$.
The derivative $\mathcal D_r$ of the shifted integral differentiates
the integration endpoint while cancelling the derivative of its
$-r^{d_r}v_r$ shift. Consequently
$\mathcal D_rI_cf=f-(\partial_r\chi_m)Jf$.
Also $\overline{Jf}=\int\overline f\,dr$, by translation invariance
of Haar measure. This proves the mean identities used next directly.

The source reconstruction after the actual wave addition changes the
physical mean pressure by
\[
 \Delta p_m=I_c(\Delta g_r-\rho_{\rm phys}\Delta P).
\]
Its source has zero radial auxiliary-mean integral by
\eqref{momop:deltaP}. Therefore
\begin{equation}\label{momop:pressureidentity}
 \partial_r\overline{\Delta p_m}
     =g-\rho_{\rm phys}\Delta P,\qquad
 \int r\,\overline{\Delta p_m}\,dr
     =-\frac12\int r^2g\,dr+c_{\rho,\rm phys}\Delta P .
\end{equation}
The first identity follows from the displayed primitive formula;
its cutoff remainder has zero auxiliary average. The second integrates
$\partial_r(r^2\overline{\Delta p_m})$ between the two zero endpoints.
The full radial residual increment after this pressure update is
$-\rho_{\rm phys}\Delta P-
(\partial_r\chi_m)J(\Delta g_r-\rho_{\rm phys}\Delta P)$.
The second term is retained. The source proves its terminal flatness
under its all-order uniform class bounds; finite smoothness alone is
not that proof.

Use different symbols for the source's two flux defects and the raw
candidate fluxes $J_2,J_1$. Their increments are exactly
\begin{align}
 \Delta J_\theta&=J_2,\nonumber\\
 \Delta J_z&=J_1-\frac12\int r^2g\,dr\nonumber\\
 &=\int r\left(C_{zz}-\frac12C_{rr}
                         -\frac12C_{\theta\theta}\right)\,dr
                    +\frac12\partial_z\int r^2C_{zr}\,dr .
 \label{momop:defectdictionary}
\end{align}
These are the differences of source (8.15): the mean-velocity terms
are unchanged and the full wave covariance changes by $C$.
The axial pressure identity \eqref{momop:pressureidentity} now proves
the exact moments after pressure reconstruction:
\begin{equation}\label{momop:afterpressure}
 M_{2,\rm new}=\partial_z\Delta J_\theta=M_2,\qquad
 M_{1,\rm new}=\partial_z(\Delta J_z+c_{\rho,\rm phys}\Delta P)
  =M_1+\partial_z\left(-\frac12\int r^2g\,dr
                                      +c_{\rho,\rm phys}\Delta P\right).
\end{equation}
These are increment identities. The old source moments and defects
are added to obtain the full updated state. No hypothesis that the
old defects vanish was used: identical old mean-velocity terms
cancel in the subtraction. The product involving $c_{\rho,\rm phys}$
stays inside $\partial_z$, including its nonzero derivative.

\subsection{New quantitative budgets for the three actual inputs}

Retain the previously proved quantities
$\mathcal W_e,\mathcal W_{e,z},\mathcal V_e,\mathcal Z_e$ for
$e=1,2$. They satisfy
\[
 \|w_*\|_{[e]}=\mathcal W_e,\quad
 \|\partial_zw_*\|_{[e]}=\mathcal W_{e,z},\quad
 \|w_y\|_{[e]}\le\mathcal V_e,\quad
 \|\partial_zw_y\|_{[e]}\le\mathcal Z_e.
\]
In particular $\mathcal V_e,\mathcal Z_e$ retain separately the
actual amplitude bounds
$Y_1=(A_0/\mu)\widehat\lambda_3^{-7/8}$ and
$Y_2=2A_0\widehat\lambda_3^{1/8}/(\sigma_2\sqrt{n_i})$
and their physical source-band factors. Define merely for the
following displayed estimates
\[
 B_e=2\mathcal W_e\mathcal V_e+\mathcal V_e^2,\qquad
 D_e=\mathcal W_{e,z}\mathcal V_e+
                    \mathcal W_e\mathcal Z_e+\mathcal V_e\mathcal Z_e.
\]
Choose $r_->0$ enclosing the radial support of the new fields and
their axial derivatives from below. Cauchy--Schwarz on the original
weighted product measures gives
\begin{equation}\label{momop:targetbounds}
 |\Delta P|\le r_-^{-2}B_1+2r_-^{-1}D_1,\qquad
 |\Delta J_\theta|\le B_2,\qquad
 |\Delta J_z|\le B_1+D_2.
\end{equation}
Here is the full estimate. The matrix
$\operatorname{diag}(-1,1,0)$ has operator norm one, so
$|C_{\theta\theta}-C_{rr}|$ is bounded after averaging by
$2|w_*||w_y|+|w_y|^2$. Replacing $r^{-1}$ by
$r_-^{-2}r$ and applying Cauchy--Schwarz gives the first part
$r_-^{-2}B_1$ of the pressure bound. Differentiating $C_{zr}$
produces the six actual cross and self product terms. Their integral
with weight $r^e$ is at most $2D_e$. Replace the unweighted
integral by $r_-^{-1}$ times the $r$-weighted integral to obtain
$2r_-^{-1}D_1$ in \eqref{momop:deltaP}.
The same product estimate without differentiation bounds
$\int r^2C_{z\theta}$ by $B_2$.
Finally $\operatorname{diag}(-1/2,-1/2,1)$ has norm one.
Apply it to the first integral in \eqref{momop:defectdictionary}
to get $B_1$; the last differentiated integral is bounded by
$\tfrac12(2D_2)=D_2$. This proves all three claims without
replacing the two amplitude components by a new unspecified constant.

\subsection{The exact five-row map for these targets}

The source map acts on the \emph{current} physical triple
$(P,J_\theta,J_z)$, after adding \eqref{momop:deltaP} and
\eqref{momop:defectdictionary} to the old source triple and
reconstructing pressure. It does not take
$(M_2,M_1)$ as its input. The five equations, in physical variables, are
\begin{align}
 \int r^2\Delta v\,dr&=0,&
 \int r\gamma_d\,dr&=0,\nonumber\\
 \int\frac{2V}{r}\Delta v\,dr&=-P,\nonumber\\
 \int r^2(G\Delta v+V\gamma_d)\,dr&=-J_\theta,\nonumber\\
 \int(2rG\gamma_d-rV\Delta v)\,dr&=-J_z .
 \label{momop:fivephysical}
\end{align}
In the last row there is exactly one factor $r$ multiplying $V\Delta v$.
The pressure-source term contributes
$-\tfrac12\int r^2(2V/r)\Delta v\,dr=-\int rV\Delta v\,dr$,
which proves that factor rather than treating the fifth row as an
independent constraint with an arbitrary weight.

Retain the source's reserved patch $x=r/\sqrt q\in I_m$. Its actual
summed base satisfies source (5.44) and (8.24):
\[
\begin{gathered}
 V_{\rm phys}=q^{-A}a(\eta)x^{-1-2\lambda},\qquad
 G_{\rm phys}=0,\\
 a(\eta)=2^{1/2+\lambda}c_{\rm patch}(1+\eta^2)^{-1},
 \quad \lambda>0,\quad |a(\eta)|\ge a_0>0.
\end{gathered}
\]
The source fixed $\lambda$ here is its reserved-profile exponent,
not a frequency of the coupled problem. In a fixed $Q$ chart,
the same physical base is
\[
 V_*=(Q/q)^A a(\eta)\bigl(\sqrt{Q/q}\,R\bigr)^{-1-2\lambda},
 \qquad G_*=0.
\]
The factor $(Q/q)^A$ has been retained.

Choose a nonnegative nonzero bump $\varphi_0\in C_c^\infty(I_m)$
supported near an interior point and a fixed spacing $d_b>0$ so small
that the three profiles
\[
 \varphi_j(x)=e^{-jd_b}\varphi_0(e^{-jd_b}x),\qquad j=0,1,2,
\]
have disjoint supports in the patch. This is possible because the patch
contains an open interval about that point: choose three distinct
scaled centers there, then decrease the support radius below half the
minimum separation. The symbol $d_b$ distinguishes this source bump
spacing from the coupled physical diffusivity $d$.
For every real $p$, set
\[
 \mu_p=\int x^p\varphi_0(x)\,dx>0,\qquad
 \int x^p\varphi_j(x)\,dx=\mu_pe^{jd_bp}.
\]
The equality is the substitution $x=e^{jd_b}\xi$ and retains the
prefactor in $\varphi_j$.

Define the two explicit moment matrices
\[
 A_\theta=
 \begin{pmatrix}
 \mu_2&\mu_2e^{2d_b}&\mu_2e^{4d_b}\\
 2a\mu_{-2-2\lambda}&
 2a\mu_{-2-2\lambda}e^{d_b(-2-2\lambda)}&
 2a\mu_{-2-2\lambda}e^{2d_b(-2-2\lambda)}\\
 -a\mu_{-2\lambda}&-a\mu_{-2\lambda}e^{-2d_b\lambda}&
                              -a\mu_{-2\lambda}e^{-4d_b\lambda}
 \end{pmatrix},
\]
\[
 A_z=
 \begin{pmatrix}
 \mu_1&\mu_1e^{d_b}\\
 a\mu_{1-2\lambda}&a\mu_{1-2\lambda}e^{d_b(1-2\lambda)}
 \end{pmatrix},\qquad a=a(\eta).
\]
Both inverses exist at every source parameter value. Indeed the
angular matrix is a diagonal row factor times the matrix with rows
$(1,t_p,t_p^2)$ for $p=2,-2-2\lambda,-2\lambda$ and
$t_p=e^{d_bp}$. Subtracting the first row from the other two
and factoring their differences gives
\[
 \det[1,t_p,t_p^2]_{p=p_1,p_2,p_3}
 =(t_{p_2}-t_{p_1})(t_{p_3}-t_{p_1})(t_{p_3}-t_{p_2}).
\]
All three powers are distinct for $\lambda>0$. The diagonal
row factor has determinant
\[
-2a^2\mu_2\mu_{-2-2\lambda}\mu_{-2\lambda}\ne0.
\]
Likewise
\[
\det A_z=a\mu_1\mu_{1-2\lambda}
(e^{d_b(1-2\lambda)}-e^{d_b})\ne0.
\]
This also gives an explicit finite inverse norm bound
$\|A^{-1}\|\le\|\operatorname{adj}A\|/|\det A|$
for either matrix. No uniform limit as $\lambda\downarrow0$ is taken.

Use the exact row normalizations
\[
 P_q=q^{2A}P,\qquad
 J_{\theta,q}=q^{2A-3/2}J_\theta,\qquad
 J_{z,q}=q^{2A-1}J_z,
\]
and define the five coefficients by
\begin{equation}\label{momop:coefficients}
 {\bf u}=A_\theta^{-1}(0,-P_q,-J_{z,q})^T,\qquad
 {\bf s}=A_z^{-1}(0,-J_{\theta,q})^T .
\end{equation}
Then the actual physical velocity increments are
\begin{equation}\label{momop:fiveincrement}
 \Delta v=q^{-A}\sum_{j=0}^2u_j\varphi_j(r/\sqrt q),
 \qquad
 \gamma_d=q^{-A}\sum_{j=0}^1s_j\varphi_j(r/\sqrt q).
\end{equation}
The first and second velocity moments have factors $q^{3/2-A}$
and $q^{1-A}$, respectively. The three inhomogeneous rows have
factors $q^{-2A}$, $q^{3/2-2A}$, and $q^{1-2A}$, respectively.
Substituting the original base into \eqref{momop:fivephysical}
and changing variables $r=\sqrt q\,x$ therefore gives exactly
the two linear systems \eqref{momop:coefficients}. This proves all
five equations, uniqueness in the selected five-dimensional span,
and smooth dependence on the actual target values. For instance
$\partial_\xi A^{-1}=-A^{-1}(\partial_\xi A)A^{-1}$ follows by
differentiating $AA^{-1}=I$; repeated product rules give every
coefficient derivative with the actual $q,\eta$ factors retained.

The desired axial increment has zero weighted integral and is
auxiliary independent. Its exact compact azimuthal potential and
radial increment are
\[
 \Psi(r,z,t)=\frac1r\int_0^r s\gamma_d(s,z,t)\,ds,\qquad
 \Delta\beta=-\partial_z\Psi,\qquad
 \Delta\gamma=(\partial_r+r^{-1})\Psi=\gamma_d.
\]
The zero moment in \eqref{momop:fivephysical} makes $\Psi$ vanish
beyond its outer support; its lower-end definition makes it vanish
below the inner support. Differentiation proves
$(\partial_r+r^{-1})\Delta\beta+\partial_z\Delta\gamma=0$.
The support includes the intervals between bump supports, which still
lie within the reserved patch.
To display the axial scaling cost explicitly, put
$\Psi_j(x)=x^{-1}\int_0^x \xi\varphi_j(\xi)\,d\xi$. Then
\begin{align}
 \Psi&=q^{1/2-A}\sum_{j=0}^1s_j\Psi_j(x),\nonumber\\
 \Delta\beta
 &=-q^{1/2-A}\sum_{j=0}^1\left[
    (\partial_zs_j)\Psi_j+
    \frac{q_z}{q}s_j\left((1/2-A)\Psi_j
                                      -\frac x2\Psi_j'\right)\right].
 \label{momop:inducedradial}
\end{align}
The two individual $\Psi_j$ can have $x^{-1}$ tails. The exact
linear relation $\sum_js_j\mu_1e^{jd_b}=0$ cancels those tails
in their sum and, after differentiation, in \eqref{momop:inducedradial}.
No individual compactness was assumed.

\subsection{What the five-row operation actually leaves}

Let $(\beta,v,\gamma)$ be the current mean correction, after whatever
preceding wave and mean updates have actually been performed. Keep the
base $(b,V,G)$ and the full current wave $w$ fixed during this
five-row step. Use the three full current defects as its targets.
Write $\mathcal D_r$ for the physical extended radial derivative and
$\mathcal D_t$ for the physical extended time derivative. With
$\Delta_0=\mathcal D_r^2+r^{-1}\mathcal D_r+\partial_z^2$, direct
subtraction of the full radial source gives
\begin{align}
 R_g:={}&g_{r,\rm new}-g_r-\frac{2V}{r}\Delta v\nonumber\\
 ={}&-\mathcal D_t\Delta\beta
 -(\mathcal D_r+r^{-1})
             \bigl(2b\Delta\beta+2\beta\Delta\beta+(\Delta\beta)^2\bigr)
 \nonumber\\
 &-\partial_z\bigl(b\Delta\gamma+G\Delta\beta+
           \beta\Delta\gamma+\gamma\Delta\beta+\Delta\beta\Delta\gamma\bigr)
 \nonumber\\
 &+\frac{2v\Delta v+(\Delta v)^2}{r}
               +\nu_{\rm NS}(\Delta_0-r^{-2})\Delta\beta .
 \label{momop:Rg}
\end{align}
Here all increments are auxiliary independent, so their
$\mathcal D_t$ is ordinary $\partial_t$. The products with the
old means still retain their auxiliary dependence. Equation
\eqref{momop:Rg} keeps the actual physical viscosity; the source
unit-viscosity formula (8.26) is its case $\nu_{\rm NS}=1$.

The full updated defects, with pressure recomputed from the new
radial source, are exactly
\begin{align}
 P_{\rm new}&=\int\overline{R_g}\,dr,\nonumber\\
 (J_\theta)_{\rm new}
    &=\int r^2\overline{\gamma\Delta v+v\Delta\gamma+
                                      \Delta\gamma\Delta v}\,dr,\nonumber\\
 (J_z)_{\rm new}
    &=\int r\overline{2\gamma\Delta\gamma+(\Delta\gamma)^2}\,dr
                      -\frac12\int r^2\overline{R_g}\,dr .
 \label{momop:recomputed}
\end{align}
To verify these formulas, expand every quadratic product in the
definitions of the current and new defects. The radial source
increment is $(2V/r)\Delta v+R_g$. The third row of
\eqref{momop:fivephysical} cancels the old $P$ with the integral
of the displayed linear term. The azimuthal flux increment contains
the base-linear terms $G\Delta v+V\Delta\gamma$, cancelled by
the fourth row because $\Delta\gamma=\gamma_d$. The axial defect
contains $2G\Delta\gamma$ and the radial-source contribution
$-\tfrac12r(2V\Delta v)$ after integration. These are precisely
the fifth row. All remaining products are the terms displayed
in \eqref{momop:recomputed}; the full wave covariance is unchanged
during this mean step. This proves the formulas without assuming
their right-hand sides vanish.

The two velocity moments remain exactly zero. The new residual
moments therefore have the source form
\[
 \int r^2\overline{E_{\theta,\rm new}}\,dr
       =\partial_z(J_\theta)_{\rm new},\qquad
 \int r\overline{E_{z,\rm new}}\,dr
       =\partial_z\bigl((J_z)_{\rm new}
                           +c_{\rho,\rm phys}P_{\rm new}\bigr).
\]
The pressure and $c_{\rho,\rm phys}$ are the actual reconstructed
objects. Thus the five equations cancel specified linear defect
contributions and expose the remaining nonlinear and differential
terms. Equations \eqref{momop:Rg}--\eqref{momop:recomputed} supply
the next actual inputs rather than a claim of complete cancellation.

\subsection{Exact source operation dictionary and the bounds it requires}

For the source's estimates, $Q=2^{-\ell}$, $\varepsilon=Q^h$,
$S_*=\ell^2$, $R=r/\sqrt Q$, $Z=z/Q^D$, $T=(1-t)/Q$.
Use a coarsest applicable common auxiliary torus $H$ as in source
Lemma 6.2; different representatives agree by its integer covering
maps. The original radial edge weights are
\[
\begin{gathered}
 \zeta(X)=\exp\left(-\frac{a_a}{\log^2(X/X_a)}
                       -\frac{a_b}{\log^2(X_b/X)}\right),\\
 \delta(X)=\min\{1,\log(X/X_a),\log(X_b/X)\},
 \quad X=\frac{QR^2}{2q},
\end{gathered}
\]
with the source fixed $a_a,a_b>0$. In the following bounds the band,
label, representative and harmonic are held fixed when differentiating.
The constants may depend on the finite correction stage and derivative
multi-index, but are uniform over bands, labels, points and rectangle
copies.

A mean coefficient $f\in M^\alpha$ has no angular dependence, descends
to the common torus, extends smoothly by zero outside the active shell,
and satisfies for each multi-index $I$ in $(R,Z,T,H)$
\[
 |\partial^If|\le C_I\varepsilon^\alpha S_*^{b_I}
                                        \zeta\,\delta^{-d_I}.
\]
It can have nonzero auxiliary modes and has no rectangle or pulse
envelope restriction. A moment coefficient $F(Z,T)\in S^\alpha$
satisfies $|\partial_{Z,T}^IF|\le C_I\varepsilon^\alpha S_*^{b_I}$.
A labelled nonzero harmonic amplitude $a_{\gamma,m}\in W^\alpha$
has no angular dependence, the source slow, transverse and shell
support, smooth zero extension across those boundaries, compatible
lift representatives, and
\[
 |\partial^Ia_{\gamma,m}|\le
 C_I\varepsilon^\alpha S_*^{b_I}\sqrt\zeta\,\delta^{-d_I}
                          {\cal P}_\gamma(v),\qquad 0\le v\le L_s.
\]
The prescribed source phase has $k_\gamma p_\gamma\in
\mathbb Z\setminus\{0\}$ and $k_\gamma=\lceil\varepsilon^{-1/2}\rceil$.
The harmonic set is finite at a fixed stage and independent of band.
Equal label--harmonic terms are grouped before testing the bounds.
The pulse envelope is the original positive exponential in (7.12),
with its proved Gaussian estimates. Actual cut-off waves additionally
have support where the source pulse cutoff $\psi$ is supported and
smooth zero extension in $v$. No divisibility by $\psi$ is required.

The source's anisotropic derivative mappings, with the phase factored
out, are
\[
\begin{gathered}
 \partial_{R,Z,T}:C^\alpha\to C^\alpha,\quad
 D_r:C^\alpha\to C^{\alpha-\kappa_s},\quad
 D_z=\varepsilon\partial_Z:C^\alpha\to C^{\alpha+1},\\
 c_{i_0}N_{i_0}:C^\alpha\to C^\alpha,\quad
 -\varepsilon\partial_T:C^\alpha\to C^{\alpha+1},
 \qquad C=M,W.
\end{gathered}
\]
Products obey $M^\alpha M^\beta\subset M^{\alpha+\beta}$ and
$M^\alpha W^\beta\subset W^{\alpha+\beta}$. Same-label wave
products have this summed order, in $M$ for zero output harmonic
and $W$ otherwise; distinct supported labels have zero products.
The zero-harmonic global $M$ statement here is for the actual temporally
cut-off waves. Before that cutoff, the source gives the mean bounds on
the labelled rectangle, not a global $M$ assertion.
These are source Proposition 6.6, not deductions from the candidate's
finite norm bounds.

Every radial target must first carry the correct measure and weight.
For $f_*=Q^{a_f}f_{\rm phys}$,
\[
 \int R^e\overline{f_*}\,dR
       =Q^{a_f-(e+1)/2}\int r^e\overline{f_{\rm phys}}\,dr .
\]
For the present inputs this reads
\begin{align*}
 P_*&=Q^{2A}P,&
 (J_\theta)_*&=Q^{2A-3/2}J_\theta,&
 (J_z)_*&=Q^{2A-1}J_z,&
 (c_\rho)_*&=Q^{-1}c_{\rho,\rm phys},\\
 (M_2)_*&=Q^{2A-1}M_2,&
 (M_1)_*&=Q^{2A-1/2}M_1 .
\end{align*}
In particular \eqref{momop:afterpressure} becomes exactly
$\varepsilon\partial_Z\Delta J_{\theta,*}$ and
$\varepsilon\partial_Z(\Delta J_{z,*}+c_{\rho,*}\Delta P_*)$.
The factors match because $h+D=1/2$. This is the extra axial
factor used in source (9.12); using the raw $J_1$ after pressure
reconstruction would omit part of its actual input.

For clarity, the following are the four ordered operations of source
Proposition 9.6, with the full current velocity and pressure
reconstructed after each operation.

\begin{enumerate}
\item For each supported nonzero harmonic coefficient $f_{\gamma,m}$
of the full current residual, solve the original transverse pulse
equation with zero entrance data and source $-f_{\gamma,m}$.
In frame coordinates its coefficient matrix is
$H_\gamma-m^2\varepsilon k_\gamma^2|n_\gamma|^2I$.
Multiply the solved potential and pressure by $\psi$ and take the
full curl. Source Proposition 7.2 maps $W^\alpha$ sources to
$W^\alpha$ velocity amplitudes and $W^{\alpha+1/2}$ pressure.
Proposition 9.1 assigns the remaining supported linear residual
order $\alpha+1/2-3\kappa_s$ and retains its Gaussian cutoff tails
additively. The actual candidate's full residual formula determines
the input coefficients. Following this wave change, replace $w_y$
by the actual total new wave increment in \eqref{momop:C}, recompute
$g_r$ and reconstruct pressure. Its old-wave cross terms must change
with the updated wave.

\item Let $F_2,F_1$ be the auxiliary-averaged physical tangential
residuals of that updated state. Their two moments are the updated
versions of \eqref{momop:afterpressure}, including the old defects.
For each $e=1,2$ retain the fixed normalized physical bump
$b_e=q^{-(e+1)/2}\widehat b_e(r/\sqrt q)$ with
$\int x^e\widehat b_e(x)\,dx=1$. Form
\[
 \mathfrak M_e=\int r^eF_e\,dr,\qquad
 \sigma_e=-r^{-e}\int_0^r s^e
                         (F_e(s)-b_e(s)\mathfrak M_e)\,ds .
\]
The zero total weighted integral proves compact support and
$(\partial_r+e/r)\sigma_e=-F_e+b_e\mathfrak M_e$.
Use $\Sigma=Q^{2A}(\sigma_2,\sigma_1)$ as the input to the
source's signed amplitude derivative (7.31), with its fixed
positive primary amplitudes. Source Proposition 7.6 maps an
auxiliary-independent $M^\alpha$ stress to $W^{\alpha-1/2}$.
Take its actual curl, retain the old-wave cross and quadratic
correction tensor (9.13), and reconstruct pressure. The source
does not set $b_e\mathfrak M_e$ to zero.

If $\mathfrak M_e=\partial_z\mathfrak J_e$, its bump term satisfies
the exact derivative identity
\[
 b_e\partial_z\mathfrak J_e
 =\partial_z(b_e\mathfrak J_e)-(\partial_zb_e)\mathfrak J_e,\qquad
 \partial_zb_e=-\frac{q_z}{2q}
                    \bigl((e+1)b_e+r\partial_rb_e\bigr).
\]
It follows by differentiating the original $q$-scaled bump.
Thus trading the moment bump for an axial tensor flux leaves the
displayed radial-profile derivative term. For the axial equation,
$\mathfrak J_1=J_z+c_{\rho,\rm phys}P$, including both product
derivatives in \eqref{momop:afterpressure}.

\item For the zero-auxiliary-mean part of the current angularly
invariant residual, first convert the physical residual to the fixed
source chart. In this step a star denotes that chart representative:
\[
\begin{gathered}
 E_*=Q^{2A+1/2}E_{\rm phys},\qquad
 E_{a,*}^\circ=E_{a,*}-\langle E_{a,*}\rangle_Y,\\
 (\Delta\beta_*,\Delta v_*,\Delta\gamma_*)
       =Q^A(\Delta\beta_{\rm phys},\Delta v_{\rm phys},
                                      \Delta\gamma_{\rm phys}).
\end{gathered}
\]
The source temporal update in these exact units is
\[
 \Delta v_*=-c_{i_0}^{-1}N_{i_0}^{-1}E_{\theta,*}^\circ,\qquad
 \gamma_{d,*}=-c_{i_0}^{-1}N_{i_0}^{-1}E_{z,*}^\circ,
 \quad
 N_{i_0}^{-1}F=
 \sum_{k\ne0}\frac{\widehat F(k)}{2\pi i v_t\cdot k}e^{2\pi ik\cdot H}.
\]
Here $v_t=(\sqrt2-1,1)$ and
$c_{i_0}=T_g^{i_0}Q^{1+h}$, with $c_{i_0}^{-1}\le CS_*$.
The source Diophantine bound gives
$\|N^{-1}F\|_{C^m}\le C_m\|F\|_{C^{m+4}}$.
Let $T_{1,*}$ and $A_{1,*}$ be the source's radial potential and
its cutoff remainder in this chart. Realize $\gamma_{d,*}$ by
$T_{1,*}\gamma_{d,*}$:
$\Delta\beta_*=-\varepsilon\partial_ZT_{1,*}\gamma_{d,*}$ and
$\Delta\gamma_*=\gamma_{d,*}-A_{1,*}\gamma_{d,*}$.
The actual axial remainder $-A_{1,*}\gamma_{d,*}$ stays in the
new velocity. The exact normalized cancellations are
\[
\begin{split}
 c_{i_0}N_{i_0}\Delta v_*&=-E_{\theta,*}^\circ,\\
 c_{i_0}N_{i_0}\Delta\gamma_*&=-E_{z,*}^\circ
              -c_{i_0}N_{i_0}A_{1,*}\gamma_{d,*}.
\end{split}
\]
Convert every resulting velocity component back by $Q^{-A}$ before
reconstructing the physical pressure or taking physical products and
defects. In particular its physical angular update is exactly
\[
 \Delta v_{\rm phys}
  =-Q^{-A}c_{i_0}^{-1}Q^{2A+1/2}
                     N_{i_0}^{-1}E_{\theta,\rm phys}^\circ
  =-T_g^{-i_0}N_{i_0}^{-1}E_{\theta,\rm phys}^\circ .
\]
The last equality uses $A+1/2=1+h$ and the original
$c_{i_0}=T_g^{i_0}Q^{1+h}$. It is also
$-N_{\rm abs}^{-1}E_{\theta,\rm phys}^\circ$, since
$N_{\rm abs}=T_g^{i_0}N_{i_0}$ on the representative.
This proves the conversion without applying a normalized inverse
directly to a physical residual.

\item Apply \eqref{momop:coefficients} to the three defects of this
updated state, use \eqref{momop:fiveincrement} and its actual
radial component \eqref{momop:inducedradial}, and reconstruct
pressure again. The exact remaining defects are
\eqref{momop:Rg}--\eqref{momop:recomputed}.
Source Lemma 8.7 maps $S^\alpha$ targets to $M^\alpha$
tangential increments and $M^{\alpha+1}$ radial increments.
Source Lemma 8.8 gives the remaining defect class
$S^{\alpha+0.9-2\kappa_s}$ using the actual bounds
$v,\gamma\in M^{0.9}$, $\beta\in M^{1.9}$, $\alpha\ge0.9$,
and its reserved-patch base derivative bounds.
\end{enumerate}

The source cycle assumes its full finite correction-state bounds
$w\in W^{1/2}$, $w-w_0^{\tan}\in W^{0.68}$,
$v,\gamma,p_m\in M^{0.9}$, $\beta\in M^{1.9}$,
with supported nonzero residual order
$B_j=1/2+\sigma_j$, mean and defect order $C_j^*=1+\sigma_j$,
and $\sigma_j=1/5+j/10$. Its ordered estimate proves gains
$B_{j+1}-B_j=C_{j+1}^*-C_j^*=1/10$.
These are requirements and conclusions for the source's construction.
The explicit candidate budgets and equations above do not by
themselves identify those exponents.

There is a precise finite-domain statement available now. Write the
full current residual as $R_{\rm old}+\Delta R$, retaining the old
source state and every term of the computed increment. The added
candidate uses finitely many labels, with its cutoffs supported
strictly inside the source rectangle and radial annulus. On those
compact coefficient supports, $\varepsilon>0$, $\zeta>0$,
$\delta>0$ and $\mathcal P_\gamma(v)>0$ have positive minima
for each fixed label. By local finiteness of the retained source-label
family, only finitely many old labels intersect this fixed compact
perturbation support. Every differentiated coefficient of $\Delta R$
is finite by its explicit smooth shapes, including the retained old-state
cross terms there. For any fixed
$\alpha$ and derivative $I$, the maximum of its derivative divided by
$\varepsilon^\alpha\sqrt\zeta\mathcal P_\gamma$ on these finitely
many supports is therefore a concrete finite constant. Extend the
finite \emph{increment} family by zero on unused labels. This proves its
finite coefficient inequalities without assigning an asymptotic gain
to the amplitude. It does not set any old source coefficient to zero:
$R_{\rm old}$ retains its actual locally finite family and its attributed
source-class bounds. The finite-sum closure in source Proposition 6.6
combines those bounds with the proved increment bounds for each full
current input. The same compact-support argument applies to the added
mean sources and the reserved bump operators. Thus the operations
above are defined on the actual old-plus-increment inputs.
Their constants can contain the inverse of a very small Gaussian
envelope or source scale; no uniform control for an increasing family
of coupled stages or bands follows from this compactness argument.
The original component budgets and the full nonlinear remainders
remain the quantities to propagate when investigating such a family.


\section{Compact completion of the full averaged residual}
\label{mcstress:section}

This section computes a compact symmetric stress and a compact pressure
for the complete averaged residual of the actual finite curl candidate.
It retains the original radial bumps depending on $q(z,t)$, both axial
fluxes, the radial force, and the full cylindrical connection terms.
The result is a stress--pressure representation of a computed force;
it does not by itself modify a velocity or identify the constructed
stress with a covariance of additional waves.

\subsection{The actual extended fields and all three averaged components}

Use right-handed cylindrical coordinates
$x=(r\cos\theta,r\sin\theta,z)$, with oriented orthonormal basis
$(e_r,e_\theta,e_z)$ and $r>0$. The retained finite source state is
$(u_*,p_*)$; the actual added field and its selected pressure are
$(v,\pi)=(w_y,\pi_y)$. On a finite collection of source labels the
added potential has the specified form
\begin{equation}\label{mcstress:actual_potential}
\begin{gathered}
 \mathcal A_y=\sum_j\operatorname{Re}
   \left[Q_j^{1/2-A}
     \frac{i\,n_j\times t_j}{k_jm_j|n_j|^2}
                      e^{ik_jm_j\Phi_j}\right],\qquad
 v=\operatorname{curl}_{\mathscr D}\mathcal A_y,\\
 t_j=\chi_jB_jy_3(t_{c,j}),\qquad
 y_3=\binom{\Theta_3}{\Omega_3},\qquad
 t_{c,j}=t_b+\alpha_jv_j,\qquad
 n_j^TB_j=0,\qquad k_jm_jp_j\in\mathbb Z\setminus\{0\}.
\end{gathered}
\end{equation}
Here $y_3$ is the actual signed third-growth trajectory, the cutoffs
and source frame are the retained functions, and $Q_j$ is constant
on each source chart. The chosen pressure is the real harmonic sum
with chart coefficient
\[
 Q_j^{-2A}
 \frac{i\{n_j^T\mathsf K_jt_j-(n'_j)^Tt_j\}}
      {k_jm_j|n_j|^2}e^{ik_jm_j\Phi_j}.
\]
All coefficient components are independent of $\theta$; their angular
dependence is the displayed nonzero integer harmonic. Their radial
and axial supports lie in a fixed compact set with $r$ bounded away
from zero on the finite time interval under consideration. Smooth
zero extension at each source cutoff is retained.

The source absolute auxiliary coordinate is $Y\in\mathbb T^2$ and
the physical evaluation is
\begin{equation}\label{mcstress:extended_derivatives}
\begin{gathered}
 Y(r,t)=v_r r^{d_r}+v_t t\pmod{\mathbb Z^2},\qquad
 v_r=(1,-b_g),\quad v_t=(b_g,1),\quad b_g=\sqrt2-1,\\
 \rho_g=\frac{\log(4-\sqrt2)}{\log(4+\sqrt2)},\quad
 \kappa_s=10^{-5},\quad d_r=2((1+h)\rho_g-h\kappa_s),\\
 \mathscr D_r=\partial_r+d_r r^{d_r-1}v_r\cdot\partial_Y,
 \qquad \mathscr D_t=\partial_t+v_t\cdot\partial_Y,
 \qquad \mathscr D_z=\partial_z.
\end{gathered}
\end{equation}
The angular derivative is $\partial_\theta$.
The exponent $h$ here is the unchanged source chart exponent. These
lifted derivatives commute, obey the product rule, and satisfy
$\mathscr D_r r=1$. Equivalently in Cartesian coordinates they are
$\partial_{x_i}+\sum_\alpha(\partial_{x_i}Y_\alpha)\partial_{Y_\alpha}$;
their commutators vanish because mixed derivatives of $Y$ commute.
Thus evaluation at the physical map preserves every differentiated
identity below.

The actual finite source has its retained representation
$u_* =\operatorname{curl}_{\mathscr D}\mathcal A_*+B_*e_\theta$,
with $\partial_\theta B_*=0$. Commutation of the lifted derivatives
proves $\operatorname{div}_{\mathscr D}\operatorname{curl}_{\mathscr D}=0$.
Also $\operatorname{div}_{\mathscr D}(B_*e_\theta)=0$.
Consequently both $u_*$ and $v$ are divergence free at every independent
$Y$, not only after physical evaluation.

Let the normalized averages act on cylindrical components before
evaluation:
\[
 \langle f\rangle_\theta=\frac1{2\pi}\int_0^{2\pi}f\,d\theta,
 \quad \langle f\rangle_Y=\int_{\mathbb T^2}f\,dY,
 \quad \langle f\rangle=\langle\langle f\rangle_\theta\rangle_Y,
 \quad\int_{\mathbb T^2}1\,dY=1.
\]
Every cylindrical curl operation retains the nonzero angular harmonic
in \eqref{mcstress:actual_potential}. For example its radial and axial
components are $r^{-1}\partial_\theta\mathcal A_z-\partial_z\mathcal A_\theta$
and $(\mathscr D_r+r^{-1})\mathcal A_\theta-r^{-1}\partial_\theta\mathcal A_r$.
Integrating the integer harmonic proves
\begin{equation}\label{mcstress:zero_mean}
 \langle v\rangle_\theta=0,\qquad
 \langle\pi\rangle_\theta=0
 \quad\hbox{at every }(r,z,t,Y).
\end{equation}
This statement includes the complete coefficient curl and every cutoff
derivative, rather than only the principal amplitude.

For a smooth periodic coefficient $f$, integration of its auxiliary
derivatives gives
\begin{equation}\label{mcstress:average_derivatives}
 \langle\mathscr D_r f\rangle=\partial_r\langle f\rangle,
 \quad \langle\mathscr D_t f\rangle=\partial_t\langle f\rangle,
 \quad \langle\partial_z f\rangle=\partial_z\langle f\rangle,
 \quad \langle\partial_\theta f\rangle=0.
\end{equation}
Applying the first identity twice also treats $\mathscr D_r^2$;
the derivative of its variable coefficient has zero auxiliary average
because that coefficient is independent of $Y$.

At arbitrary fixed physical viscosity $\nu_{\rm NS}>0$, the exact
residual increment is
\begin{equation}\label{mcstress:residual_increment}
 \Delta R=\mathscr D_t v+(u_*\cdot\nabla_{\mathscr D})v
 +(v\cdot\nabla_{\mathscr D})u_*
 +(v\cdot\nabla_{\mathscr D})v
 -\nu_{\rm NS}\Delta_{\mathscr D}v+\nabla_{\mathscr D}\pi.
\end{equation}
This is the residual of the actual modified state minus the source
state's residual at the same viscosity. Put
\begin{equation}\label{mcstress:actual_tensor}
 \bar u_* =\langle u_*\rangle_\theta,\qquad
 w_*=u_*-\bar u_*,\qquad
 S=\langle w_*\otimes v+v\otimes w_*+v\otimes v\rangle.
\end{equation}
It is a smooth real symmetric axisymmetric tensor compactly supported
in the annular and axial set at issue. Each summand contains a factor
of the actual compact field $v$.

\begin{theorem}\label{mcstress:full_mean}
The complete averaged residual is the full Cartesian divergence of
the actual tensor $S$:
\begin{equation}\label{mcstress:Ffull}
 \langle\Delta R\rangle
 =F_r e_r+F_2e_\theta+F_1e_z
 =\operatorname{div}S,
\end{equation}
where
\begin{equation}\label{mcstress:all_components}
\begin{aligned}
 F_r&=(\partial_r+r^{-1})S_{rr}-r^{-1}S_{\theta\theta}
                                      +\partial_zS_{rz},\\
 F_2&=(\partial_r+2r^{-1})S_{r\theta}+\partial_zS_{z\theta},\\
 F_1&=(\partial_r+r^{-1})S_{rz}+\partial_zS_{zz}.
\end{aligned}
\end{equation}
\end{theorem}
\begin{proof}
For divergence-free fields, direct Cartesian component expansion
rewrites the three advection terms of \eqref{mcstress:residual_increment}
as the lifted divergence of
$u_*\otimes v+v\otimes u_*+v\otimes v$.
The mean-flow cross tensor has zero angular average by
\eqref{mcstress:zero_mean}, even when $\bar u_*$ depends on $Y$.
Thus its averaged divergence is zero by
\eqref{mcstress:average_derivatives} and the cylindrical divergence
formula.

To check the linear radial component explicitly, let
\[
 \mathscr L_0=\mathscr D_r^2+r^{-1}\mathscr D_r
             +r^{-2}\partial_\theta^2+\partial_z^2.
\]
The complete cylindrical vector Laplacian is
\[
 \begin{aligned}
 (\Delta_{\mathscr D}v)_r&=(\mathscr L_0-r^{-2})v_r
                       -2r^{-2}\partial_\theta v_\theta,\\
 (\Delta_{\mathscr D}v)_\theta&=(\mathscr L_0-r^{-2})v_\theta
                       +2r^{-2}\partial_\theta v_r,\qquad
 (\Delta_{\mathscr D}v)_z=\mathscr L_0v_z.
 \end{aligned}
\]
Each average is zero by \eqref{mcstress:zero_mean} and
\eqref{mcstress:average_derivatives}. The radial pressure average is
$\langle\mathscr D_r\pi\rangle=\partial_r\langle\pi\rangle=0$;
the angular and axial components vanish as well. The averaged time
derivative is zero. These verifications retain the radial connection
term and physical viscosity before taking the average.

Finally, for a symmetric tensor independent of $\theta$ in cylindrical
components, differentiation of
$e_r=(\cos\theta,\sin\theta,0)$ and
$e_\theta=(-\sin\theta,\cos\theta,0)$ gives exactly the three
expressions in \eqref{mcstress:all_components}. One direct verification
is to write its Cartesian matrix as $O_\theta S O_\theta^T$, where
$O_\theta=[e_r,e_\theta,e_z]$, and use
$\partial_x=\cos\theta\partial_r-r^{-1}\sin\theta\partial_\theta$
and $\partial_y=\sin\theta\partial_r+r^{-1}\cos\theta\partial_\theta$.
The derivatives of $O_\theta$ give $-S_{\theta\theta}/r$ in the radial
row, $2S_{r\theta}/r$ in the angular row, and $S_{rz}/r$ in the axial
row. Averaging the lifted version replaces $\mathscr D_r$ by
$\partial_r$. This proves the full Cartesian identity.
\end{proof}

\subsection{The compact fluxes and the original varying radial bumps}

For $e=1,2$, write $\tau_1=z$, $\tau_2=\theta$, and set
\begin{equation}\label{mcstress:fluxes}
 M_e(z,t)=\int_0^\infty r^eF_e(r,z,t)\,dr,
 \qquad
 J_e(z,t)=\int_0^\infty r^eS_{z\tau_e}(r,z,t)\,dr.
\end{equation}
Multiplication of the corresponding row of
\eqref{mcstress:all_components} by $r^e$ makes its radial term
$\partial_r(r^e S_{r\tau_e})$. The radial boundary values vanish
by compact annular support. Therefore
\begin{equation}\label{mcstress:flux_identity}
 M_e=\partial_zJ_e,\qquad
 \int_{\mathbb R}M_e(z,t)\,dz=0\quad(e=1,2).
\end{equation}
Both $J_e$ are smooth and compactly supported in $z$. No claim that
either $M_e$ vanishes pointwise is used.

Retain the original $q$-dependent bumps, for $q(z,t)>0$ on the compact
source patch:
\begin{equation}\label{mcstress:bumps}
 b_e(r,z,t)=q(z,t)^{-(e+1)/2}
             \widehat b_e(r/\sqrt{q(z,t)}),\qquad
 \widehat b_e\in C_c^\infty((0,\infty)),\qquad
 \int_0^\infty x^e\widehat b_e(x)\,dx=1.
\end{equation}
The fixed profiles may be taken to be the already selected ones.
The change of variable $r=\sqrt q\,x$ proves
$\int r^e b_e\,dr=1$ for every $(z,t)$. Choose a fixed annulus
containing the supports of $S$ and of the bumps on the relevant compact
$(z,t)$ set. This is possible because $q$ has positive lower and finite
upper bounds there. Outside the compact axial set where a flux or
force is nonzero all the formulas below extend by zero.

The exact derivatives and cumulative bump are
\begin{equation}\label{mcstress:bump_derivatives}
\begin{gathered}
 \partial_zb_e=-\frac{q_z}{2q}
                  \bigl((e+1)b_e+r\partial_rb_e\bigr),\\
 \beta_e(r,z,t)=\int_0^r s^eb_e(s,z,t)\,ds
              =\int_0^{r/\sqrt q}x^e\widehat b_e(x)\,dx,\\
 \partial_z\beta_e=-\frac{q_z}{2q}r^{e+1}b_e(r,z,t),\qquad
 \int_0^\infty r^e\partial_zb_e\,dr=0.
\end{gathered}
\end{equation}
The first identity is the product and chain rule. The second is the
same radial substitution as the normalization. Differentiating its
upper limit proves the third. Differentiating the normalization under
the uniformly compact radial integral proves the fourth. In particular
none of these identities assigns a sign to $q_z$.

\subsection{Corrected radial primitives and the complete symmetric tensor}

Define the actual corrected primitives
\begin{equation}\label{mcstress:sigma}
 \sigma_e(r,z,t)=-r^{-e}\int_0^r s^e
       \bigl(F_e(s,z,t)-\partial_z(b_e(s,z,t)J_e(z,t))\bigr)\,ds,
 \qquad e=1,2.
\end{equation}
Their exact radial equations are
\begin{equation}\label{mcstress:sigma_equation}
 (\partial_r+e/r)\sigma_e
      =-F_e+\partial_z(b_eJ_e).
\end{equation}
This follows by multiplying \eqref{mcstress:sigma} by $r^e$ and
using the fundamental theorem of calculus.

For comparison with the initially proposed primitive, put
\[
 \sigma_e^{\rm old}=-r^{-e}\int_0^r s^e
                       (F_e-b_e\partial_zJ_e)\,ds.
\]
Every missing term can be displayed exactly:
\begin{equation}\label{mcstress:primitive_correction}
 \sigma_e=\sigma_e^{\rm old}+r^{-e}J_e\partial_z\beta_e
         =\sigma_e^{\rm old}-\frac{rq_z}{2q}b_eJ_e.
\end{equation}
Indeed expanding $\partial_z(b_eJ_e)$ in
\eqref{mcstress:sigma} proves the first equality, and
\eqref{mcstress:bump_derivatives} proves the second. With the old
primitive the angular or axial stress divergence below would instead
be $-F_e-(\partial_zb_e)J_e$. Thus the correction is necessary for the
retained varying bump. It vanishes in the special case $q_z=0$.

There is an exact expression entirely in the original tensor entries:
\begin{equation}\label{mcstress:compact_flux_correction}
\begin{gathered}
 C_e(r,z,t)=\int_0^r s^eS_{z\tau_e}(s,z,t)\,ds
                                    -\beta_e(r,z,t)J_e(z,t),\\
 \sigma_e=-S_{r\tau_e}-r^{-e}\partial_zC_e.
\end{gathered}
\end{equation}
To prove it, integrate the corresponding component of
\eqref{mcstress:all_components} only up to $r$, obtaining
$\int_0^r s^eF_e\,ds=r^eS_{r\tau_e}
+\partial_z\int_0^r s^eS_{z\tau_e}\,ds$. Substitute this identity
and $\int_0^r s^e b_e\,ds=\beta_e$ into
\eqref{mcstress:sigma}. Below both radial supports $C_e=0$;
above both supports its two summands are $J_e$ and $J_e$, so it is
again zero. This proves its compact support as well as the exact
relation of the constructed radial stress to the actual covariance
tensor.

\begin{theorem}\label{mcstress:completion}
The real symmetric axisymmetric tensor with cylindrical entries
\begin{equation}\label{mcstress:tensor}
\begin{gathered}
 T_{rr}=0,\qquad T_{r\theta}=T_{\theta r}=\sigma_2,\qquad
 T_{rz}=T_{zr}=\sigma_1,\\
 T_{z\theta}=T_{\theta z}=-b_2J_2,\qquad
 T_{zz}=-b_1J_1,\qquad
 T_{\theta\theta}=rF_r+r\partial_z\sigma_1
\end{gathered}
\end{equation}
is smooth and compactly supported in an annular cylinder and obeys
the full Cartesian identity
\begin{equation}\label{mcstress:divergence}
 \operatorname{div}T=-(F_r e_r+F_2e_\theta+F_1e_z).
\end{equation}
If the term $rF_r$ is omitted only in $T_{\theta\theta}$, the exact
divergence is instead $-(0,F_2,F_1)$ in cylindrical components.
\end{theorem}
\begin{proof}
The integral of the numerator in \eqref{mcstress:sigma} over the
entire radial half-line is
\[
 M_e-\partial_z\left(J_e\int_0^\infty s^eb_e\,ds\right)
  =\partial_zJ_e-\partial_zJ_e=0.
\]
Thus the numerator primitive is zero below its support and zero above
it. Division by $r^e$ introduces no singularity because it is zero on
a fixed neighborhood of the axis. All derivatives under its fixed
compact radial integral exist, and the normalization identity holds
smoothly in $(z,t)$, so its extension by zero is smooth. Compact axial
support follows because both $F_e$ and $J_e$ vanish outside the fixed
compact axial set. The same statements hold for all entries of
\eqref{mcstress:tensor}. Passing to the Cartesian matrix
$O_\theta T O_\theta^T$ preserves smoothness away from the axis; near
the axis it is identically zero. Hence this is a smooth compactly
supported Cartesian tensor.

Apply the full symmetric axisymmetric divergence formula already
proved in \eqref{mcstress:all_components}. Its three rows are
\[
 \begin{aligned}
 (\operatorname{div}T)_r
   &=-\frac{rF_r+r\partial_z\sigma_1}{r}
                                     +\partial_z\sigma_1=-F_r,\\
 (\operatorname{div}T)_\theta
   &=(\partial_r+2/r)\sigma_2-\partial_z(b_2J_2)=-F_2,\\
 (\operatorname{div}T)_z
   &=(\partial_r+1/r)\sigma_1-\partial_z(b_1J_1)=-F_1.
 \end{aligned}
\]
The last two identities are \eqref{mcstress:sigma_equation}, including
every derivative of $b_e$. This proves the Cartesian vector identity
and, on omitting $rF_r$, its stated radial-zero variant.
\end{proof}

\subsection{Compact pressure, tracefree stress, and the exact gauge relation}

The trace and its prescribed pressure are explicitly
\begin{equation}\label{mcstress:pressure}
 \operatorname{tr}T=rF_r+r\partial_z\sigma_1-b_1J_1,
 \qquad
 p_T=\frac13\bigl(rF_r+r\partial_z\sigma_1-b_1J_1\bigr),
 \qquad T^0=T-p_TI_3.
\end{equation}
Then $\operatorname{tr}T^0=0$, and every field is smooth and compactly
supported in the same annular cylinder. The exact full cancellation is
\begin{equation}\label{mcstress:tracefree_cancellation}
 F+\operatorname{div}T^0+\nabla p_T=0.
\end{equation}
To check its sign, for any scalar $p$ Cartesian differentiation gives
$\operatorname{div}(pI_3)=\nabla p$. Therefore
$\operatorname{div}T^0=\operatorname{div}T-\nabla p_T
=-F-\nabla p_T$, which proves the displayed equation.
In components the new diagonal entries are
\[
 T^0_{rr}=-p_T,\quad
 T^0_{\theta\theta}=rF_r+r\partial_z\sigma_1-p_T,\quad
 T^0_{zz}=-b_1J_1-p_T.
\]
The axial stress divergence alone changes by $-\partial_zp_T$;
the pressure gradient contributes $+\partial_zp_T$. Thus the combined
stress--pressure correction preserves the exact tangential equations.
There is no claim that their stress term alone stays unchanged.

More generally $T-pI_3$ with any smooth compact scalar $p$ satisfies
$F+\operatorname{div}(T-pI_3)+\nabla p=0$.
The choice in \eqref{mcstress:pressure} is the unique one making
this stress tracefree. The map from symmetric tensors to
tracefree-tensor--scalar pairs,
\[
 T\longmapsto\bigl(T-\tfrac13\operatorname{tr}T\,I_3,
                         \tfrac13\operatorname{tr}T\bigr),
\]
has exact inverse $(U,p)\mapsto U+pI_3$ when $\operatorname{tr}U=0$.
It intertwines $\operatorname{div}$ with
$(U,p)\mapsto\operatorname{div}U+\nabla p$. This supplies the full
map between the two presentations, including compact support and signs.

\subsection{All total force and torque components and the divergence image}

The Cartesian components of the averaged force are
\[
 F=(F_r\cos\theta-F_2\sin\theta,
       F_r\sin\theta+F_2\cos\theta,F_1).
\]
Integration in $\theta$ with the physical volume element
$dx=r\,dr\,d\theta\,dz$ gives
\begin{equation}\label{mcstress:force_torque}
\begin{gathered}
 \int_{\mathbb R^3}F\,dx
  =2\pi e_z\int_{\mathbb R}M_1(z,t)\,dz=0,\\
 \int_{\mathbb R^3}x\times F\,dx
  =2\pi e_z\int_{\mathbb R}M_2(z,t)\,dz=0.
\end{gathered}
\end{equation}
For the torque calculation, direct multiplication in the oriented
basis gives
\[
 x\times F=-zF_2e_r+(zF_r-rF_1)e_\theta+rF_2e_z.
\]
The first two terms integrate to zero in $\theta$, and the third
gives the second equation of \eqref{mcstress:force_torque}.
Thus all three force and all three torque components vanish. The last
equalities use the already proved compact flux identities
\eqref{mcstress:flux_identity}. The same torque vanishes about every
fixed origin $x_0$, because subtracting $x_0$ changes it by
$-x_0\times\int F\,dx=0$. On a finite axial slab $a<z<b$, the
axial force and torque are respectively
$2\pi(J_1(b,t)-J_1(a,t))$ and
$2\pi(J_2(b,t)-J_2(a,t))$. The quantities need not vanish on an
individual slice.

There is also a direct Cartesian stress verification. Since
$F_j=-\partial_kT_{jk}$ and $T$ is compact, integration by parts gives
\begin{equation}\label{mcstress:first_moments}
 \int F_j\,dx=0,\qquad
 \int x_iF_j\,dx=\int T_{ji}\,dx=-\int S_{ji}\,dx.
\end{equation}
Symmetry makes the antisymmetric part of the first-moment matrix zero,
which is precisely all three torque components. Its symmetric part
is not asserted to vanish. The second equality with $S$ uses
$F=\operatorname{div}S$. Equivalently
$\operatorname{div}(T+S)=0$; multiplying this equation by $x_i$ and
integrating proves $\int(T+S)_{ji}\,dx=0$. In the tracefree-pressure
presentation the same first moment is
$\int T^0_{ji}\,dx+\delta_{ij}\int p_T\,dx$.
All these identities persist under any finite time differentiation on
a common compact support, by differentiation under the integral.

For completeness, this construction proves the exact image of compact
symmetric-stress divergence in the axisymmetric class. Let
an open axial interval $I_z$ contain the actual axial supports, with
$q>0$ smooth on a neighborhood of its relevant compact subsets at the
fixed time. In the following two spaces all supports are compact
subsets of $\{r>0,\ z\in I_z\}$. Thus the original bump is defined
everywhere it is used, and integrating between axial support endpoints
stays inside $I_z$.
Let
$\mathcal F_0$ be the space of smooth real axisymmetric compact forces
supported away from $r=0$ and satisfying exactly
\[
 \int_{\mathbb R}\int_0^\infty rF_1\,dr\,dz=0,
 \qquad
 \int_{\mathbb R}\int_0^\infty r^2F_2\,dr\,dz=0.
\]
Let $\mathcal S_c$ be the space of smooth real compact symmetric
axisymmetric Cartesian tensors supported away from $r=0$. Define
$\mathcal D:\mathcal S_c\to\mathcal F_0$ by
$\mathcal D(U)=-\operatorname{div}U$.
Integration by parts proves that its values satisfy both defining
identities of $\mathcal F_0$, and axisymmetry makes the other force
and torque components vanish as above.

For any $F\in\mathcal F_0$, compute its actual radial moments and set
\[
 J_e(z)=\int_{-\infty}^z\int_0^\infty r^eF_e(r,s)\,dr\,ds.
\]
This is a smooth compact axial function: its derivative is $M_e$, it
vanishes below the axial support, and the defining zero total integral
makes it vanish above that support. Use the fixed bumps
\eqref{mcstress:bumps}, the corrected primitive
\eqref{mcstress:sigma}, and all entries of
\eqref{mcstress:tensor}. The proof of Theorem~\ref{mcstress:completion}
uses only these identities and proves a compact tensor
$\mathcal C_bF$ with
\begin{equation}\label{mcstress:right_inverse}
 \mathcal D\mathcal C_bF=F\qquad(F\in\mathcal F_0).
\end{equation}
Thus the two total moments are necessary and sufficient, with an
explicit right inverse. For the actual force in
\eqref{mcstress:Ffull}, this primitive in $z$ equals the already
computed $J_e$ from $S$, because both have derivative $M_e$ and vanish
at the negative axial end. No additional compatibility condition has
been imposed on the actual candidate.

For fixed $q$ and bump profiles, all operations defining
$\mathcal C_b$ are linear. Hence
$\mathcal P_b=\mathcal C_b\mathcal D$ is a projection on
$\mathcal S_c$: \eqref{mcstress:right_inverse} gives
$\mathcal P_b^2=\mathcal P_b$ and
$\ker\mathcal P_b=\ker\mathcal D$.
In particular the completed actual tensor is
$T=\mathcal P_b(-S)$, and $T+S$ is a compact divergence-free symmetric
tensor. This gives an exact relation to the actual averaged covariance,
not merely equality of its two radial moments.

\subsection{Exact source chart, bump, stress, and pressure powers}

Fix one source chart and retain all of its constants:
\begin{equation}\label{mcstress:chart}
 r=\sqrt Q\,R,\qquad z=Q^D Z,\qquad1-t=QT,
 \qquad A=\tfrac12+h,\quad D=\tfrac12-h,\quad
 \varepsilon=Q^h.
\end{equation}
The dyadic label $Q$ is fixed under differentiation. The averaged quantities
are independent of $Y$ and $\theta$, so in this
subsection the complete normalized spatial derivatives are
$\partial_R$, $R^{-1}\partial_\theta$ and $\varepsilon\partial_Z$,
with the same cylindrical basis derivatives as before.
Define the exact chart representatives
\begin{equation}\label{mcstress:chart_representatives}
\begin{gathered}
 S^*=Q^{2A}S,\qquad F^*=Q^{2A+1/2}F,\qquad
 T^*=Q^{2A}T,\qquad \Sigma_e=Q^{2A}\sigma_e,\\
 J_e^*=Q^{2A-(e+1)/2}J_e
             =\int_0^\infty R^e S^*_{z\tau_e}\,dR,\qquad
 M_e^*=Q^{2A-e/2}M_e=\varepsilon\partial_ZJ_e^*,\\
 q_*=q/Q,\qquad
 b_e^*=Q^{(e+1)/2}b_e(\sqrt Q R,z,t)
       =q_*^{-(e+1)/2}\widehat b_e(R/\sqrt{q_*}).
\end{gathered}
\end{equation}
The two moment powers follow directly from
$r^e\,dr=Q^{(e+1)/2}R^e\,dR$ and
$\partial_z=Q^{-D}\partial_Z$; the remaining factor in $M_e^*$
is $Q^{1/2-D}=\varepsilon$. In particular the scaled bump satisfies
$\int R^e b_e^*\,dR=1$.

Every term of the corrected primitive has the same exact stress unit:
\begin{equation}\label{mcstress:chart_sigma}
\begin{gathered}
 \Sigma_e=-R^{-e}\int_0^R X^e
             \bigl(F_e^*-\varepsilon\partial_Z(b_e^*J_e^*)\bigr)\,dX,\\
 (\partial_R+e/R)\Sigma_e
       =-F_e^*+\varepsilon\partial_Z(b_e^*J_e^*),\\
 \Sigma_e=\Sigma_e^{\rm old}
       -\frac{\varepsilon R\partial_Zq_*}{2q_*}b_e^*J_e^*.
\end{gathered}
\end{equation}
To verify the first line, insert
$F_e=Q^{-2A-1/2}F_e^*$ and $b_eJ_e=Q^{-2A}b_e^*J_e^*$ in
\eqref{mcstress:sigma}, and change the integration variable to
$s=\sqrt Q X$. The second follows by differentiation. For the third,
the physical correction in \eqref{mcstress:primitive_correction}
has factor $\sqrt Q\,q_z/q$; using
$q_z=Q^{1-D}\partial_Zq_*$ gives
$\sqrt Q\,q_z/q=\varepsilon\partial_Zq_*/q_*$.
This verifies the original $q$ derivative, including its sign.
Also, with $C_e^*=Q^{2A-(e+1)/2}C_e$,
\[
 \Sigma_e=-S^*_{r\tau_e}
                     -\varepsilon R^{-e}\partial_Z C_e^*.
\]

The completed stress and its pressure become exactly
\begin{equation}\label{mcstress:chart_tensor}
\begin{gathered}
 T^*_{rr}=0,\quad T^*_{r\theta}=\Sigma_2,\quad
 T^*_{rz}=\Sigma_1,\quad T^*_{z\theta}=-b_2^*J_2^*,\quad
 T^*_{zz}=-b_1^*J_1^*,\\
 T^*_{\theta\theta}=RF_r^*+\varepsilon R\partial_Z\Sigma_1,
 \qquad
 p_T^*=Q^{2A}p_T
  =\tfrac13(RF_r^*+\varepsilon R\partial_Z\Sigma_1-b_1^*J_1^*),\\
 (T^0)^*=T^*-p_T^*I_3,\qquad
 F^*+\operatorname{div}_*(T^0)^*+\nabla_*p_T^*=0.
\end{gathered}
\end{equation}
For example the radial row is
$-T^*_{\theta\theta}/R+\varepsilon\partial_Z\Sigma_1=-F_r^*$,
and the tangential rows use \eqref{mcstress:chart_sigma}. The pressure
identity follows from the same Cartesian isotropic-tensor calculation
as before. Physical divergence and pressure gradient both have factor
$Q^{-2A-1/2}$ because $\nabla_x=Q^{-1/2}\nabla_*$ on the chart.
Thus the radial residual, hoop stress, both surviving axial fluxes,
varying-bump correction and compact pressure all have their exact
source units. No decay or uniform estimate at a singular endpoint is
inferred from these finite compact-support identities.


\section{Original-scale derivative costs of the actual compatibility defects}
\label{mcq:section}

The state and increment in this section are exactly the finite source state
and actual third-growth curl constructed in the preceding reader.
Write \(u_*=\langle u_*\rangle_\theta+w_*\), \(v=w_y\), and
\[
 C=\langle w_*\otimes v+v\otimes w_*+v\otimes v\rangle,
\]
where the outer brackets are the normalized angular and auxiliary
averages before physical evaluation. Every component below is in the
original oriented cylindrical orthonormal basis. The full moment
completion and source pressure calculation prove the formulas used here;
we also derive the scalar compatibility formulas directly to fix all signs.
Throughout this section the unadorned symbols \(P,J_\theta,J_z\)
denote the defect increments \(\Delta P,\Delta J_\theta,\Delta J_z\)
caused by this actual wave addition. They are not the full old-plus-new
source defects. The unchanged old defects must be added when applying
the source correction to the whole updated state.

\subsection{Actual mixed derivatives of the third-growth curl}

Retain the fixed finite set of source labels \(j\), the actual clocks
\(t_{c,j}=t_b+\alpha_jv_j\), and the two full shape fields \(W_{ja}\).
The original data \(t_b,r_a,S_3,u_i\) are independent of these spatial
labels. In particular
\[
 v=\sum_{j,a}Q_j^{-A}y_a(t_{c,j})W_{ja},\quad
 A=\tfrac12+h,\quad D=\tfrac12-h,\quad
 \partial_t t_{c,j}=\alpha_jQ_j^{-1-h},\quad \partial_z t_{c,j}=0 .
\]
Let \(Y_{1,n}=T_n\) and \(Y_{2,n}=O_n\) be the complete finite
derivative arrays already proved for the actual coupled solution.
Thus \(Y_{1,0}=(A_0/\mu)\widehat\lambda_3^{-7/8}\),
\(Y_{2,0}=2A_0\widehat\lambda_3^{1/8}/(\sigma_2\sqrt{n_i})\),
\(Y_{1,1}=2\gamma_iY_{1,0}\), and
\(Y_{2,1}=\gamma_iA_0\widehat\lambda_3^{1/8}/(\sigma_2\sqrt{n_i})\).
Higher entries are those closed recurrences in the retained moving-parent
data; no independent amplitude bound is introduced here.

For a pair \(I=(a,b)\in\mathbb N^2\), put
\(\partial^I=\partial_z^a\partial_t^b\) and
\(I_z=I+(1,0)\). On a fixed chart the full time operator is
\(\mathfrak t_{*,j}=-\varepsilon_j\partial_{T_j}+c_{i,j}N_{i,j}\).
The exact mixed derivative formula is
\begin{equation}\label{mcq:actual_derivative}
\begin{split}
 \partial_z^a\partial_t^b v
 =\sum_{j,c}\sum_{s=0}^b&
 \binom bs Q_j^{-A-aD-b(1+h)}\alpha_j^s
 y_c^{(s)}(t_{c,j})\\
 &\hspace{4mm}\cdot
 \partial_{Z_j}^a\mathfrak t_{*,j}^{\,b-s}W_{jc}.
\end{split}
\end{equation}
All operators on the shapes include their phase, frame and cutoff
derivatives. The shape time operator is the full derivative after
auxiliary evaluation, not a derivative that holds the fast phase fixed.
To prove the identity, differentiate the physical product \(b\) times
in time. If \(s\) of the derivatives act on \(y_c\), its factor is
\((\alpha_jQ_j^{-1-h})^sy_c^{(s)}\); the other \(b-s\) derivatives
give \(Q_j^{-(b-s)(1+h)}\mathfrak t_{*,j}^{b-s}W_{jc}\).
There are \(\binom bs\) such choices. The \(a\) axial derivatives act
only on the shape and give \(Q_j^{-aD}\partial_{Z_j}^a\).
The axial derivative commutes with the stated time operator, since
its coefficients are fixed label constants. Combining the powers proves
\eqref{mcq:actual_derivative} without dropping any derivative.

For each \(e\in\{-1,0,1,2\}\), define the precise weighted norm
\[
 \|f\|_{[e]}^2=\int_0^\infty r^e\langle |f|^2\rangle\,dr,\qquad
 \|F\|_{[e],j}^2=\int_0^\infty R_j^e\langle |F|^2\rangle\,dR_j.
\]
Every field being integrated is supported in a compact annulus away
from zero, so the negative weight introduces no singular integral.
Define the explicit actual-shape budget
\begin{equation}\label{mcq:V}
\begin{split}
 \mathcal V_{e,(a,b)}
 =\sum_{j,c}\sum_{s=0}^b&
 \binom bs \alpha_j^sY_{c,s}
 Q_j^{(e+1)/4-A-aD-b(1+h)}\\
 &\hspace{4mm}\cdot
 \|\partial_{Z_j}^a\mathfrak t_{*,j}^{\,b-s}W_{jc}\|_{[e],j}.
\end{split}
\end{equation}
Then \(\|\partial^I v\|_{[e]}\le\mathcal V_{e,I}\).
Indeed the measure identity
\(r^e\,dr=Q_j^{(e+1)/2}R_j^e\,dR_j\), valid also for \(e=-1\),
and the triangle inequality applied to \eqref{mcq:actual_derivative}
give this bound. There is no inter-label orthogonality assumption.
Put \(\mathcal W_{e,I}=\|\partial^I w_*\|_{[e]}\), restricted to a
fixed radial set enclosing the supports of the actual new fields
and all derivatives at issue. The old wave includes its existing curl
corrections. These are norms of the specified source state, not missing
amplitude variables.

\subsection{A contraction estimate that retains every product derivative}

For \(J=(j_1,j_2)\le I=(a,b)\), set
\(\binom IJ=\binom a{j_1}\binom b{j_2}\), and define
\begin{equation}\label{mcq:B}
\begin{split}
 \mathcal B_{e,I}=\sum_{J\le I}\binom IJ
 \bigl(&\mathcal W_{e,J}\mathcal V_{e,I-J}
       +\mathcal V_{e,J}\mathcal W_{e,I-J}\\
       &+\mathcal V_{e,J}\mathcal V_{e,I-J}\bigr).
\end{split}
\end{equation}
If \(E\) is any fixed real symmetric three-by-three matrix in
cylindrical components, then
\begin{equation}\label{mcq:contraction}
 \left|\partial^I\int r^e(E:C)\,dr\right|
       \le \|E\|_{\mathrm{op}}\mathcal B_{e,I}.
\end{equation}
For a single outer product,
\(E:(f\otimes g)=f^TEg\), whose absolute value is at most
\(\|E\|_{\mathrm{op}}|f||g|\). The repeated product rule gives exactly
the sum over \(J\le I\) in \eqref{mcq:B}. Cauchy--Schwarz on the
original weighted angular-auxiliary-radial measure gives each displayed
product of norms. Compact smooth support justifies differentiation
under the integral. Axial and time derivatives do not differentiate
the cylindrical basis. This proves \eqref{mcq:contraction} for every
specified order.
For the time derivative, the full lifted physical operator is
\(\mathcal D_t=\partial_t+N_{\rm abs}\), with \(N_{\rm abs}\) a
constant torus direction. The Haar integral of \(N_{\rm abs}f\) is
zero by periodicity. Therefore
\(\partial_t\langle f\rangle=\langle\mathcal D_tf\rangle\).
The same identity applies repeatedly, and axial differentiation commutes
with this average and this constant torus direction. Consequently the
ordinary mixed derivatives of the averaged tensor in
\eqref{mcq:contraction} are exactly the averages of the full physical
derivatives estimated in \eqref{mcq:actual_derivative}.

\subsection{The actual three defects and the pressure-corrected axial flux}

Since the wave increment has zero angular mean and is divergence-free,
the change of the radial mean pressure source is
\[
 g=-(\partial_r+r^{-1})C_{rr}-\partial_zC_{zr}+r^{-1}C_{\theta\theta}.
\]
The three actual source compatibility-defect increments are therefore
\begin{equation}\label{mcq:defects}
\begin{split}
 P&=\int\frac{C_{\theta\theta}-C_{rr}}r\,dr
                    -\partial_z\int C_{zr}\,dr,\\
 J_\theta&=\int r^2C_{z\theta}\,dr,\\
 J_z&=\int r\left(C_{zz}-\tfrac12C_{rr}
                              -\tfrac12C_{\theta\theta}\right)\,dr
                  +\tfrac12\partial_z\int r^2C_{zr}\,dr .
\end{split}
\end{equation}
The first formula is \(\int g\,dr\), using the zero boundary integral
of \(\partial_r C_{rr}\). For the last, integration by parts gives
\[
 \int r^2 g\,dr
  =\int r(C_{rr}+C_{\theta\theta})\,dr
                      -\partial_z\int r^2C_{zr}\,dr .
\]
Substitute it into the source definition
\(J_z=\int rC_{zz}\,dr-\frac12\int r^2g\,dr\).
This proves all three identities in \eqref{mcq:defects}.
In particular the axial source defect is not replaced by the raw flux
\(\int rC_{zz}\,dr\).

Let the source pressure bump be the original
\[
 \rho(r,z,t)=q^{-1/2}\widehat\rho(r/\sqrt q),\qquad
 \int\widehat\rho(x)\,dx=1,\qquad
 c_0=\tfrac12\int x^2\widehat\rho(x)\,dx .
\]
It has \(\int\rho\,dr=1\) and
\(c_\rho=\frac12\int r^2\rho\,dr=c_0q\).
The auxiliary average of the compact source pressure increment is
\[
 \overline{\Delta p_m}(r)=\int_0^r(g(s)-\rho(s)P)\,ds .
\]
Its derivative is \(g-\rho P\), and its integrand has zero total
radial integral, giving compact support. Its weighted pressure moment is
\(\int r\overline{\Delta p_m}\,dr=-\frac12\int r^2g\,dr+c_\rho P\).
Here is the exact map to the full auxiliary-dependent source operation.
Let \(\Delta g_r\) be the angularly averaged source increment before
auxiliary averaging, so that \(\langle\Delta g_r\rangle_Y=g\), and
put \(f=\Delta g_r-\rho P\). The source's shifted primitive uses
\[
 \begin{split}
 \mathcal I f(r,Y)&=\int_0^r
       f(s,Y+(s^{d_r}-r^{d_r})v_r)\,ds,\\
 \mathcal J f(r,Y)&=\int_0^\infty
       f(s,Y+(s^{d_r}-r^{d_r})v_r)\,ds,\qquad
 \mathcal I_cf=\mathcal If-\chi_m\mathcal Jf .
 \end{split}
\]
The fixed source exponent and direction are the original \(d_r,v_r\),
and \(\chi_m\) is its inner-zero, outer-one radial cutoff. The full
pressure increment is \(\Delta p_m=\mathcal I_cf\).
Haar translation invariance gives
\(\langle\mathcal Jf\rangle_Y=\int(g-\rho P)\,dr=0\) and
\(\langle\mathcal If\rangle_Y=\int_0^r(g-\rho P)\,ds\).
This proves the asserted averaged pressure formula from the full map.
The lifted radial derivative cancels the shift derivative and yields
\[
 \mathcal D_r\Delta p_m-\Delta g_r
       =-\rho P-(\partial_r\chi_m)\mathcal Jf .
\]
The second term has zero auxiliary average but need not vanish before
physical evaluation. No estimate for its singular-time flatness is
inferred from its zero average.
Consequently the double-averaged axial residual-moment increment after
this reconstruction is
\[
 \partial_zK_z,\qquad K_z=J_z+c_0qP .
\]
The full product, including the derivatives of \(q\), is retained.

The following simultaneous quantitative estimates hold at every finite
mixed order:
\begin{equation}\label{mcq:defect_bounds}
\begin{split}
 |\partial^IP|&\le
       \mathcal B_{-1,I}+\tfrac12\mathcal B_{0,I_z},\\
 |\partial^IJ_\theta|&\le\tfrac12\mathcal B_{2,I},\\
 |\partial^IJ_z|&\le
       \mathcal B_{1,I}+\tfrac14\mathcal B_{2,I_z},\\
 |\partial^IK_z|&\le
       \mathcal B_{1,I}+\tfrac14\mathcal B_{2,I_z}\\
 &\quad+|c_0|\sum_{J\le I}\binom IJ
       |\partial^Jq|\left(
       \mathcal B_{-1,I-J}+\tfrac12\mathcal B_{0,(I-J)_z}\right).
\end{split}
\end{equation}
For the first line the diagonal contraction matrix is
\(\operatorname{diag}(-1,1,0)\), of norm one, whereas the \(rz\)
entry is the contraction against the symmetric matrix with \(rz,zr\)
entries \(1/2\), of norm \(1/2\). For the second line use the
analogous \(z\theta\) matrix. For the third the diagonal matrix is
\(\operatorname{diag}(-1/2,-1/2,1)\), of norm one, and the
coefficient \(1/2\) multiplying \(C_{zr}\) contributes an additional
factor \(1/2\). These facts and \eqref{mcq:contraction} prove the first
three bounds. The exact product rule for \(qP\) proves the fourth.
Replacing \(I\) by \(I_z\) gives bounds for the two residual moments
that actually enter the source radial inverse. Thus every pressure,
radial and axial coupling has a bound in the original shape fields
and third-growth constants.

\subsection{Exact derivatives of the original moving scale and bumps}

The source scale is determined by
\[
 q-z^2q^{2h}=1-t,\qquad z=q^D\eta,\qquad
 D=\tfrac12-h,\qquad 0<h<1/100,\quad -1<\eta<1 .
\]
Set \(B(\eta)=1-2h\eta^2>1-2h>0\). Implicit differentiation of the
first equation, followed by the displayed formula for \(z\), gives
\begin{equation}\label{mcq:q_derivatives}
\begin{gathered}
 q_z=\frac{2\eta q^A}{B},\qquad q_t=-\frac1B,\\
 \eta_z=\frac{q^{-D}(1-\eta^2)}B,\qquad
 \eta_t=\frac{D\eta q^{-1}}B .
\end{gathered}
\end{equation}
For example \(q_z=2zq^{2h}/B\), and \(D+2h=A\).
Differentiating \(\eta=zq^{-D}\) gives the last two formulas;
the identity \(2D+2h=1\) gives the factor \(1-\eta^2\).
This proves the derivative dictionary without assigning a sign to \(q_z\).

Put \(x=r/\sqrt q\). For a smooth function \(f(x,\eta)\) define the
explicit differential expressions
\[
 \mathcal Z_\sigma f=
 \frac{2\sigma\eta f-\eta x f_x+(1-\eta^2)f_\eta}{B},\qquad
 \mathcal T_\sigma f=
 \frac{-\sigma f+\frac12x f_x+D\eta f_\eta}{B}.
\]
Equations \eqref{mcq:q_derivatives} and the ordinary chain rule prove
\begin{equation}\label{mcq:dilation_dictionary}
 \partial_z(q^\sigma f)=q^{\sigma-D}\mathcal Z_\sigma f,\quad
 \partial_t(q^\sigma f)=q^{\sigma-1}\mathcal T_\sigma f,\quad
 \partial_r(q^\sigma f)=q^{\sigma-1/2}f_x .
\end{equation}
For the last formula \(q,\eta\) are independent of \(r\).
For the first two, differentiation of \(x\) gives respectively
\(-\eta xq^{-D}/B\) and \(xq^{-1}/(2B)\), proving every term.

In particular let the original radial moment bump be
\[
 b_e=q^{-(e+1)/2}\widehat b_e(x),\qquad
 \int x^e\widehat b_e(x)\,dx=1 .
\]
For any \(k,a,b\ge0\), its derivative is exactly
\begin{equation}\label{mcq:bump_derivatives}
 \partial_r^k\partial_z^a\partial_t^b b_e
 =q^{-(e+1+k)/2-aD-b} f_{e;k,a,b}(x,\eta).
\end{equation}
Here the coefficient is given by a finite recursion: start at
\(\sigma=-(e+1+k)/2\) and \(f=\widehat b_e^{(k)}\);
repeat \(f\leftarrow\mathcal Z_\sigma f\),
\(\sigma\leftarrow\sigma-D\), \(a\) times, and then
\(f\leftarrow\mathcal T_\sigma f\),
\(\sigma\leftarrow\sigma-1\), \(b\) times.
The result is \(f_{e;k,a,b}\). This proves the identity by induction
using \eqref{mcq:dilation_dictionary}. Smooth mixed derivatives commute,
so this fixed order computes the stated derivative.
Every coefficient is explicitly a finite rational expression in \(x,\eta\),
the retained constants, and derivatives of the specified bump, with
nonzero denominator a power of \(B\). On its compact \(x\) support
and \(|\eta|\le1\) its maximum is finite and computable directly from
that expression. Thus \eqref{mcq:bump_derivatives} supplies the original
power of \(q\) at every order.
The pressure bump \(\rho\) uses exactly this recursion with \(e=0\).
Likewise every derivative in \eqref{mcq:defect_bounds} is explicit:
\(\partial_z^a\partial_t^bq=q^{1-aD-b}Q_{a,b}(\eta)\), where the same
recursion starts with \(\sigma=1,f=1\).

For completeness, the weighted cumulative bump has the exact form
\[
 \beta_e(r,z,t)=\int_0^rs^eb_e(s,z,t)\,ds
              =\int_0^{r/\sqrt q}x^e\widehat b_e(x)\,dx .
\]
Its axial derivative is
\begin{equation}\label{mcq:bump_movement}
 \partial_z\beta_e
 =-\frac{q_z}{2q}\,r^{e+1}b_e .
\end{equation}
This follows by differentiating its upper endpoint, which is
\(-rq_z/(2q^{3/2})\). Inserting the definition of \(b_e\) gives
the displayed original physical factors. Therefore replacing the
primitive of \(F_e-b_e\partial_zJ_e\) by the exact primitive of
\(F_e-\partial_z(b_eJ_e)\) changes the stress coefficient by
\[
 -\frac{rq_z}{2q}\,b_eJ_e .
\]
This is the moving-bump term required in the complete stress construction.
Its sign and all further derivatives are fixed by
\eqref{mcq:q_derivatives}--\eqref{mcq:bump_movement}; it is not set to zero
by holding the physical source scale fixed during axial differentiation.

These calculations supply finite quantitative inputs for the actual
source pressure and moment operations. They do not identify the
distinct direct candidate with a completed infinite source correction
sequence: the complete source-operation formulas retain the new
velocity, covariance, pressure and nonlinear defects that such a step
actually creates.


\section{Physical derivative and full vector-force costs of the five-row mean operation}
\label{mcmean:section}

This section completes the finite mean operation on the actual current
source state. All coordinates are the original right-handed cylindrical
coordinates. The physical viscosity is the fixed number
\(\nu_{\rm NS}>0\). Write the full current velocity as
\(U=(b+\beta,V+v,G+\gamma)+w\), where \(w\) includes every current
wave and curl correction. During the operation below this wave is held
fixed. The three targets are the full current defects
\begin{equation}\label{mcmean:full_targets}
 \mathcal P=P_{\rm old}+\Delta P,\qquad
 \mathcal J_\theta=(J_\theta)_{\rm old}+\Delta J_\theta,\qquad
 \mathcal J_z=(J_z)_{\rm old}+\Delta J_z.
\end{equation}
Here the increments are the pressure-corrected actual wave increments
of \eqref{mcq:defects}. If other finite source operations have already
been performed, ``old'' means their actual entering state and the
increments are recomputed for the actual addition. Thus none of the
following formulas substitutes an increment for a full target.

Retain, without a change of scale or profile constant,
\begin{equation}\label{mcmean:original_scale}
\begin{gathered}
 q-z^2q^{2h}=1-t,\quad z=q^D\eta,\quad x=r/\sqrt q,\quad
 A=\tfrac12+h,\quad D=\tfrac12-h,\\
 0<h<1/100,\quad q>0,\quad -1<\eta<1,\quad
 B(\eta)=1-2h\eta^2,\\
 V=q^{-A}a(\eta)x^{-1-2\lambda},\quad G=0
 \quad\hbox{on the reserved mean patch},\\
 a(\eta)=\frac{2^{1/2+\lambda}c_{\rm patch}}{1+\eta^2},
 \qquad \lambda>0,\qquad |a(\eta)|\ge a_0>0.
\end{gathered}
\end{equation}
The symbol \(a(\eta)\) in this display is a scalar base coefficient;
below the bold symbol \(\boldsymbol a\) denotes the vector increment.
The original constant \(c_{\rm patch}\) can have either nonzero sign.

\subsection{Exact inverse profiles and the five physical rows}

Use precisely the bumps and spacing of the five-row operation:
\[
 \varphi_j(x)=e^{-jd_b}\varphi_0(e^{-jd_b}x),\qquad
 \mu_p=\int_0^\infty x^p\varphi_0(x)\,dx>0,
 \qquad t_p=e^{d_bp},\quad d_b>0.
\]
The nonnegative nonzero smooth bump \(\varphi_0\) and its three
dilates have support in the reserved open interval \(I_m\Subset(0,\infty)\).
The substitution \(x=e^{jd_b}\xi\) proves exactly
\(\int x^p\varphi_j(x)\,dx=\mu_p t_p^j\).
Number matrix rows and columns starting at zero, and define
\begin{equation}\label{mcmean:vandermonde}
\begin{gathered}
 V_\theta=[t_p^j]_{p=2,-2-2\lambda,-2\lambda; j=0,1,2},
 \qquad B_\theta=V_\theta^{-1},\\
 V_z=[t_p^j]_{p=1,1-2\lambda; j=0,1},
 \qquad B_z=V_z^{-1},\\
 A_\theta=\operatorname{diag}
  (\mu_2,2a\mu_{-2-2\lambda},-a\mu_{-2\lambda})V_\theta,
 \qquad
 A_z=\operatorname{diag}(\mu_1,a\mu_{1-2\lambda})V_z.
\end{gathered}
\end{equation}
For three row values \(t_0,t_1,t_2\), subtraction of the first row
from each later row, followed by factoring, gives
\(\det V_\theta=(t_1-t_0)(t_2-t_0)(t_2-t_1)\).
For two row values the determinant is \(t_1-t_0\).
The original exponents in each displayed matrix are distinct because
\(\lambda>0\), and the exponential with \(d_b>0\) is injective.
The row factors are nonzero. This proves both inverses exist at the
actual fixed source parameters. Explicitly each entry of either
\(B\) is its corresponding transposed cofactor divided by the stated
determinant; this is a finite formula in the original exponentials.

Define the following actual profiles, including their \(\eta\) dependence:
\begin{equation}\label{mcmean:profiles}
\begin{split}
 V_P(x,\eta)&=-\frac{1}{2a(\eta)\mu_{-2-2\lambda}}
                   \sum_{j=0}^2B_\theta[j,1]\varphi_j(x),\\
 V_{J_z}(x,\eta)&=\frac{1}{a(\eta)\mu_{-2\lambda}}
                   \sum_{j=0}^2B_\theta[j,2]\varphi_j(x),\\
 G_{J_\theta}(x,\eta)&=-\frac{1}{a(\eta)\mu_{1-2\lambda}}
                   \sum_{j=0}^1B_z[j,1]\varphi_j(x).
\end{split}
\end{equation}
The signs, the factor two, and all three negative powers in the moment
denominators are part of the formulas. Multiplying the columns by
the exact rows in \eqref{mcmean:vandermonde} proves the eight scalar
profile identities
\begin{equation}\label{mcmean:profile_moments}
\begin{array}{c|cc}
 &V_P&V_{J_z}\\ \hline
 \int x^2(\cdot)\,dx&0&0\\
 2a\int x^{-2-2\lambda}(\cdot)\,dx&-1&0\\
 -a\int x^{-2\lambda}(\cdot)\,dx&0&-1
\end{array}
\qquad
\begin{array}{c|c}
 &G_{J_\theta}\\ \hline
 \int x(\cdot)\,dx&0\\
 a\int x^{1-2\lambda}(\cdot)\,dx&-1
\end{array}.
\end{equation}
There are eight entries in these two tables; each follows from
\(V_\theta B_\theta=I_3\) or \(V_zB_z=I_2\), with the row factor
left in place. The complete matrix identities include all nine and
all four scalar inverse identities, respectively.

The source coefficients are
\[
 \boldsymbol u=A_\theta^{-1}
       (0,-q^{2A}\mathcal P,-q^{2A-1}\mathcal J_z)^T,
 \qquad
 \boldsymbol s=A_z^{-1}
       (0,-q^{2A-3/2}\mathcal J_\theta)^T.
\]
Since \((DV)^{-1}=V^{-1}D^{-1}\), substitution of these coefficients
in the original physical increments gives
\begin{equation}\label{mcmean:tangential_increment}
 \delta v=q^A\mathcal P V_P+q^{A-1}\mathcal J_zV_{J_z},
 \qquad
 \delta\gamma=q^{A-3/2}\mathcal J_\theta G_{J_\theta}.
\end{equation}
For example the fifth-row coefficient in \(\boldsymbol u\) is
\((-q^{2A-1}\mathcal J_z)/(-a\mu_{-2\lambda})\), which proves
the positive sign in \(V_{J_z}\). Multiplication by the original
outer factor \(q^{-A}\) proves its exponent \(A-1\).
The other two terms follow by the same displayed diagonal products.

Changing variables \(r=\sqrt q\,x\) in each integral, and using
\eqref{mcmean:profile_moments}, now proves the full five equations
\begin{equation}\label{mcmean:five_rows}
\begin{gathered}
 \int r^2\delta v\,dr=0,\qquad
 \int r\delta\gamma\,dr=0,\qquad
 \int\frac{2V}{r}\delta v\,dr=-\mathcal P,\\
 \int r^2(G\delta v+V\delta\gamma)\,dr=-\mathcal J_\theta,\\
 \int(2rG\delta\gamma-rV\delta v)\,dr=-\mathcal J_z.
\end{gathered}
\end{equation}
Indeed the pressure term with \(V_P\) has scale
\(q^{-A}q^A=1\); the angular flux with \(G_{J_\theta}\) has scale
\(q^{3/2}q^{-A}q^{A-3/2}=1\); and the fifth row with
\(V_{J_z}\) has scale \(q q^{-A}q^{A-1}=1\).
Every other term is one of the zero entries in the tables. The two
homogeneous velocity moments use the first rows of those tables.
This verifies the physical equations, including the single factor
\(r\) multiplying \(V\delta v\) in the fifth row.

\subsection{Compact potential, induced radial velocity, and support}

Define the combined potential profile
\begin{equation}\label{mcmean:H}
 H_\theta(x,\eta)=\frac1x\int_0^x\xi G_{J_\theta}(\xi,\eta)\,d\xi.
\end{equation}
Choose \(0<x_-<x_+\) such that the convex hull of all the bump
supports lies in \([x_-,x_+]\Subset I_m\). For \(x<x_-\) the
integral in \eqref{mcmean:H} is zero. For \(x>x_+\) it is the
zero axial moment in \eqref{mcmean:profile_moments}.
It follows that \(H_\theta\), together with every derivative, is
supported in this fixed compact interval. Its smoothness at the two
collars follows from the fact that the integrand is identically zero
there and that its total integral is zero. In particular it extends
smoothly by zero to \(x\le0\).

The primitive of an individual \(\varphi_j\) is
\(\psi_j=x^{-1}\int_0^x\xi\varphi_j(\xi)\,d\xi\), which equals
\(\mu_1e^{jd_b}/x\) past its support. Thus an individual primitive
can have a tail. The exact identity
\(\sum_{j=0}^1 B_z[j,1]e^{jd_b}=0\) cancels these tails in
\(H_\theta=-(a\mu_{1-2\lambda})^{-1}\sum_j B_z[j,1]\psi_j\).
The identity holds for every \(\eta\), so its \(\eta\) derivatives
also vanish. This proves compactness for the combined profile and
all derivatives without assuming compactness for the summands.

The physical potential and radial increment are exactly
\begin{equation}\label{mcmean:potential_radial}
\begin{split}
 \Psi(r,z,t)&=\frac1r\int_0^r s\delta\gamma(s,z,t)\,ds
             =q^{A-1}\mathcal J_\theta H_\theta(x,\eta),\\
 \delta\beta&=-\partial_z\Psi\\
 &=-q^{A-1}(\partial_z\mathcal J_\theta)H_\theta
   -q^{A-1-D}\mathcal J_\theta\mathcal Z_{A-1}H_\theta.
\end{split}
\end{equation}
The substitution in the integral supplies one factor \(q\) and its
outer \(1/r\) supplies \(q^{-1/2}\), giving the exponent \(A-1\).
The operator in the second line is explicitly
\[
 \mathcal Z_\sigma f
 =\frac{2\sigma\eta f-\eta x f_x+(1-\eta^2)f_\eta}{B(\eta)}.
\]
The chain rule \eqref{mcq:dilation_dictionary} proves the last line,
including the \(\eta\) derivative of the original factor \(1/a\).
Direct differentiation of \(xH_\theta\) gives
\((\partial_x+x^{-1})H_\theta=G_{J_\theta}\), so
\((\partial_r+r^{-1})\Psi=\delta\gamma\). Commuting the smooth
physical \(r,z\) derivatives proves
\begin{equation}\label{mcmean:divergence}
 (\partial_r+r^{-1})\delta\beta+\partial_z\delta\gamma=0.
\end{equation}
The three velocity increments and \(\Psi\) are auxiliary independent
and angularly independent; their support lies in
\(\sqrt q[x_-,x_+]\), including any intervals between the bumps.

\subsection{Every mixed physical derivative, with original scale powers}

For clarity the original differential dictionary is
\[
 q_z=\frac{2\eta q^A}{B},\quad q_t=-\frac1B,\quad
 \eta_z=\frac{q^{-D}(1-\eta^2)}B,\quad
 \eta_t=\frac{D\eta q^{-1}}B,
 \qquad
 \mathcal T_\sigma f=
 \frac{-\sigma f+\tfrac12x f_x+D\eta f_\eta}{B}.
\]
These equations follow by implicit differentiation of
\eqref{mcmean:original_scale}; in particular
\(1-2hz^2q^{2h-1}=B\) and \(D+2h=A\).
Differentiating \(x=rq^{-1/2}\) supplies
\(x_z=-\eta xq^{-D}/B\) and \(x_t=xq^{-1}/(2B)\).
Consequently
\[
 \partial_r(q^\sigma F)=q^{\sigma-1/2}F_x,\quad
 \partial_z(q^\sigma F)=q^{\sigma-D}\mathcal Z_\sigma F,\quad
 \partial_t(q^\sigma F)=q^{\sigma-1}\mathcal T_\sigma F.
\]
Here \(F=F(x,\eta)\). These are physical derivatives of the moving
scale; the time loss is the original one in this scale.

For \(k,m,n\ge0\) define \(F_{\sigma;k,m,n}\) by the following
finite recursion. Start with \(f=\partial_x^kF\) and
\(s=\sigma-k/2\). Repeat
\(f\leftarrow\mathcal Z_s f,\ s\leftarrow s-D\) exactly \(m\)
times; then repeat \(f\leftarrow\mathcal T_s f,\ s\leftarrow s-1\)
exactly \(n\) times. The final \(f\) is \(F_{\sigma;k,m,n}\).
For any actual scalar target \(X=X(z,t)\), the exact product formula is
\begin{equation}\label{mcmean:all_derivatives}
\begin{split}
 \partial_r^k\partial_z^m\partial_t^n(q^\sigma XF)
  =\sum_{i=0}^m\sum_{j=0}^n&\binom mi\binom nj
     (\partial_z^i\partial_t^jX)\\
  &\cdot q^{\sigma-k/2-(m-i)D-(n-j)}
       F_{\sigma;k,m-i,n-j}(x,\eta).
\end{split}
\end{equation}
To prove it, first apply \(k\) radial derivatives, which do not act
on \(X,q,\eta\). The repeated product rule for the \(m\) axial and
\(n\) time derivatives assigns \(i,j\) derivatives to \(X\), with
the two binomial multiplicities displayed. Apply the three chain
rules above successively to the remaining factor. Each axial
derivative lowers its current exponent by \(D\), and each time
derivative lowers it by one, while its current exponent is the index
in the next operator. This is precisely the stated recursion.
Smooth mixed derivatives commute, so the chosen radial-then-axial-
then-time order computes the derivative written in the formula.

This proves every mixed derivative of \(\delta v,\delta\gamma,\Psi\)
by using, respectively, the two terms and the one term in
\eqref{mcmean:tangential_increment} and \eqref{mcmean:potential_radial}.
It proves every derivative of \(\delta\beta\) either by applying
\eqref{mcmean:all_derivatives} to its two displayed terms or by
\begin{equation}\label{mcmean:radial_all_derivatives}
 \partial_r^k\partial_z^m\partial_t^n\delta\beta
 =-\partial_r^k\partial_z^{m+1}\partial_t^n
          (q^{A-1}\mathcal J_\theta H_\theta).
\end{equation}
This equality retains, in particular, the derivative of the full
target and the additional original factor \(q^{-D}\).

The recursion is an explicit finite rational expression in
\(x,\eta,h\), the retained profile derivatives, and powers of
\(B^{-1}\). The other profile constants are exactly the inverse
Vandermonde entries and moments in \eqref{mcmean:profiles}.
In fact the reciprocal base coefficient has the exact derivatives
\begin{equation}\label{mcmean:reciprocal_a}
 \partial_\eta^s(a^{-1})=
 \frac{1}{2^{1/2+\lambda}c_{\rm patch}}
 \begin{cases}1+\eta^2&s=0,\\2\eta&s=1,\\2&s=2,\\0&s\ge3.
 \end{cases}
\end{equation}
Also \(\varphi_j^{(k)}(x)=e^{-(k+1)jd_b}
\varphi_0^{(k)}(e^{-jd_b}x)\).
These two formulas give every mixed derivative of the three profiles
without an unspecified inverse-matrix derivative constant.

For a similarly explicit derivative of \(H_\theta\), write
\(F_0(x,\eta)=\int_0^x\xi G_{J_\theta}(\xi,\eta)\,d\xi\).
For \(\ell\ge1\) the fundamental theorem of calculus and the
product rule give
\[
 \partial_x^\ell F_0=x\partial_x^{\ell-1}G_{J_\theta}
              +(\ell-1)\partial_x^{\ell-2}G_{J_\theta},
\]
where the second term is defined to be zero for \(\ell=1\).
Therefore, for \(x>0\),
\begin{equation}\label{mcmean:H_derivatives}
 \partial_x^k\partial_\eta^sH_\theta
 =\sum_{\ell=0}^k\binom k\ell(-1)^{k-\ell}(k-\ell)!
    x^{-1-k+\ell}\partial_x^\ell\partial_\eta^sF_0.
\end{equation}
The \(\ell=0\) factor is the explicitly specified integral; every
other factor is the displayed profile derivative. For example its
absolute integral is at most
\(\int_{x_-}^{x_+}\xi|\partial_\eta^sG_{J_\theta}(\xi,\eta)|\,d\xi\).
This supplies finite computable profile costs on the exact support.
The bound \(B\ge1-2h>0\), together with \(x\ge x_->0\), proves
that each fixed recursively specified coefficient is bounded on
\([x_-,x_+]\times[-1,1]\). The original moments, spacing, \(h\),
\(c_{\rm patch}\), and \(\lambda\) remain in those coefficients.

\subsection{Explicit costs in full target derivative arrays}

For \(I=(m,n)\), put \(I_z=(m+1,n)\). Use the original actual-shape
arrays \(\mathcal B_{e,I}\) from \eqref{mcq:B}, with their two
retained coupled amplitude components and old-wave cross terms.
Define the full target majorants
\begin{equation}\label{mcmean:target_arrays}
\begin{split}
 \mathsf P_I&=|\partial_z^m\partial_t^nP_{\rm old}|
          +\mathcal B_{-1,I}+\tfrac12\mathcal B_{0,I_z},\\
 \mathsf J_{\theta,I}
      &=|\partial_z^m\partial_t^n(J_\theta)_{\rm old}|
          +\tfrac12\mathcal B_{2,I},\\
 \mathsf J_{z,I}
      &=|\partial_z^m\partial_t^n(J_z)_{\rm old}|
          +\mathcal B_{1,I}+\tfrac14\mathcal B_{2,I_z}.
\end{split}
\end{equation}
The triangle inequality in \eqref{mcmean:full_targets} and the
proved increment bounds \eqref{mcq:defect_bounds} give
\(|\partial^I\mathcal P|\le\mathsf P_I\),
\(|\partial^I\mathcal J_\theta|\le\mathsf J_{\theta,I}\), and
\(|\partial^I\mathcal J_z|\le\mathsf J_{z,I}\).
Thus the old absolute derivatives appear explicitly in every cost.

For a nonnegative target array \(\mathsf X\), define the following
pointwise, fully specified finite sum:
\begin{equation}\label{mcmean:cost_operator}
\begin{split}
 \mathfrak C_{\sigma,F;k,m,n}[\mathsf X](r,z,t)
  =\sum_{i=0}^m\sum_{j=0}^n&\binom mi\binom nj\mathsf X_{(i,j)}(z,t)\\
 &\cdot q^{\sigma-k/2-(m-i)D-(n-j)}
       |F_{\sigma;k,m-i,n-j}(x,\eta)|.
\end{split}
\end{equation}
All powers are powers of the actual positive \(q\), even when their
exponents are negative. Define
\begin{equation}\label{mcmean:velocity_arrays}
\begin{split}
 \mathsf A_{\theta;k,m,n}
    &=\mathfrak C_{A,V_P;k,m,n}[\mathsf P]
        +\mathfrak C_{A-1,V_{J_z};k,m,n}[\mathsf J_z],\\
 \mathsf A_{z;k,m,n}
    &=\mathfrak C_{A-3/2,G_{J_\theta};k,m,n}[\mathsf J_\theta],\\
 \mathsf S_{k,m,n}
    &=\mathfrak C_{A-1,H_\theta;k,m,n}[\mathsf J_\theta],\\
 \mathsf A_{r;k,m,n}
    &=\mathfrak C_{A-1,H_\theta;k,m,n}[\mathsf J_\theta^{+z}]
      +\mathfrak C_{A-1-D,\mathcal Z_{A-1}H_\theta;k,m,n}
                                                 [\mathsf J_\theta],
 \qquad (\mathsf J_\theta^{+z})_{(i,j)}=\mathsf J_{\theta,(i+1,j)}.
\end{split}
\end{equation}
Equations \eqref{mcmean:all_derivatives}--\eqref{mcmean:potential_radial}
and the triangle inequality prove, for every \(k,m,n\),
\begin{equation}\label{mcmean:velocity_bounds}
\begin{gathered}
 |\partial_r^k\partial_z^m\partial_t^n\delta v|
       \le\mathsf A_{\theta;k,m,n},\qquad
 |\partial_r^k\partial_z^m\partial_t^n\delta\gamma|
       \le\mathsf A_{z;k,m,n},\\
 |\partial_r^k\partial_z^m\partial_t^n\Psi|
       \le\mathsf S_{k,m,n},\qquad
 |\partial_r^k\partial_z^m\partial_t^n\delta\beta|
       \le\mathsf A_{r;k,m,n}.
\end{gathered}
\end{equation}
The alternative upper bound \(\mathsf S_{k,m+1,n}\) for the last
quantity follows directly from \eqref{mcmean:radial_all_derivatives}.
For example the zero-order costs read
\[
\begin{split}
 |\delta v|&\le q^A\mathsf P_{0,0}|V_P|
                    +q^{A-1}\mathsf J_{z,0,0}|V_{J_z}|,\\
 |\delta\gamma|&\le q^{A-3/2}\mathsf J_{\theta,0,0}|G_{J_\theta}|,
 \qquad |\Psi|\le q^{A-1}\mathsf J_{\theta,0,0}|H_\theta|,\\
 |\delta\beta|&\le q^{A-1}\mathsf J_{\theta,1,0}|H_\theta|
       +q^{A-1-D}\mathsf J_{\theta,0,0}|\mathcal Z_{A-1}H_\theta|.
\end{split}
\]
These display separately the target derivative and moving-scale radial
cost. Replacing each profile value in \eqref{mcmean:cost_operator}
by its maximum over its fixed \(x,\eta\) support supplies a pointwise
bound depending only on \(z,t\) and the stated fixed profile data.

The same formula gives weighted radial costs without a loss hidden in
a constant. For any real \(e\), replace the final profile absolute
value in \eqref{mcmean:cost_operator} by
\[
 q^{(e+1)/4}
 \left(\int_{x_-}^{x_+}x^e
          |F_{\sigma;k,m-i,n-j}(x,\eta)|^2\,dx\right)^{1/2}.
\]
Call the resulting sum \(\mathfrak C^{[e]}\). The exact measure
identity \(r^e\,dr=q^{(e+1)/2}x^e\,dx\), followed by the triangle
inequality in this weighted \(L^2\) space, proves all four bounds
\eqref{mcmean:velocity_bounds} with the left sides replaced by their
\([e]\) norms and the right sides by the corresponding
\(\mathfrak C^{[e]}\) sums. The support is bounded away from zero,
so these integrals are finite for every fixed real \(e\).

\subsection{Actual pressure reconstruction and the surviving defects}

Set \(\boldsymbol a=(a_r,a_\theta,a_z)
=(\delta\beta,\delta v,\delta\gamma)\). Before pressure is
changed, let \(\delta g_r\) denote the exact change in the source's
angularly averaged negative radial residual. The angular average here
retains auxiliary dependence. The actual pressure-defect change is
\begin{equation}\label{mcmean:pressure_defect_step}
 \delta P_{\rm step}=\int_0^\infty\overline{\delta g_r}\,dr
                    =P_{\rm new}-\mathcal P.
\end{equation}
This is the defect change caused by the present mean operation;
it is different from \(\Delta P\) in \eqref{mcmean:full_targets}.
The source pressure bump is still
\(\rho=q^{-1/2}\widehat\rho(x)\),
\(\int\widehat\rho\,dx=1\). The actual pressure increment for this
operation is
\begin{equation}\label{mcmean:actual_pressure}
 \delta p=\mathcal I_c f,\qquad
 f=\delta g_r-\rho\,\delta P_{\rm step}.
\end{equation}
The original physical shifted primitive is
\[
\begin{split}
 \mathcal If(r,Y)&=\int_0^r
  f(s,Y+(s^{d_r}-r^{d_r})v_r,z,t)\,ds,\\
 \mathcal Jf(r,Y)&=\int_0^\infty
  f(s,Y+(s^{d_r}-r^{d_r})v_r,z,t)\,ds,\\
 \mathcal I_cf&=\mathcal If-\chi_m\mathcal Jf,\\
 d_r&=2((1+h)\rho_g-h\kappa_s),\quad
 \rho_g=\frac{\log(4-\sqrt2)}{\log(4+\sqrt2)},\quad
 \kappa_s=10^{-5},\quad v_r=(1,-(\sqrt2-1)).
\end{split}
\]
Here \(\chi_m\) is the actual inner-zero, outer-one source cutoff;
the two suppressed arguments of each integrand are exactly \(z,t\).
The pressure entering the physical force is the physical evaluation of
this source function, with its original phase evaluation and all its
derivatives. Thus it includes any nonzero auxiliary modes of the
source reconstruction.

To verify the reconstruction identity, the source physical radial
operator \(\mathcal D_r\) cancels the derivative of the term
\(-r^{d_r}v_r\) in the shift. Endpoint differentiation gives
\(\mathcal D_r\mathcal If=f\) and
\(\mathcal D_r\mathcal Jf=0\), and hence
\begin{equation}\label{mcmean:pressure_remainder}
 \mathcal D_r\delta p-\delta g_r
 =-\rho\,\delta P_{\rm step}
                  -(\partial_r\chi_m)\mathcal Jf.
\end{equation}
Haar translation invariance and \eqref{mcmean:pressure_defect_step}
give \(\overline{\mathcal Jf}=0\), but this does not set
\(\mathcal Jf\) itself to zero. In the inner collar \(\mathcal If=0\)
and \(\chi_m=0\); in the outer collar \(\mathcal If=\mathcal Jf\)
and \(\chi_m=1\). This proves compact support of the actual pressure
and all derivatives. In particular its gradient is retained in the
force calculation below.

For completeness the exact remaining source terms can be written
entirely in the original current mean fields. With
\(\Delta_0=\mathcal D_r^2+r^{-1}\mathcal D_r+\partial_z^2\),
direct expansion of the radial conservative source gives
\begin{equation}\label{mcmean:Rg}
\begin{split}
 R_g:={}&\delta g_r-\frac{2V}{r}\delta v\\
 ={}&-\partial_t\delta\beta
 -(\mathcal D_r+r^{-1})
       \bigl(2b\delta\beta+2\beta\delta\beta+(\delta\beta)^2\bigr)\\
 &-\partial_z\bigl(b\delta\gamma+G\delta\beta+
       \beta\delta\gamma+\gamma\delta\beta+
       \delta\beta\delta\gamma\bigr)\\
 &+r^{-1}\bigl(2v\delta v+(\delta v)^2\bigr)
       +\nu_{\rm NS}(\Delta_0-r^{-2})\delta\beta.
\end{split}
\end{equation}
Indeed the radial conservative quadratic expression before the change
is \(-(\mathcal D_r+r^{-1})(b+\beta)^2
-\partial_z((b+\beta)(G+\gamma))+(V+v)^2/r\), together with
the wave covariance, time term, and viscosity. Subtract it from its
value after the mean change. The unchanged wave covariance cancels;
each remaining quadratic product expands to the terms displayed.
The term \(2V\delta v/r\) has been written separately in the
definition of \(R_g\). The ordinary time derivative on the increment
is exact because it is auxiliary independent.

The five physical rows then give the exact new defects
\begin{equation}\label{mcmean:nonlinear_defects}
\begin{split}
 P_{\rm new}&=\int\overline{R_g}\,dr,\\
 (J_\theta)_{\rm new}
   &=\int r^2\overline{\gamma\delta v+v\delta\gamma+
                                    \delta\gamma\delta v}\,dr,\\
 (J_z)_{\rm new}
   &=\int r\overline{2\gamma\delta\gamma+(\delta\gamma)^2}\,dr
                        -\tfrac12\int r^2\overline{R_g}\,dr.
\end{split}
\end{equation}
For the first line add \(\mathcal P\) to
\(\int(2V\delta v/r+\overline{R_g})\,dr\) and use the third
row. For the second, expanding \((G+\gamma+\delta\gamma)
(V+v+\delta v)-(G+\gamma)(V+v)\) gives base-linear terms
\(G\delta v+V\delta\gamma\), cancelled by the fourth row, and
exactly the three remaining terms displayed. For the third, expand
\((G+\gamma+\delta\gamma)^2-(G+\gamma)^2\) and the radial-source
moment \(-\tfrac12\int r^2\delta g_r\,dr\); the base-linear
terms are the fifth row. This proves all three equalities.
Neither their right sides nor the cutoff term in
\eqref{mcmean:pressure_remainder} has been removed.

\subsection{Exact full cylindrical physical force increment}

For the actual physical fields, put
\(\mathcal R(U,p)=\partial_tU+(U\cdot\nabla)U
-\nu_{\rm NS}\Delta U+\nabla p\). Ordinary derivatives in this
formula are derivatives after the physical source evaluation; they
therefore include all phase, frame, cutoff and scale derivatives of
the actual full \(U\). Expanding the two quadratic products proves
the vector identity
\begin{equation}\label{mcmean:vector_force}
 \delta\mathcal R=\partial_t\boldsymbol a
 +(U\cdot\nabla)\boldsymbol a+(\boldsymbol a\cdot\nabla)U
 +(\boldsymbol a\cdot\nabla)\boldsymbol a
 -\nu_{\rm NS}\Delta\boldsymbol a+\nabla\delta p.
\end{equation}
The pressure here is precisely \eqref{mcmean:actual_pressure}.

To derive every cylindrical component, the basis derivatives are
\(\partial_\theta e_r=e_\theta\),
\(\partial_\theta e_\theta=-e_r\), and
\(\partial_\theta e_z=0\). Thus for two arbitrary vector fields
\(X,Y\) the component vector of \((X\cdot\nabla)Y\) is
\[
 X_r\partial_rY+\frac{X_\theta}{r}\partial_\theta Y
 +X_z\partial_zY+\frac{X_\theta}{r}(-Y_\theta,Y_r,0).
\]
Apply this equality to \((U,\boldsymbol a)\),
\((\boldsymbol a,U)\), and \((\boldsymbol a,\boldsymbol a)\).
The scalar components of \(\boldsymbol a\) have zero angular
derivative, whereas those of the actual \(U\) generally do not.
Writing \(L_0=\partial_r^2+r^{-1}\partial_r+\partial_z^2\), the
complete result is
\begin{equation}\label{mcmean:force_r}
\begin{split}
 \delta\mathcal R_r={}&\partial_ta_r+U_r\partial_ra_r+U_z\partial_za_r
       +a_r\partial_rU_r+a_z\partial_zU_r
       +\frac{a_\theta}{r}\partial_\theta U_r\\
 &+a_r\partial_ra_r+a_z\partial_za_r
       -\frac{2U_\theta a_\theta+a_\theta^2}{r}
       -\nu_{\rm NS}(L_0-r^{-2})a_r+\partial_r\delta p,
\end{split}
\end{equation}
\begin{equation}\label{mcmean:force_theta}
\begin{split}
 \delta\mathcal R_\theta={}&\partial_ta_\theta
       +U_r\partial_ra_\theta+U_z\partial_za_\theta
       +a_r\partial_rU_\theta+a_z\partial_zU_\theta
       +\frac{a_\theta}{r}\partial_\theta U_\theta\\
 &+a_r\partial_ra_\theta+a_z\partial_za_\theta
       +\frac{U_ra_\theta+a_rU_\theta+a_ra_\theta}{r}
       -\nu_{\rm NS}(L_0-r^{-2})a_\theta
       +\frac1r\partial_\theta\delta p,
\end{split}
\end{equation}
\begin{equation}\label{mcmean:force_z}
\begin{split}
 \delta\mathcal R_z={}&\partial_ta_z+U_r\partial_ra_z+U_z\partial_za_z
       +a_r\partial_rU_z+a_z\partial_zU_z
       +\frac{a_\theta}{r}\partial_\theta U_z\\
 &+a_r\partial_ra_z+a_z\partial_za_z
       -\nu_{\rm NS}L_0a_z+\partial_z\delta p.
\end{split}
\end{equation}
The self terms in the radial equation include
\(-a_\theta^2/r\); those in the angular equation include
\(a_ra_\theta/r\). The old-wave cross terms occur in the actual
\(U\) and in its angular derivatives in all three equations.
The mean pressure is angularly invariant, so its angular derivative
is zero for this source operation; its radial and axial derivatives
are the actual reconstructed ones and are generally nonzero.

For the viscosity formula, applying the scalar cylindrical Laplacian
to the moving basis yields, for a general vector field \(X\),
\[
 \Delta X=
 \left(\Delta_sX_r-r^{-2}X_r-2r^{-2}\partial_\theta X_\theta,
       \Delta_sX_\theta-r^{-2}X_\theta+2r^{-2}\partial_\theta X_r,
       \Delta_sX_z\right),
 \quad\Delta_s=L_0+r^{-2}\partial_\theta^2.
\]
This follows by differentiating the two basis identities twice and
applying the product rule. Substituting the angularly independent
components of \(\boldsymbol a\) gives exactly
\(((L_0-r^{-2})a_r,(L_0-r^{-2})a_\theta,L_0a_z)\), proving all
three displayed viscous terms.
The radial expression before the pressure term, when angularly
averaged in the extended source domain and negated, is precisely
\(\delta g_r\): this follows from the definition of the negative
radial source residual and linearity of the average. It is also an
explicit formula for the input to \eqref{mcmean:actual_pressure} in
terms of the same full velocity and mean increments.

\subsection{Finite all-order force costs with every inverse-radius term}

Here are explicit derivative bounds for all three components, including
angular derivatives of the actual current velocity. Let
\(\alpha=(k,\ell,m,n)\in\mathbb N^4\) index the scalar component
derivative \(\partial^\alpha=\partial_r^k\partial_\theta^\ell
\partial_z^m\partial_t^n\), and let
\(e_r,e_\theta,e_z,e_t\) be the four coordinate multi-indices.
Extend the arrays in \eqref{mcmean:velocity_arrays} by
\[
 \mathsf A_{i,\alpha}=
 \begin{cases}\mathsf A_{i;k,m,n}&\ell=0,\\0&\ell>0,
 \end{cases}
 \qquad
 \mathsf U_{i,\alpha}=|\partial^\alpha U_i|.
\]
The pressure array \(\mathsf\Pi\) will be the following explicit
majorant for the derivatives of the actual specified function
\eqref{mcmean:actual_pressure}, after physical evaluation.

The original physical radial phase satisfies
\(Y_{\rm phys}(s,z,t)-Y_{\rm phys}(r,z,t)
=(s^{d_r}-r^{d_r})v_r\). Substituting this exact equality in the
two shifted integrals proves the evaluated formula
\begin{equation}\label{mcmean:evaluated_pressure}
\begin{split}
 \delta p(r,z,t)
 &=\int_0^r f_{\rm phys}(s,z,t)\,ds
        -\chi_m(r,z,t)\int_0^\infty f_{\rm phys}(s,z,t)\,ds,\\
 f_{\rm phys}&=(\delta g_r)_{\rm phys}
                       -\rho\,\delta P_{\rm step}.
\end{split}
\end{equation}
Every original axial and temporal phase dependence is retained in
\(f_{\rm phys}\). The second integral is independent of \(r\)
after this evaluation. This equality does not assert that it is zero.

Fix any compact interior physical domain for the finite operation,
and choose one physical interval \([R_-,R_+]\Subset(0,\infty)\)
containing the supports of \(\delta g_r\) and \(\rho\) for that
domain and every auxiliary phase. Put \(L_R=R_+-R_-\), with these
actual endpoints kept fixed during differentiation. The smooth zero
extensions imply that all their derivatives have support in the
same interval. Define \(\mathsf G_{k,m,n}(z,t)\) as the supremum
over \(r\in[R_-,R_+]\) and the original auxiliary phases of the
angular average of the following nonnegative expression: the radial
cost \(\mathsf F_{r;(k,0,m,n)}\) in \eqref{mcmean:force_costs}
with its last pressure term deleted. This definition involves only
\(\mathsf A\), the actual current velocity derivatives,
\(\nu_{\rm NS}\), and the inverse-radius arrays; it has no
dependence on \(\mathsf\Pi\). When auxiliary phases are retained,
the velocity derivative arrays mean the corresponding full lifted
physical derivatives before evaluation. The chain rule for source
evaluation identifies these with derivatives of the evaluated fields
at each phase. The force identities and the derivative proof below
therefore imply
\[
 |\partial_r^k\partial_z^m\partial_t^n(\delta g_r)_{\rm phys}|
       \le\mathsf G_{k,m,n},\qquad
 |\partial_z^m\partial_t^n\delta P_{\rm step}|
       \le\mathsf D_{m,n}:=L_R\mathsf G_{0,m,n}.
\]
For the first inequality use the negative angularly averaged radial
force before pressure, as proved above, and the triangle inequality
under that average. For the second, differentiate
\eqref{mcmean:pressure_defect_step} under its compact radial
integral, commute axial and temporal differentiation with Haar
average, and integrate the same bound over the interval of width
\(L_R\). The Haar integral of each original constant torus
direction derivative is zero by periodicity, which justifies that
commutation for the full physical time derivative as well.

The derivative of the pressure bump is already an explicit original-
scale coefficient in \eqref{mcq:bump_derivatives}. Define its exact
support supremum
\[
 \mathsf H_{k,m,n}(z,t)
 =q^{-(1+k)/2-mD-n}
    \sup_{r\in[R_-,R_+]}
       |f_{0;k,m,n}(r/\sqrt q,\eta)|,
\]
where \(f_{0;k,m,n}\) is the recursion starting with the actual
\(\widehat\rho\), rather than another normalized bump.
It follows that \(|\partial_r^k\partial_z^m\partial_t^n\rho|
\le\mathsf H_{k,m,n}\). Set
\begin{equation}\label{mcmean:pressure_source_cost}
 \mathsf E_{k,m,n}
 =\mathsf G_{k,m,n}
 +\sum_{i=0}^m\sum_{j=0}^n\binom mi\binom nj
        \mathsf H_{k,m-i,n-j}\mathsf D_{i,j}.
\end{equation}
Since \(\delta P_{\rm step}\) has no radial dependence, all
\(k\) radial derivatives in its product with \(\rho\) act on
\(\rho\). The repeated axial and temporal product rule proves
\(|\partial_r^k\partial_z^m\partial_t^n f_{\rm phys}|
\le\mathsf E_{k,m,n}\) on the full interval.

Write \(\mathsf X_{k,i,j}(r,z,t)
=|\partial_r^k\partial_z^i\partial_t^j\chi_m(r,z,t)|\), retaining
the actual cutoff and its moving-scale derivatives. Define
\begin{equation}\label{mcmean:explicit_pressure_cost}
\begin{split}
 \mathsf\Pi_{(k,0,m,n)}={}&
 \begin{cases}L_R\mathsf E_{0,m,n}&k=0,\\
               \mathsf E_{k-1,m,n}&k\ge1
 \end{cases}\\
 &+L_R\sum_{i=0}^m\sum_{j=0}^n\binom mi\binom nj
           \mathsf X_{k,i,j}\mathsf E_{0,m-i,n-j},
 \qquad \mathsf\Pi_{(k,\ell,m,n)}=0\quad(\ell>0).
\end{split}
\end{equation}
This proves the explicit bound
\(|\partial^\alpha\delta p|\le\mathsf\Pi_\alpha\).
In fact differentiating \eqref{mcmean:evaluated_pressure} yields
the exact first-integral derivative
\[
 \partial_r^k\partial_z^m\partial_t^n
       \int_0^r f_{\rm phys}(s,z,t)\,ds
 =\begin{cases}
   \displaystyle\int_0^r\partial_z^m\partial_t^n
                         f_{\rm phys}(s,z,t)\,ds&k=0,\\
   \partial_r^{k-1}\partial_z^m\partial_t^n f_{\rm phys}(r,z,t)&k\ge1.
  \end{cases}
\]
The absolute values are bounded by the first term in
\eqref{mcmean:explicit_pressure_cost}. In the cutoff product all
radial derivatives act on \(\chi_m\), whereas the axial and temporal
derivatives distribute with the two displayed binomial factors.
Every derivative of its radial integral is bounded by
\(L_R\mathsf E_{0,m-i,n-j}\). This proves its second term, with
every cutoff derivative and the actual support width retained.
Angular invariance proves the stated zero entries. This construction
first computes the pressure-free \(\mathsf G\), then
\(\mathsf D,\mathsf H,\mathsf E,\mathsf\Pi\), and finally the
full force array below, so the pressure budget is a finite calculation
without circularly assuming a pressure bound.

For nonnegative arrays define their exact Leibniz convolution and
derivative shift by
\begin{equation}\label{mcmean:convolution}
 (X\star Y)_\alpha=\sum_{\beta\le\alpha}
               \binom\alpha\beta X_\beta Y_{\alpha-\beta},
 \qquad (S_\delta X)_\alpha=X_{\alpha+\delta}.
\end{equation}
Products with three factors use either parenthesization; multiplying
the binomial coefficients gives
\(\alpha!/(\beta!\gamma!(\alpha-\beta-\gamma)!)\), proving the
two parenthesizations agree. For the exact inverse-radius arrays put
\begin{equation}\label{mcmean:radius_array}
 (\mathsf R_j)_\alpha=
 \begin{cases}
 (j)_k r^{-j-k}&\ell=m=n=0,\\0&\text{otherwise},
 \end{cases}
 \quad (j)_0=1,\quad(j)_k=j(j+1)\cdots(j+k-1),\quad j=1,2.
\end{equation}
Induction gives \(\partial_r^k r^{-j}=(-1)^k(j)_kr^{-j-k}\),
which proves that \(\mathsf R_j\) is its exact absolute derivative
array. No radial derivative of an inverse-radius coefficient is
omitted in the following costs.

For \(i=r,\theta,z\) define the common transport cost
\begin{equation}\label{mcmean:transport_cost}
\begin{split}
 \mathsf T_i={}&S_{e_t}\mathsf A_i
 +\mathsf U_r\star S_{e_r}\mathsf A_i
 +\mathsf U_z\star S_{e_z}\mathsf A_i
 +\mathsf A_r\star S_{e_r}\mathsf U_i
 +\mathsf A_z\star S_{e_z}\mathsf U_i\\
 &+\mathsf R_1\star\mathsf A_\theta\star S_{e_\theta}\mathsf U_i
 +\mathsf A_r\star S_{e_r}\mathsf A_i
 +\mathsf A_z\star S_{e_z}\mathsf A_i,
\end{split}
\end{equation}
and the exact finite upper-cost arrays
\begin{equation}\label{mcmean:force_costs}
\begin{split}
 \mathsf F_r={}&\mathsf T_r
    +2\mathsf R_1\star\mathsf U_\theta\star\mathsf A_\theta
    +\mathsf R_1\star\mathsf A_\theta\star\mathsf A_\theta\\
 &+\nu_{\rm NS}\bigl(S_{2e_r}\mathsf A_r
       +\mathsf R_1\star S_{e_r}\mathsf A_r
       +S_{2e_z}\mathsf A_r+\mathsf R_2\star\mathsf A_r\bigr)
       +S_{e_r}\mathsf\Pi,\\
 \mathsf F_\theta={}&\mathsf T_\theta
    +\mathsf R_1\star\bigl(\mathsf U_r\star\mathsf A_\theta
                    +\mathsf A_r\star\mathsf U_\theta
                    +\mathsf A_r\star\mathsf A_\theta\bigr)\\
 &+\nu_{\rm NS}\bigl(S_{2e_r}\mathsf A_\theta
       +\mathsf R_1\star S_{e_r}\mathsf A_\theta
       +S_{2e_z}\mathsf A_\theta+\mathsf R_2\star\mathsf A_\theta\bigr)
       +\mathsf R_1\star S_{e_\theta}\mathsf\Pi,\\
 \mathsf F_z={}&\mathsf T_z
    +\nu_{\rm NS}\bigl(S_{2e_r}\mathsf A_z
       +\mathsf R_1\star S_{e_r}\mathsf A_z
       +S_{2e_z}\mathsf A_z\bigr)+S_{e_z}\mathsf\Pi.
\end{split}
\end{equation}
Every addition and scalar multiplication here is entrywise. For every
finite \(\alpha\), these arrays prove the full physical component
inequalities
\begin{equation}\label{mcmean:all_force_bounds}
 |\partial^\alpha\delta\mathcal R_i|\le\mathsf F_{i,\alpha},
 \qquad i=r,\theta,z.
\end{equation}
Indeed the repeated Leibniz rule gives
\(|\partial^\alpha(fg)|\le(X\star Y)_\alpha\) whenever the
respective derivative absolute values are bounded by \(X,Y\).
The three-factor version is its iterated rule, proved by the explicit
multinomial coefficient above. Apply these rules term by term to
\eqref{mcmean:force_r}--\eqref{mcmean:force_z}, using
\eqref{mcmean:velocity_bounds} for the increment, the actual
derivatives for \(U,\delta p\), and \eqref{mcmean:radius_array} for
each inverse radius. The constant viscosity passes through every
derivative. The three resulting sums are precisely
\eqref{mcmean:transport_cost}--\eqref{mcmean:force_costs}, proving
the claim at every order.

These are component derivatives in the original orthonormal basis.
For an explicit bound on the actual Cartesian vector after angular
differentiation, the moving basis gives
\[
 \left|\partial_r^k\partial_\theta^\ell\partial_z^m\partial_t^n
                 \delta\mathcal R\right|
 \le\sum_{j=0}^\ell\binom\ell j
       \bigl(\mathsf F_{r;(k,j,m,n)}
             +\mathsf F_{\theta;(k,j,m,n)}\bigr)
       +\mathsf F_{z;(k,\ell,m,n)}.
\]
To prove it, expand each angular derivative by the binomial rule
between a scalar component and its basis vector. Every derivative of
\(e_r,e_\theta\) has norm one, whereas the derivatives of \(e_z\)
of positive order vanish. The triangle inequality gives exactly the
displayed bound. Thus the component formulas also provide a bound
for the physical vector, with its angular basis differentiation.

There is also an exact recursion for every Cartesian derivative.
Let \(O_\theta\) be the orthogonal matrix with columns
\((e_r,e_\theta,e_z)\), and let
\(J(f_r,f_\theta,f_z)=(-f_\theta,f_r,0)\). Use
\(x_{\rm C}=r\cos\theta,y_{\rm C}=r\sin\theta\) for the physical
Cartesian coordinates, distinct from the original scaled radius
\(x=r/\sqrt q\). Define the operators on component vectors
\begin{equation}\label{mcmean:cartesian_recursion}
\begin{split}
 \mathcal Xf&=\cos\theta\,\partial_rf
        -\frac{\sin\theta}{r}(\partial_\theta f+Jf),\\
 \mathcal Yf&=\sin\theta\,\partial_rf
        +\frac{\cos\theta}{r}(\partial_\theta f+Jf),
 \qquad \mathcal Zf=\partial_zf.
\end{split}
\end{equation}
The scalar coordinate rules
\(\partial_{x_{\rm C}}=\cos\theta\partial_r
-r^{-1}\sin\theta\partial_\theta\) and
\(\partial_{y_{\rm C}}=\sin\theta\partial_r
+r^{-1}\cos\theta\partial_\theta\), together with
\(\partial_\theta O_\theta=O_\theta J\), give exactly
\(\partial_{x_{\rm C}}(O_\theta f)=O_\theta\mathcal Xf\),
\(\partial_{y_{\rm C}}(O_\theta f)=O_\theta\mathcal Yf\), and
\(\partial_z(O_\theta f)=O_\theta\mathcal Zf\).
Applying these identities repeatedly proves
\[
 \partial_{x_{\rm C}}^i\partial_{y_{\rm C}}^j
 \partial_z^k\partial_t^n(O_\theta f)
   =O_\theta\mathcal X^i\mathcal Y^j\mathcal Z^k\partial_t^nf.
\]
The order written is a fixed composition order and is exact because
it was obtained from commuting smooth Cartesian derivatives.
For a fully explicit majorant recursion, if \(\mathsf C_r,
\mathsf C_\theta,\mathsf C_z\) majorize the component derivative
arrays of \(f\), set \((\mathsf J C)_r=\mathsf C_\theta\),
\((\mathsf J C)_\theta=\mathsf C_r\),
\((\mathsf J C)_z=0\). Let \(\mathsf c_\alpha
=|\partial^\alpha\cos\theta|\),
\(\mathsf s_\alpha=|\partial^\alpha\sin\theta|\), whose entries
vanish unless \(k=m=n=0\) and whose remaining entries are the
actual derivatives of these elementary functions. Then one
\(\mathcal X\) application is bounded componentwise by
\[
 \mathsf c\star S_{e_r}\mathsf C_i
 +\mathsf s\star\mathsf R_1\star
             (S_{e_\theta}\mathsf C_i+(\mathsf J C)_i),
\]
and one \(\mathcal Y\) application by the same expression with
\(\mathsf c,\mathsf s\) interchanged. A \(\mathcal Z\) or time
application is the array shift \(S_{e_z}\) or \(S_{e_t}\).
These rules follow from the exact product rule
\eqref{mcmean:convolution}. Starting with \(\mathsf F\), iterating
them the indicated finite number of times, and taking the sum of
the three resulting zero-order component bounds gives a finite bound
for each stated Cartesian vector derivative. Thus all Cartesian
inverse-radius and angular-basis costs are also explicitly supplied.

All costs above are finite for the actual finite state on each compact
interior physical domain. Their dependence on the original scale,
target derivatives, inverse-radius factors, viscosity, full current
velocity, and reconstructed pressure is explicit. The exact nonlinear
defects \eqref{mcmean:nonlinear_defects} remain after this operation.
No uniform estimate over an increasing family of source stages,
no singular-time limit, and no infinite correction gain is asserted
by these finite derivative identities and costs.
